% Arithmétique — Mathématiques 3ème — Cours signé OnBuch+
% Programme officiel MINESEC. Compiler avec : tectonic N01-arithmetique.tex
\newcommand{\DOCMATIERE}{Mathématiques}
\newcommand{\DOCNIVEAU}{3ème}
\newcommand{\DOCMODULE}{Mathématiques – classe de 3ème}
\newcommand{\DOCLECON}{N01}
\newcommand{\DOCTITRE}{Arithmétique}
% OnBuch+ — préambule « cours signés » ENRICHI (compilé avec tectonic / XeLaTeX)
% Système visuel « Pop » d'OnBuch+ : fond crème (jamais blanc), bordures encre
% épaisses, ombres « dures » décalées, titres Archivo Black en capitales.
% Version MATHÉMATIQUES : graphes (pgfplots/tikz), boîtes pédagogiques
% (définition, propriété, méthode, attention, le savais-tu, exercice résolu,
% exercice, TP, bilan), page de garde et signature OnBuch+ ancrée en bas.
%
% Macros attendues dans main.tex AVANT \input{preamble} :
%   \DOCMATIERE  \DOCNIVEAU  \DOCMODULE  \DOCLECON (numéro)  \DOCTITRE
\documentclass[11pt]{article}

\usepackage{fontspec}
\usepackage[french]{babel}
\usepackage[a4paper,margin=2cm,top=3.4cm,bottom=2.8cm,headheight=1.4cm,footskip=1.3cm]{geometry}
\usepackage{amsmath,amssymb,amsfonts}
\usepackage{mathtools} % \xmapsto, \coloneqq... souvent utilisés par les modèles
\usepackage{cancel} % simplifications barrées
% \abs{z} (module, valeur absolue) et \norm{u} (norme), utilisés par les modèles
\providecommand{\abs}{}\let\abs\relax\DeclarePairedDelimiter{\abs}{\lvert}{\rvert}
\providecommand{\norm}{}\let\norm\relax\DeclarePairedDelimiter{\norm}{\lVert}{\rVert}
\usepackage[table]{xcolor} % [table] : fournit aussi \rowcolors (alternance de lignes)
\usepackage{pagecolor}
\usepackage{tikz}
\usetikzlibrary{arrows.meta,positioning,calc,shapes.geometric,shapes.misc,shapes.arrows,decorations.pathmorphing,decorations.pathreplacing,decorations.markings,patterns,shadows,mindmap,trees,fit,backgrounds,matrix,intersections,angles,quotes}
\usepackage{pgfplots}
\usepackage{pgf-pie} % diagrammes circulaires \pie (statistiques)
\pgfplotsset{compat=1.18}
\usepgfplotslibrary{fillbetween,groupplots} % aires entre courbes (calcul intégral)
\usepackage{enumitem}
\usepackage{fancyhdr}
% [most] charge toutes les bibliothèques tcolorbox (listings, minted, raster,
% documentation...) et gonfle inutilement la mémoire TeX sur un document long
% (~300 boîtes) : on ne charge que ce qui est réellement utilisé.
\usepackage[skins,breakable]{tcolorbox}
\usepackage{graphicx}
\usepackage{colortbl}
\usepackage{booktabs}
\usepackage{tabularx}
\usepackage{diagbox} % case d'en-tête diagonale des tableaux à double entrée
% Colonnes extensibles centrée / gauche / droite (C, L, R), souvent utilisées par les modèles
\newcolumntype{C}{>{\centering\arraybackslash}X}
\newcolumntype{L}{>{\raggedright\arraybackslash}X}
\newcolumntype{R}{>{\raggedleft\arraybackslash}X}
\usepackage{array}
\usepackage{multicol}
\usepackage{siunitx}
% parse-numbers=false : accepte aussi \SI{2,5 \times 10^{-3}}{...}
\sisetup{locale=FR,output-decimal-marker={,},group-minimum-digits=4,parse-numbers=false}
\DeclareSIUnit\ohm{\ensuremath{\symup{\Omega}}} % U+2126 absent de la police
\usepackage{qrcode}
% \degree : commande fréquemment utilisée par les modèles pour un angle ou une
% température en dehors de \SI{}{} (non fournie par les paquets chargés).
\providecommand{\degree}{^{\circ}}
% \qed (fin de démonstration), utilisé par les modèles sans amsthm
\providecommand{\qed}{\hfill$\square$}
% \barre : double barre des valeurs interdites dans un tableau de variations (tkz-tab)
\providecommand{\barre}{\ensuremath{\|}}
% environnement proof (amsthm non chargé) utilisé par les modèles
\ifdefined\proof\else\newenvironment{proof}[1][Démonstration]{\par\noindent\textit{#1.}\ }{\qed\par}\fi
\usepackage{lastpage}
\usepackage{hyperref}
\hypersetup{hidelinks}

% ============================================================
% Palette « Pop » OnBuch+ — mêmes hex que la classe P (plus_ui.dart)
% ============================================================
\definecolor{poporange}{HTML}{F59321}
\definecolor{popdark}{HTML}{E07A0C}
\definecolor{popcream}{HTML}{FFF6E8}
\definecolor{popink}{HTML}{1C1714}
\definecolor{poppurple}{HTML}{7A5AE0}
\definecolor{popgreen}{HTML}{2E9E6A}
\definecolor{poppink}{HTML}{E0457B}
\definecolor{popblue}{HTML}{2F7DE1}
\definecolor{popgold}{HTML}{C98A12}
\definecolor{popmuted}{HTML}{8A7F76}
\definecolor{popline}{HTML}{EADFD2}
\definecolor{popgray}{HTML}{8A7F76}
\colorlet{popgrey}{popgray}
% Alias tolérants (le modèle nomme parfois les couleurs différemment)
\colorlet{popwhite}{white}
\colorlet{popblack}{popink}
\colorlet{popred}{poppink}
\colorlet{popyellow}{popgold}
% Teintes claires pour fonds de tableaux / zones
\colorlet{popblueL}{popblue!10!white}
\colorlet{poporangeL}{poporange!14!white}
\colorlet{popgreenL}{popgreen!12!white}
\colorlet{poppurpleL}{poppurple!12!white}
\colorlet{poppinkL}{poppink!12!white}
\colorlet{popgoldL}{popgold!14!white}

\pagecolor{popcream}

% ============================================================
% Typographie
% ============================================================
\defaultfontfeatures{Ligatures=TeX}
\setmainfont{PlusJakartaSans-Medium.ttf}[Path=../../../fonts/,BoldFont=PlusJakartaSans-Bold.ttf]
% Maths en sans-serif (Fira Math) avec chiffres et lettres droites de Plus
% Jakarta, pour que les formules se fondent dans le texte.
\usepackage[math-style=ISO,bold-style=ISO]{unicode-math}
\setmathfont{FiraMath-Regular.otf}[Path=../../../fonts/]
\setmathfont{PlusJakartaSans-Medium.ttf}[Path=../../../fonts/,range={up/num,up/{Latin,latin},\mathup}]
\usepackage{icomma} % virgule décimale sans espace en mode math
% Caractères mathématiques Unicode absents des polices de texte (Plus Jakarta,
% Archivo Black) : rendus par leur équivalent mathématique, en texte comme en
% formule — sinon ils s'affichent en carré vide (ex. « Arithmétique dans ℤ »).
\usepackage{newunicodechar}
\newunicodechar{ℝ}{\ensuremath{\mathbb{R}}}
\newunicodechar{ℤ}{\ensuremath{\mathbb{Z}}}
\newunicodechar{ℂ}{\ensuremath{\mathbb{C}}}
\newunicodechar{ℕ}{\ensuremath{\mathbb{N}}}
\newunicodechar{ℚ}{\ensuremath{\mathbb{Q}}}
\newunicodechar{𝔻}{\ensuremath{\mathbb{D}}}
\newunicodechar{α}{\ensuremath{\alpha}}
\newunicodechar{β}{\ensuremath{\beta}}
\newunicodechar{γ}{\ensuremath{\gamma}}
\newunicodechar{δ}{\ensuremath{\delta}}
\newunicodechar{ε}{\ensuremath{\varepsilon}}
\newunicodechar{θ}{\ensuremath{\theta}}
\newunicodechar{λ}{\ensuremath{\lambda}}
\newunicodechar{μ}{\ensuremath{\mu}}
\newunicodechar{σ}{\ensuremath{\sigma}}
\newunicodechar{τ}{\ensuremath{\tau}}
\newunicodechar{φ}{\ensuremath{\varphi}}
\newunicodechar{ψ}{\ensuremath{\psi}}
\newunicodechar{ω}{\ensuremath{\omega}}
\newunicodechar{Σ}{\ensuremath{\Sigma}}
% majuscules produites par \MakeUppercase dans les titres (α-aminés -> Α)
\newunicodechar{Α}{\ensuremath{\alpha}}
\newunicodechar{Β}{\ensuremath{\beta}}
\newunicodechar{Γ}{\ensuremath{\Gamma}}
\newunicodechar{Δ}{\ensuremath{\Delta}}
\newunicodechar{Ε}{\ensuremath{\varepsilon}}
\newunicodechar{Θ}{\ensuremath{\Theta}}
\newunicodechar{Λ}{\ensuremath{\Lambda}}
\newunicodechar{Μ}{\ensuremath{\mu}}
\newunicodechar{Τ}{\ensuremath{\tau}}
\newunicodechar{Φ}{\ensuremath{\Phi}}
\newunicodechar{Ψ}{\ensuremath{\Psi}}
\newunicodechar{Ω}{\ensuremath{\Omega}}
\newunicodechar{↦}{\ensuremath{\mapsto}}
\newunicodechar{∈}{\ensuremath{\in}}
\newunicodechar{∉}{\ensuremath{\notin}}
\newunicodechar{∧}{\ensuremath{\wedge}}
\newunicodechar{⇔}{\ensuremath{\Leftrightarrow}}
\newunicodechar{⟺}{\ensuremath{\Longleftrightarrow}}
\newunicodechar{⇒}{\ensuremath{\Rightarrow}}
\newunicodechar{⟹}{\ensuremath{\Longrightarrow}}
\newunicodechar{∘}{\ensuremath{\circ}}
\newunicodechar{∪}{\ensuremath{\cup}}
\newunicodechar{∩}{\ensuremath{\cap}}
\newunicodechar{≡}{\ensuremath{\equiv}}
\newunicodechar{⊂}{\ensuremath{\subset}}
\newunicodechar{⊥}{\ensuremath{\perp}}
\newunicodechar{∃}{\ensuremath{\exists}}
\newunicodechar{∀}{\ensuremath{\forall}}
\newunicodechar{∛}{\ensuremath{\sqrt[3]{\,}}}
\newunicodechar{∜}{\ensuremath{\sqrt[4]{\,}}}
\newunicodechar{⃗}{\ensuremath{{}^{\to}}}
\newunicodechar{ⁿ}{\ifmmode^{n}\else\textsuperscript{\ensuremath{n}}\fi}
\newunicodechar{ˣ}{\ifmmode^{x}\else\textsuperscript{\ensuremath{x}}\fi}
\newunicodechar{ᵇ}{\ifmmode^{b}\else\textsuperscript{\ensuremath{b}}\fi}
\newunicodechar{ʳ}{\ifmmode^{r}\else\textsuperscript{\ensuremath{r}}\fi}
\newunicodechar{ᵘ}{\ifmmode^{u}\else\textsuperscript{\ensuremath{u}}\fi}
\newunicodechar{ʸ}{\ifmmode^{y}\else\textsuperscript{\ensuremath{y}}\fi}
\newunicodechar{ᵃ}{\ifmmode^{a}\else\textsuperscript{\ensuremath{a}}\fi}
\newunicodechar{ᵉ}{\ifmmode^{e}\else\textsuperscript{\ensuremath{e}}\fi}
\newunicodechar{ᵏ}{\ifmmode^{k}\else\textsuperscript{\ensuremath{k}}\fi}
\newunicodechar{ᵖ}{\ifmmode^{p}\else\textsuperscript{\ensuremath{p}}\fi}
\newunicodechar{ᴺ}{\ifmmode^{N}\else\textsuperscript{\ensuremath{N}}\fi}
\newunicodechar{ᶜ}{\ifmmode^{c}\else\textsuperscript{\ensuremath{c}}\fi}
\newunicodechar{ʰ}{\ifmmode^{h}\else\textsuperscript{\ensuremath{h}}\fi}
\newunicodechar{ᵠ}{\ifmmode^{q}\else\textsuperscript{\ensuremath{q}}\fi}
\newunicodechar{ₙ}{\ifmmode_{n}\else\textsubscript{\ensuremath{n}}\fi}
\newunicodechar{ₐ}{\ifmmode_{a}\else\textsubscript{\ensuremath{a}}\fi}
\newunicodechar{ᵢ}{\ifmmode_{i}\else\textsubscript{\ensuremath{i}}\fi}
\newunicodechar{ᵣ}{\ifmmode_{r}\else\textsubscript{\ensuremath{r}}\fi}
\newunicodechar{ₖ}{\ifmmode_{k}\else\textsubscript{\ensuremath{k}}\fi}
\newunicodechar{ᵦ}{\ifmmode_{\beta}\else\textsubscript{\ensuremath{\beta}}\fi}
\newunicodechar{ₓ}{\ifmmode_{x}\else\textsubscript{\ensuremath{x}}\fi}
\newfontfamily\popdisplay{ArchivoBlack-Regular.ttf}[Path=../../../fonts/]
\newfontfamily\poplogo{SpaceGrotesk-Bold.ttf}[Path=../../../fonts/]
\newfontfamily\popsemi{PlusJakartaSans-SemiBold.ttf}[Path=../../../fonts/]
\newfontfamily\popxbold{PlusJakartaSans-ExtraBold.ttf}[Path=../../../fonts/]
\newfontfamily\popxboldls{PlusJakartaSans-ExtraBold.ttf}[Path=../../../fonts/,LetterSpace=3.0]

\setlist[enumerate]{leftmargin=1.7em,itemsep=5pt,topsep=4pt}
\setlist[itemize]{leftmargin=1.7em,itemsep=4pt,topsep=4pt,label={\color{poporange}\textbullet}}
\setlength{\parindent}{0pt}
\setlength{\parskip}{5pt}
\renewcommand{\arraystretch}{1.25}

% pgfplots : style Pop par défaut
\pgfplotsset{
  every axis/.append style={axis line style={popink,line width=1pt},tick style={popink},
    grid style={popline,line width=0.5pt},label style={font=\small},tick label style={font=\footnotesize}},
  popcurve/.style={poporange,line width=1.8pt,no markers,smooth},
}
% Même style disponible dans un \draw TikZ simple (hors pgfplots)
\tikzset{popcurve/.style={poporange,line width=1.8pt}}

% ============================================================
% Pill / squiggle
% ============================================================
\newcommand{\poppill}[3]{%
  \tikz[baseline=-0.55ex]\node[draw=popink,line width=1.2pt,rounded corners=2.6pt,%
    fill=#2,inner xsep=6.5pt,inner ysep=3pt]{%
    \popxboldls\fontsize{7}{7}\selectfont\color{#3}\MakeUppercase{#1}};%
}
\newcommand{\popsquiggle}[1]{%
  \tikz[baseline=-0.4ex]{
    \draw[#1,line width=2.6pt,line cap=round]
      plot [smooth] coordinates {(0,0.09) (0.34,0.01) (0.68,0.21) (1.02,0.01) (1.36,0.10)};
  }%
}
% Mot-clé surligné (marqueur orange sous le texte)
\newcommand{\cle}[1]{{\popxbold\color{popdark}#1}}

% ============================================================
% En-tête / pied
% ============================================================
\pagestyle{fancy}
\fancyhf{}
\renewcommand{\headrulewidth}{0pt}
\renewcommand{\footrulewidth}{0pt}
\lhead{\raisebox{-0.15ex}{\poplogo\fontsize{13}{13}\selectfont\color{popink}OnBuch\color{poporange}+}}
\rhead{\poppill{\DOCMATIERE\ \textbullet\ \DOCNIVEAU\ \textbullet\ Leçon \DOCLECON}{white}{popink}}
\lfoot{\popxbold\fontsize{7.6}{7.6}\selectfont\color{popmuted}\MakeUppercase{Cours signé OnBuch+}\ \textbullet\ {\popsemi\DOCTITRE}}
\rfoot{\tikz[baseline=-0.6ex]\node[rounded corners=8.5pt,fill=popink,minimum height=17pt,minimum width=17pt,inner xsep=5pt,inner sep=0]{\popxbold\fontsize{8}{8}\selectfont\color{popcream}\thepage\,/\,\pageref*{LastPage}};}

% ============================================================
% Titres
% ============================================================
\newcommand{\chaptitre}[2]{%
  \begin{tcolorbox}[enhanced,colback=poporange,colframe=popink,boxrule=2.6pt,arc=11pt,
      drop shadow=popink,left=15pt,right=15pt,top=12pt,bottom=11pt,before skip=3pt,after skip=18pt]
    \poppill{#2}{popink}{popcream}\par\vspace{9pt}
    {\raggedright\hyphenpenalty=10000\exhyphenpenalty=10000\popdisplay\fontsize{22}{24}\selectfont\color{popcream}\MakeUppercase{#1}\par}
  \end{tcolorbox}
}
% Section numérotée : pastille orange avec le numéro + titre Archivo
\newcounter{popsec}
\newcommand{\coursec}[1]{%
  \stepcounter{popsec}\setcounter{popsub}{0}%
  \par\vspace{16pt}\noindent
  \begin{minipage}{\linewidth}
  \tikz[baseline=-0.7ex]\node[draw=popink,line width=1.6pt,fill=poporange,rounded corners=4pt,minimum size=22pt,inner sep=0]{\popdisplay\fontsize{12}{12}\selectfont\color{popcream}\thepopsec};%
  \hspace{8pt}{\popdisplay\fontsize{15}{17}\selectfont\color{popink}\MakeUppercase{#1}}\par
  \vspace{4pt}\noindent{\color{poporange}\rule{36pt}{3.6pt}}
  \end{minipage}\par\nopagebreak\vspace{8pt}
}
\newcounter{popsub}
\newcommand{\courssub}[1]{%
  \stepcounter{popsub}%
  \par\vspace{9pt}\noindent{\popxbold\fontsize{11.5}{13}\selectfont\color{popink}{\color{poporange}\thepopsec.\thepopsub}\hspace{6pt}#1}\par\nopagebreak\vspace{4pt}
}

% ============================================================
% Boîtes pédagogiques — même recette : fond blanc, bordure épaisse colorée,
% ombre dure encre, badge plein. Titre optionnel [#1].
% ============================================================
\tcbset{popbase/.style={breakable,enhanced,colback=white,boxrule=2.3pt,arc=9pt,
  shadow={3pt}{-3pt}{0pt}{popink},left=14pt,right=14pt,top=11pt,bottom=11pt,before skip=11pt,after skip=15pt,
  fontupper=\popsemi\fontsize{10.6}{14.5}\selectfont\color{popink},
  fontlower=\popsemi\fontsize{10.6}{14.5}\selectfont\color{popink}}}
% Coche dessinée (le glyphe ✓ n'existe pas dans les polices du document)
\newcommand{\popcheck}[1]{\tikz[baseline=-0.5ex]\draw[#1,line width=1.6pt,line cap=round,line join=round](-0.13,0.0)--(-0.04,-0.09)--(0.13,0.1);}
\newcommand{\popboxhead}[3]{\poppill{#1}{#2}{white}\ifx\relax#3\relax\else\hspace{6pt}{\popxbold\fontsize{10.6}{12}\selectfont\color{popink}#3}\fi\par\vspace{7pt}}

\newtcolorbox{aretenir}[1][]{popbase,colframe=popgreen,before upper={\popboxhead{À retenir}{popgreen}{#1}}}
\newtcolorbox{exemplebox}[1][]{popbase,colframe=poporange,before upper={\popboxhead{Exemple}{poporange}{#1}}}
\newtcolorbox{definition}[1][]{popbase,colframe=popblue,before upper={\popboxhead{Définition}{popblue}{#1}}}
\newtcolorbox{propriete}[1][]{popbase,colframe=poppurple,before upper={\popboxhead{Propriété}{poppurple}{#1}}}
\newtcolorbox{methode}[1][]{popbase,colframe=popgold,before upper={\popboxhead{Méthode}{popgold}{#1}}}
\newtcolorbox{attention}[1][]{popbase,colframe=poppink,before upper={\popboxhead{Attention}{poppink}{#1}}}
\newtcolorbox{experience}[1][]{popbase,colframe=popblue,colback=popblueL,before upper={\popboxhead{Expérience}{popblue}{#1}}}
% « Le savais-tu ? » : carte violette pleine (contexte, culture, Cameroun)
\newtcolorbox{savaistu}[1][]{popbase,colback=poppurple,colframe=popink,
  fontupper=\popsemi\fontsize{10.4}{14.2}\selectfont\color{white},
  before upper={\poppill{Le savais-tu\,?}{white}{poppurple}\ifx\relax#1\relax\else\hspace{6pt}{\popxbold\color{white}#1}\fi\par\vspace{7pt}}}
% Exercice résolu : énoncé en haut, \tcblower puis la solution sur fond crème
\newcounter{popexo}\newcounter{popexr}
\newtcolorbox{exoresolu}[1][]{popbase,colframe=poporange,colbacklower=popcream,
  segmentation style={popink,line width=1.2pt,dashed},
  before upper={\stepcounter{popexr}\popboxhead{Exercice résolu \thepopexr}{poporange}{#1}},
  before lower={\poppill{Solution}{popink}{popcream}\par\vspace{6pt}}}
% Exercice (non résolu en place) — la correction est dans la partie Corrigés
\newtcolorbox{exercice}[1][]{popbase,colframe=popink,
  before upper={\stepcounter{popexo}\popboxhead{Exercice \thepopexo}{popink}{#1}}}
% Corrigé
\newtcolorbox{corrige}[1][]{popbase,colframe=popgreen,colback=popgreenL,
  before upper={\popboxhead{Corrigé}{popgreen}{#1}}}
% Objectifs / prérequis / bilan
\newtcolorbox{objectifs}{popbase,colframe=popink,colback=poporangeL,
  before upper={\popboxhead{Objectifs de la leçon}{popink}{}}}
\newtcolorbox{prerequis}{popbase,colframe=popink,colback=popblueL,
  before upper={\popboxhead{Prérequis}{popblue}{}}}
\newtcolorbox{bilan}{popbase,colframe=popink,colback=white,boxrule=2.8pt,
  before upper={\popboxhead{Fiche bilan}{poporange}{L'essentiel à connaître pour l'examen}}}
% Figure encadrée avec légende
\newtcolorbox{popfigure}[1][]{enhanced,breakable=false,colback=white,colframe=popline,boxrule=1.4pt,arc=8pt,
  left=10pt,right=10pt,top=10pt,bottom=8pt,before skip=10pt,after skip=12pt,halign=center,
  lower separated=false,halign lower=center,
  fontlower=\popsemi\fontsize{9}{11.5}\selectfont\color{popmuted},
  #1}
\newcommand{\legende}[1]{\par\vspace{6pt}{\popsemi\fontsize{9}{11.5}\selectfont\color{popmuted}#1}}

% ============================================================
% Signature OnBuch+
% ============================================================
\newcommand{\poplogoinline}[1]{{\poplogo\fontsize{#1}{#1}\selectfont\color{popink}OnBuch\color{poporange}+}}
\newcommand{\signatureonbuch}{%
  \begin{tcolorbox}[enhanced,colback=popink,colframe=popink,boxrule=2.2pt,arc=10pt,
      drop shadow=poporange,left=14pt,right=14pt,top=10pt,bottom=10pt,before skip=0pt,after skip=0pt]
    \tikz[baseline=-0.7ex]\node[circle,fill=popgreen,draw=popcream,line width=1.3pt,minimum size=22pt,inner sep=0]{\popcheck{white}};%
    \hspace{9pt}\begin{minipage}[c]{0.62\linewidth}
      {\popxbold\fontsize{12}{13}\selectfont\color{popcream}Signé\ \textbullet\ {\poplogo OnBuch\color{poporange}+}}\par\vspace{2pt}
      {\raggedright\popsemi\fontsize{8}{10.5}\selectfont\color{popline}\MakeUppercase{Équipe pédagogique OnBuch+}\\\MakeUppercase{Conforme au programme officiel MINESEC}\par}
    \end{minipage}\hfill
    {\poplogo\fontsize{22}{22}\selectfont\color{popcream}OnBuch\color{poporange}+}
  \end{tcolorbox}%
}

% ============================================================
% Page de garde
% #1 titre · #2 matière · #3 niveau · #4 module · #5 n° leçon · #6 durée ·
% #7 description / objectifs (texte ou itemize)
% ============================================================
\newcommand{\pagegarde}[7]{%
  \thispagestyle{empty}
  \newgeometry{a4paper,margin=1.7cm,top=1.6cm,bottom=1.4cm}%
  \begin{tikzpicture}[remember picture,overlay]
    \fill[poporange!18] ([xshift=-3.2cm,yshift=-3.6cm]current page.north east) circle (4.6cm);
    \fill[poppurple!14] ([xshift=2.4cm,yshift=7.6cm]current page.south west) circle (3.8cm);
  \end{tikzpicture}%
  {\poplogo\fontsize{30}{30}\selectfont\color{popink}OnBuch\color{poporange}+}\hfill
  \raisebox{0.6ex}{\poppill{Cours signé \textbullet\ Premium}{popink}{popcream}}\par
  \vspace{1pt}\popsquiggle{poppurple}\par
  \vspace{8pt}
  \begin{tcolorbox}[enhanced,colback=poporange,colframe=popink,boxrule=2.8pt,arc=13pt,
      drop shadow=popink,left=18pt,right=18pt,top=12pt,bottom=13pt,after skip=10pt]
    \poppill{#2 \textbullet\ #3}{popink}{popcream}\hspace{5pt}\poppill{#4}{white}{popink}\par\vspace{8pt}
    {\popdisplay\fontsize{14}{14}\selectfont\color{popink}LEÇON #5}\par\vspace{4pt}
    {\raggedright\hyphenpenalty=10000\exhyphenpenalty=10000\popdisplay\fontsize{25}{27}\selectfont\color{popcream}\MakeUppercase{#1}\par}
    \vspace{8pt}
    {\popxbold\fontsize{9.5}{9.5}\selectfont\color{popink}\MakeUppercase{Durée conseillée : #6}}
  \end{tcolorbox}
  \begin{tcolorbox}[enhanced,colback=white,colframe=popink,boxrule=2.2pt,arc=10pt,
      drop shadow=popink,left=16pt,right=16pt,top=10pt,bottom=10pt,after skip=8pt]
    \poppill{Dans cette leçon}{poppurple}{white}\par\vspace{5pt}
    \popsemi\fontsize{9.3}{12.6}\selectfont\color{popink}#7
  \end{tcolorbox}
  \noindent\begin{minipage}[t]{0.487\linewidth}
    \begin{tcolorbox}[enhanced,colback=popgreen,colframe=popink,boxrule=2.2pt,arc=10pt,
        drop shadow=popink,left=11pt,right=11pt,top=10pt,bottom=10pt,height=3.1cm,valign=center]
      \tcbox[colback=white,colframe=popink,boxrule=1.3pt,arc=5pt,boxsep=0pt,nobeforeafter,
        left=3pt,right=3pt,top=3pt,bottom=3pt,valign=center]{\qrcode[height=1.9cm]{https://play.google.com/store/apps/details?id=com.luvvix.onbuch&hl=fr}}%
      \hspace{8pt}\begin{minipage}[c]{0.5\linewidth}
        {\popdisplay\fontsize{11}{12.5}\selectfont\color{white}\MakeUppercase{Télécharge OnBuch}}\par\vspace{4pt}
        {\raggedright\popsemi\fontsize{8.4}{11}\selectfont\color{white}Gratuit \textbullet\ Google Play\\Tuteur IA, annales, concours, résultats\par}
      \end{minipage}
    \end{tcolorbox}
  \end{minipage}\hfill
  \begin{minipage}[t]{0.487\linewidth}
    \begin{tcolorbox}[enhanced,colback=poppurple,colframe=popink,boxrule=2.2pt,arc=10pt,
        drop shadow=popink,left=11pt,right=11pt,top=10pt,bottom=10pt,height=3.1cm,valign=center]
      \tikz[baseline=-0.7ex]\node[circle,fill=white,draw=popink,line width=1.3pt,minimum size=1.6cm,inner sep=0]{\popdisplay\fontsize{24}{24}\selectfont\color{poppurple}+};%
      \hspace{8pt}\begin{minipage}[c]{0.6\linewidth}
        {\popdisplay\fontsize{11}{12.5}\selectfont\color{white}\MakeUppercase{Passe à OnBuch+}}\par\vspace{4pt}
        {\raggedright\popsemi\fontsize{8.4}{11}\selectfont\color{white}Tous les cours signés, Tuteur IA illimité, fiches PDF — onglet OnBuch+ de l'app\par}
      \end{minipage}
    \end{tcolorbox}
  \end{minipage}
  \vspace{8pt}
  \popcontenu
  \vfill
  \signatureonbuch
  \restoregeometry
  \newpage
}

% Rangée « Ce cours contient » (page de garde)
\newcommand{\poptile}[3]{% #1 couleur #2 gros texte #3 légende
  \begin{tcolorbox}[enhanced,colback=#1,colframe=popink,boxrule=2pt,arc=9pt,shadow={2.5pt}{-2.5pt}{0pt}{popink},
      left=6pt,right=6pt,top=9pt,bottom=9pt,halign=center,valign=center,height=1.75cm,width=\linewidth,nobeforeafter]
    {\popdisplay\fontsize{12.5}{13}\selectfont\color{white}\MakeUppercase{#2}}\par\vspace{3pt}
    {\centering\popsemi\fontsize{7.8}{9.5}\selectfont\color{white}#3\par}
  \end{tcolorbox}}
\newcommand{\popcontenu}{%
  \par\noindent{\popxboldls\fontsize{7.5}{7.5}\selectfont\color{popmuted}CE COURS SIGNÉ CONTIENT}\par\vspace{7pt}
  \noindent\begin{minipage}[t]{0.235\linewidth}\poptile{popblue}{Cours}{complet, illustré, pas à pas}\end{minipage}\hfill
  \begin{minipage}[t]{0.235\linewidth}\poptile{poporange}{Méthodes}{et exercices résolus}\end{minipage}\hfill
  \begin{minipage}[t]{0.235\linewidth}\poptile{poppink}{Exercices}{type examen + corrigés}\end{minipage}\hfill
  \begin{minipage}[t]{0.235\linewidth}\poptile{popgreen}{Fiche bilan}{l'essentiel à retenir}\end{minipage}\par}

% Sommaire
\newtcolorbox{sommaire}{enhanced,colback=white,colframe=popink,boxrule=2.2pt,arc=10pt,
  drop shadow=popink,left=18pt,right=18pt,top=13pt,bottom=13pt,before skip=6pt,after skip=20pt,
  before upper={\poppill{Sommaire}{poppurple}{white}\par\vspace{9pt}\popsemi\fontsize{10.6}{14.5}\selectfont\color{popink}}}

% Signature de fin de cours — poussée tout en bas de la dernière page
\newcommand{\findecours}{%
  \par\vfill\nopagebreak
  \begin{center}\popsquiggle{poporange}\end{center}\vspace{4pt}
  \signatureonbuch
}
\AtBeginDocument{\def\llbracket{\mathopen{[\mkern-3.5mu[}}\def\rrbracket{\mathclose{]\mkern-3.5mu]}}} % intervalles d'entiers (glyphe ⟦ absent de la police)
% Glyphes absents de Fira Math : redéfinis à partir de glyphes disponibles
\AtBeginDocument{%
  \def\setminus{\mathbin{\backslash}}\let\smallsetminus\setminus
  \def\vdots{\mathord{\vbox{\baselineskip3.2pt\lineskiplimit0pt\kern4pt\hbox{.}\hbox{.}\hbox{.}}}}%
  \def\ddots{\mathinner{\mkern1mu\raise6.5pt\hbox{.}\mkern2mu\raise3.6pt\hbox{.}\mkern2mu\raise0.7pt\hbox{.}\mkern1mu}}%
  \def\triangle{\mathord{\scalebox{0.9}{$\symup{\Delta}$}}}%
  \def\nabla{\mathord{\raisebox{\depth}{\scalebox{1}[-1]{$\symup{\Delta}$}}}}%
  \def\checkmark{\popcheck{popdark}}%
  \def\star{\mathbin{\ast}}\def\bigstar{\mathord{\ast}}%
  \def\diamond{\mathbin{\scalebox{0.6}{$\Diamond$}}}%
  \def\bigcup{\mathop{\mathchoice{\raisebox{-0.3em}{\scalebox{1.7}{$\cup$}}}{\scalebox{1.25}{$\cup$}}{\cup}{\cup}}\displaylimits}%
  \def\bigcap{\mathop{\mathchoice{\raisebox{-0.3em}{\scalebox{1.7}{$\cap$}}}{\scalebox{1.25}{$\cap$}}{\cap}{\cap}}\displaylimits}%
  \def\lesseqqgtr{\mathrel{\lessgtr}}\def\bigoplus{\mathop{\oplus}\displaylimits}%
}
\providecommand{\ssi}{si et seulement si} % abréviation fréquente
\AtBeginDocument{\providecommand{\wideparen}{\overparen}} % arc de cercle

%% FIGFIX-BEGIN v5 — correctifs globaux des figures (docs/onbuch-generation/tools/figfix)
\usepackage{adjustbox}\usepackage{etoolbox}\tcbuselibrary{raster}
% toute figure TikZ/pgfplots plus large que la ligne est réduite à la largeur disponible
\BeforeBeginEnvironment{tikzpicture}{\begin{adjustbox}{max width=\linewidth}}
\AfterEndEnvironment{tikzpicture}{\end{adjustbox}}
% légendes pgfplots lisibles par-dessus les courbes
\pgfplotsset{every axis/.append style={legend style={fill=white,fill opacity=0.92,text opacity=1,draw=black!15,inner sep=3pt}}}
% étiquettes d'arêtes (syntaxe edge["..."]) avec fond blanc
\tikzset{every edge quotes/.append style={fill=white,inner sep=1.2pt}}
%% FIGFIX-END
\begin{document}

\pagegarde{Arithmétique}{Mathématiques}{3ème}{Mathématiques – classe de 3ème}{N01}{3 séances (fiche de progression harmonisée)}{Cette leçon approfondit l'arithmétique des entiers naturels : recherche du PGCD par les listes de diviseurs, par l'algorithme des soustractions successives et par l'algorithme d'Euclide, puis étude du PPCM et de la relation fondamentale PGCD × PPCM = a × b. Les notions sont systématiquement appliquées à des situations concrètes de la vie courante camerounaise (partages, pavages, rendez-vous périodiques).
\vspace{4pt}
\begin{multicols}{2}\small
\begin{itemize}
\item À la fin de cette leçon, je sais déterminer l'ensemble des diviseurs communs de deux entiers naturels
\item À la fin de cette leçon, je sais calculer le PGCD de deux entiers par la méthode des listes de diviseurs
\item À la fin de cette leçon, je sais calculer le PGCD par l'algorithme des soustractions successives
\item À la fin de cette leçon, je sais calculer le PGCD par l'algorithme d'Euclide et choisir la méthode la plus efficace
\item À la fin de cette leçon, je sais reconnaître et justifier que deux entiers sont premiers entre eux
\item À la fin de cette leçon, je sais déterminer le PPCM de deux entiers naturels
\end{itemize}\end{multicols}}

\begin{sommaire}
\begin{multicols}{2}
\begin{enumerate}
\item Rappels : diviseurs, multiples et division euclidienne
\item PGCD : définition et méthode des listes
\item Algorithme des soustractions successives
\item Algorithme d'Euclide
\item Nombres premiers entre eux
\item PPCM et relation avec le PGCD
\item Résolution de problèmes concrets et synthèse
\item Activité d'intégration
\item Méthodes et exercices résolus
\item Exercices
\item Corrigés des exercices
\item Fiche bilan
\end{enumerate}
\end{multicols}
\end{sommaire}

\begin{objectifs}
\begin{itemize}
\item À la fin de cette leçon, je sais déterminer l'ensemble des diviseurs communs de deux entiers naturels
\item À la fin de cette leçon, je sais calculer le PGCD de deux entiers par la méthode des listes de diviseurs
\item À la fin de cette leçon, je sais calculer le PGCD par l'algorithme des soustractions successives
\item À la fin de cette leçon, je sais calculer le PGCD par l'algorithme d'Euclide et choisir la méthode la plus efficace
\item À la fin de cette leçon, je sais reconnaître et justifier que deux entiers sont premiers entre eux
\item À la fin de cette leçon, je sais déterminer le PPCM de deux entiers naturels
\item À la fin de cette leçon, je sais utiliser la relation PGCD × PPCM = a × b pour trouver l'un connaissant l'autre
\item À la fin de cette leçon, je sais modéliser et résoudre un problème concret de partage ou de périodicité à l'aide du PGCD ou du PPCM
\item À la fin de cette leçon, je sais rédiger et justifier clairement ma démarche et ma conclusion
\end{itemize}
\end{objectifs}

\begin{prerequis}
\begin{itemize}
\item Division euclidienne : dividende, diviseur, quotient, reste (6ème-5ème)
\item Diviseurs et multiples d'un entier naturel ; critères de divisibilité (6ème)
\item Nombres premiers et décomposition en produit de facteurs premiers (4ème)
\item Notion de PGCD et de PPCM vus en 4ème avec les décompositions en facteurs premiers
\end{itemize}
\end{prerequis}

\newpage

\coursec{Rappels : diviseurs, multiples et division euclidienne}

\courssub{Une situation pour ouvrir le chapitre}

C'est la veille de la fête de fin d'année de l'école publique du quartier. Maman Aïcha a préparé \textbf{84 beignets} et \textbf{126 cacahuètes grillées}. Elle veut les répartir dans des sachets \emph{tous identiques} — chaque sachet doit contenir le même nombre de beignets et le même nombre de cacahuètes — et surtout, elle ne veut \textbf{rien laisser de côté}. Sa question est simple : \emph{quel est le plus grand nombre de sachets qu'elle peut préparer ?}

Pendant ce temps, à la gare routière de Mvan, le chef de gare organise les départs. Les cars à destination de \textbf{Bafoussam} partent toutes les \textbf{12 minutes}, et ceux à destination de \textbf{Bafang} toutes les \textbf{18 minutes}. Ce matin, à 6 h précises, deux cars sont partis en même temps, un pour chaque destination. Le chef de gare se demande : \emph{à quelle heure deux cars partiront-ils de nouveau ensemble ?}

Ces deux problèmes semblent n'avoir rien en commun : l'un parle de beignets, l'autre de cars. Pourtant, ils se résolvent avec les \emph{mêmes outils} de l'arithmétique : le \cle{PGCD} (plus grand commun diviseur) pour Maman Aïcha, et le \cle{PPCM} (plus petit commun multiple) pour le chef de gare. Avant de découvrir ces outils, nous devons réactiver trois notions que tu connais déjà depuis la classe de 6\textsuperscript{e} : les \textbf{diviseurs}, les \textbf{multiples} et la \textbf{division euclidienne}. C'est l'objet de cette première section.

\begin{attention}[Problématique de la leçon]
Comment déterminer efficacement le plus grand diviseur commun (PGCD) et le plus petit multiple commun (PPCM) de deux entiers naturels, et comment les utiliser pour résoudre des problèmes concrets de partage et de coïncidence ?
\end{attention}

\courssub{Diviseurs et multiples d'un entier naturel}

\textbf{Intuition.} Dire qu'un nombre \emph{divise} un autre, c'est dire qu'on peut le partager en parts égales de cette taille \emph{sans reste}. Par exemple, Maman Aïcha peut répartir ses 84 beignets dans 7 sachets de 12 beignets exactement : on dit que 7 divise 84 (et que 12 divise 84 aussi).

\begin{definition}[Diviseur et multiple]
Soient $a$ et $b$ deux entiers naturels, avec $a \neq 0$. On dit que $a$ \cle{divise} $b$ (ou que $a$ est un \cle{diviseur} de $b$, ou encore que $b$ est un \cle{multiple} de $a$) lorsqu'il existe un entier naturel $k$ tel que :
\[ b = a \times k. \]
On note parfois $a \mid b$, qui se lit « $a$ divise $b$ ».
\end{definition}

\begin{exemplebox}[Lecture de la définition sur des exemples]
\begin{itemize}
\item $84 = 7 \times 12$, donc $7$ divise $84$ et $12$ divise $84$. On dit aussi que $84$ est un multiple de $7$ et de $12$.
\item $126 = 9 \times 14$, donc $9$ et $14$ sont des diviseurs de $126$.
\item $5$ ne divise pas $84$ : en effet, $84 = 5 \times 16 + 4$, il reste $4$ ; aucun entier $k$ ne vérifie $84 = 5 \times k$.
\end{itemize}
\end{exemplebox}

\begin{attention}[Ne pas inverser le sens !]
Erreur très fréquente au BEPC : confondre « $a$ divise $b$ » et « $a$ est divisible par $b$ ».
\begin{itemize}
\item « $7$ divise $84$ » signifie $84 = 7 \times k$ : c'est \textbf{vrai} ($k = 12$).
\item « $7$ est divisible par $84$ » signifierait $7 = 84 \times k$ : c'est \textbf{faux} pour un entier naturel $k$.
\end{itemize}
Retiens : dans « $a$ divise $b$ », c'est le \emph{petit} qui divise le \emph{grand} (sauf cas d'égalité). Le diviseur est plus petit ou égal au nombre ; le multiple est plus grand ou égal.
\end{attention}

\textbf{Deux remarques utiles.}
\begin{itemize}
\item Tout entier naturel $n \geq 1$ admet toujours au moins deux diviseurs : $1$ et lui-même, car $n = 1 \times n$.
\item La liste des \emph{multiples} d'un entier $a \neq 0$ est \textbf{infinie} : ce sont $0, a, 2a, 3a, \dots$ Par exemple, les multiples de $12$ sont $0 ; 12 ; 24 ; 36 ; 48 ; 60 ; 72 ; 84 ; \dots$ — on ne s'arrête jamais. C'est exactement la liste des horaires de départ des cars de Bafoussam (en minutes après 6 h) !
\end{itemize}

\courssub{Comment lister sans oubli tous les diviseurs d'un entier}

Trouver \emph{tous} les diviseurs d'un nombre comme 84 sans en oublier un seul demande une méthode. L'idée repose sur une observation simple : les diviseurs viennent \textbf{par paires}. Si $a$ divise $n$, alors $n = a \times k$, et donc $k$ divise aussi $n$. Le diviseur $a$ et son « compagnon » $k = \dfrac{n}{a}$ se correspondent.

\begin{methode}[Lister tous les diviseurs d'un entier $n$ par paires]
\begin{enumerate}
\item On teste les entiers $1, 2, 3, 4, \dots$ dans l'ordre, en utilisant les critères de divisibilité (par 2, 3, 5, 9, 10) et en posant les divisions.
\item Chaque fois qu'un entier $a$ divise $n$, on écrit \textbf{les deux diviseurs} de la paire : $a$ et $\dfrac{n}{a}$.
\item On s'arrête dès que le nombre testé dépasse $\sqrt{n}$ (c'est-à-dire dès que les deux membres de la paire « se croisent »).
\item On range ensuite tous les diviseurs trouvés dans l'ordre croissant.
\end{enumerate}
\end{methode}

\begin{exemplebox}[Les diviseurs de 84, pas à pas]
Appliquons la méthode à $n = 84$ (les beignets de Maman Aïcha) :
\begin{itemize}
\item $84 = 1 \times 84$ : paire $(1 ; 84)$ ;
\item $84 = 2 \times 42$ : paire $(2 ; 42)$ ;
\item $84 = 3 \times 28$ : paire $(3 ; 28)$ (car $8 + 4 = 12$ est divisible par 3) ;
\item $84 = 4 \times 21$ : paire $(4 ; 21)$ ;
\item $84 = 6 \times 14$ : paire $(6 ; 14)$ ;
\item $84 = 7 \times 12$ : paire $(7 ; 12)$.
\end{itemize}
Le nombre suivant à tester serait $8$, $9$, $10$, $11$ : aucun ne divise $84$, et à partir de $10$ on dépasse $\sqrt{84} \approx 9,2$ : on s'arrête. Conclusion :
\[ \text{Diviseurs de } 84 = \{1 ; 2 ; 3 ; 4 ; 6 ; 7 ; 12 ; 14 ; 21 ; 28 ; 42 ; 84\}. \]
On compte \textbf{12 diviseurs}, soit 6 paires.
\end{exemplebox}

\begin{exoresolu}[Diviseurs de 126 et diviseurs communs à 84 et 126]
\textbf{Énoncé.} (a) Déterminer la liste de tous les diviseurs de $126$. (b) En déduire la liste des diviseurs communs à $84$ et $126$, puis leur plus grand diviseur commun.
\tcblower
\textbf{Solution.} (a) On applique la méthode des paires à $n = 126$ :
\begin{itemize}
\item $126 = 1 \times 126$ : paire $(1 ; 126)$ ;
\item $126 = 2 \times 63$ : paire $(2 ; 63)$ ;
\item $126 = 3 \times 42$ : paire $(3 ; 42)$ (car $1 + 2 + 6 = 9$ est divisible par 3) ;
\item $126 = 6 \times 21$ : paire $(6 ; 21)$ ;
\item $126 = 7 \times 18$ : paire $(7 ; 18)$ ;
\item $126 = 9 \times 14$ : paire $(9 ; 14)$.
\end{itemize}
On s'arrête car $\sqrt{126} \approx 11,2$ et ni $10$ ni $11$ ne divisent $126$. Donc :
\[ \text{Diviseurs de } 126 = \{1 ; 2 ; 3 ; 6 ; 7 ; 9 ; 14 ; 18 ; 21 ; 42 ; 63 ; 126\}. \]
(b) On compare avec les diviseurs de $84$ : les diviseurs \textbf{communs} sont
\[ \{1 ; 2 ; 3 ; 6 ; 7 ; 14 ; 21 ; 42\}. \]
Le plus grand d'entre eux est $\boxed{42}$ : c'est le PGCD de 84 et 126. Maman Aïcha peut donc préparer au maximum \textbf{42 sachets}, chacun contenant $84 \div 42 = 2$ beignets et $126 \div 42 = 3$ cacahuètes. Le problème de l'accroche est déjà résolu... mais la méthode des listes devient très longue pour de grands nombres : c'est pourquoi nous découvrirons dans les sections suivantes deux algorithmes bien plus rapides.
\end{exoresolu}

Le tableau suivant récapitule la situation : les diviseurs communs aux deux nombres y sont surlignés.

\begin{popfigure}
\begin{center}
\renewcommand{\arraystretch}{1.3}
\begin{tabular}{|c|*{12}{>{\centering\arraybackslash}p{0.55cm}|}}
\hline
\rowcolor{popblueL} $84$ & \cellcolor{popgoldL}1 & \cellcolor{popgoldL}2 & \cellcolor{popgoldL}3 & 4 & \cellcolor{popgoldL}6 & \cellcolor{popgoldL}7 & 12 & \cellcolor{popgoldL}14 & \cellcolor{popgoldL}21 & 28 & \cellcolor{popgoldL}42 & 84 \\
\hline
\rowcolor{popblueL} $126$ & \cellcolor{popgoldL}1 & \cellcolor{popgoldL}2 & \cellcolor{popgoldL}3 & \cellcolor{popgoldL}6 & \cellcolor{popgoldL}7 & 9 & \cellcolor{popgoldL}14 & 18 & \cellcolor{popgoldL}21 & \cellcolor{popgoldL}42 & 63 & 126 \\
\hline
\end{tabular}
\end{center}
\legende{Diviseurs de 84 et de 126 rangés en ordre croissant. Les cases dorées indiquent les diviseurs communs : 1 ; 2 ; 3 ; 6 ; 7 ; 14 ; 21 ; 42. Le plus grand est 42.}
\end{popfigure}

\courssub{La division euclidienne}

\textbf{Intuition.} Revenons aux beignets. Si Maman Aïcha voulait répartir ses 84 beignets dans des sachets de 9, combien de sachets remplirait-elle ? On pose la division : $84 = 9 \times 9 + 3$. Elle remplit 9 sachets complets... et il lui reste 3 beignets. Ce « reste » est le cœur de la division euclidienne, et c'est lui qui décide si le partage est possible sans reste.

\begin{propriete}[Division euclidienne]
Soient $a$ un entier naturel et $b$ un entier naturel non nul ($b \neq 0$). Il existe un \textbf{unique} couple d'entiers naturels $(q ; r)$ tel que :
\[ a = b \times q + r \qquad \text{avec} \qquad 0 \leq r < b. \]
L'entier $q$ s'appelle le \cle{quotient}, l'entier $r$ s'appelle le \cle{reste}, $a$ est le \emph{dividende} et $b$ le \emph{diviseur}.
\end{propriete}

La condition $0 \leq r < b$ est \textbf{essentielle} : le reste doit être strictement plus petit que le diviseur. Si tu trouves un reste supérieur ou égal à $b$, c'est que tu peux encore « faire un tour de plus » : le quotient n'est pas le bon.

\begin{exemplebox}[Deux divisions euclidiennes]
\begin{itemize}
\item Division de $84$ par $9$ : $84 = 9 \times 9 + 3$, avec $0 \leq 3 < 9$. Quotient : $9$, reste : $3$.
\item Division de $126$ par $18$ : $126 = 18 \times 7 + 0$, avec $0 \leq 0 < 18$. Quotient : $7$, reste : $0$.
\end{itemize}
\end{exemplebox}

\begin{popfigure}
\begin{center}
\begin{tikzpicture}[scale=0.7,>=Stealth]
% Axe horizontal
\draw[->,popline] (-0.5,0) -- (13.5,0);
% Graduations et labels
\foreach \x/\lab in {0/{$0$}, 2/{$b$}, 4/{$2b$}, 6/{$3b$}, 8/{$qb$}, 10/{$a$}, 12/{$(q+1)b$}} {
  \draw[popdark] (\x,0.1) -- (\x,-0.1) node[below=3pt] {\lab};
}
% Segments b (bleu)
\draw[popblue,thick] (0,0.3) -- (2,0.3) node[midway,above=5pt,popblue] {$b$};
\draw[popblue,thick] (2,0.3) -- (4,0.3) node[midway,above=5pt,popblue] {$b$};
\draw[popmuted] (5,0.5) node {$\cdots$};
\draw[popblue,thick] (6,0.3) -- (8,0.3) node[midway,above=5pt,popblue] {$b$};
% Segment r (orange)
\draw[poporange,thick] (8,0.3) -- (10,0.3) node[midway,above=5pt,poporange] {$r$};
% Segment b suivant (vert) pour montrer r < b
\draw[popgreen,thick] (10,0.3) -- (12,0.3) node[midway,above=5pt,popgreen] {$b$};
% Points clés
\fill[popblue] (8,0) circle (2pt);
\fill[poporange] (10,0) circle (2pt);
% Annotation formule
\node[popdark,align=center] at (10.5,-1.2) {$a = b \times q + r$\\ avec $0 \leq r < b$};
\end{tikzpicture}
\end{center}
\legende{Lecture géométrique de la division euclidienne : sur la droite graduée, $a$ se situe entre deux multiples consécutifs de $b$ : $qb \leq a < (q+1)b$. Le reste $r$ est la distance entre $qb$ et $a$ ; il est strictement plus petit que $b$.}
\end{popfigure}

\begin{aretenir}[Lien fondamental entre division euclidienne et divisibilité]
Dans la division euclidienne de $a$ par $b$ ($b \neq 0$) :
\[ b \text{ divise } a \quad \Longleftrightarrow \quad \text{le reste } r \text{ est nul.} \]
En effet, si $r = 0$, alors $a = b \times q$, donc $b$ divise $a$. Réciproquement, si $b$ divise $a$, alors $a = b \times k$ et l'unicité du quotient et du reste impose $q = k$ et $r = 0$.
\end{aretenir}

Ainsi, dire que Maman Aïcha peut répartir ses 126 cacahuètes dans des sachets de 18 sans reste, c'est exactement dire que $126 = 18 \times 7 + 0$ : le reste est nul, donc 18 divise 126.

\courssub{Une propriété clé pour la suite}

Voici la propriété qui justifiera, dans les sections 3 et 4, les algorithmes de calcul du PGCD. Nous l'admettons pour l'instant ; elle sera démontrée en section 3.

\begin{propriete}[Propriété admise — démontrée en section 3]
Dans la division euclidienne $a = b \times q + r$, tout diviseur commun à $a$ et à $b$ divise aussi le reste $r$. Réciproquement, tout diviseur commun à $b$ et à $r$ divise aussi $a$. Par conséquent, $a$ et $b$ d'une part, $b$ et $r$ d'autre part, ont \emph{exactement les mêmes} diviseurs communs — et donc le même PGCD.
\end{propriete}

\begin{exemplebox}[Vérification sur un exemple]
Prenons $a = 126$ et $b = 84$ : la division euclidienne donne $126 = 84 \times 1 + 42$.
\begin{itemize}
\item $42$ est un diviseur commun à $126$ et $84$ ; il divise bien le reste $42$ (évidemment !).
\item Les diviseurs communs à $84$ et $42$ sont $\{1 ; 2 ; 3 ; 6 ; 7 ; 14 ; 21 ; 42\}$ : exactement les mêmes que les diviseurs communs à $126$ et $84$, trouvés plus haut. La propriété est confirmée.
\end{itemize}
C'est ce phénomène qui rendra l'algorithme d'Euclide si puissant : remplacer le couple $(126 ; 84)$ par le couple $(84 ; 42)$, plus petit, sans changer la réponse.
\end{exemplebox}

\courssub{Multiples communs de deux entiers}

Pour le problème du chef de gare de Mvan, ce sont les \emph{multiples} qui interviennent. Les départs vers Bafoussam ont lieu aux instants (en minutes après 6 h) :
\[ 0 ; 12 ; 24 ; 36 ; 48 ; 60 ; 72 ; 84 ; \dots \quad \text{(multiples de } 12\text{)} \]
et ceux vers Bafang :
\[ 0 ; 18 ; 36 ; 54 ; 72 ; 90 ; \dots \quad \text{(multiples de } 18\text{)}. \]

\begin{definition}[Multiple commun]
Un \cle{multiple commun} à deux entiers naturels non nuls $a$ et $b$ est un entier naturel qui est à la fois un multiple de $a$ et un multiple de $b$.
\end{definition}

\begin{exemplebox}[Multiples communs de 12 et 18]
En comparant les deux listes ci-dessus, les premiers multiples communs de $12$ et $18$ sont :
\[ 0 ; 36 ; 72 ; 108 ; \dots \]
Le plus petit multiple commun \emph{non nul} est $36$ : c'est le PPCM de 12 et 18. Les deux cars repartiront donc ensemble \textbf{36 minutes après 6 h, c'est-à-dire à 6 h 36}, puis à 7 h 12, etc.
\end{exemplebox}

\begin{attention}[Deux listes de nature différente]
Ne confonds pas les deux situations :
\begin{itemize}
\item la liste des \textbf{diviseurs} d'un entier est \emph{finie} (celle de 84 a 12 éléments) ;
\item la liste des \textbf{multiples} d'un entier non nul est \emph{infinie} : on ne peut jamais l'écrire entièrement, on écrit les premiers termes suivis de pointillés.
\end{itemize}
De même, deux entiers ont toujours un \emph{plus grand} diviseur commun (le PGCD) et un \emph{plus petit} multiple commun non nul (le PPCM).
\end{attention}

\courssub{Bilan de la section}

\begin{aretenir}[À retenir de cette section]
\begin{itemize}
\item $a$ divise $b$ (avec $a \neq 0$) signifie : il existe un entier $k$ tel que $b = a \times k$. On dit aussi que $b$ est un multiple de $a$.
\item On liste tous les diviseurs d'un entier \textbf{par paires}, en s'arrêtant à sa racine carrée : diviseurs de $84 = \{1 ; 2 ; 3 ; 4 ; 6 ; 7 ; 12 ; 14 ; 21 ; 28 ; 42 ; 84\}$.
\item Division euclidienne de $a$ par $b \neq 0$ : $a = bq + r$ avec $0 \leq r < b$. Le reste est nul si et seulement si $b$ divise $a$.
\item Tout diviseur commun à $a$ et $b$ divise le reste $r$ : c'est la clé des algorithmes à venir.
\item Les multiples communs de $12$ et $18$ sont $0 ; 36 ; 72 ; \dots$ : le plus petit non nul, $36$, est leur PPCM.
\end{itemize}
\end{aretenir}

\begin{exercice}[Pour vérifier tes acquis]
\begin{enumerate}
\item Dresser la liste complète des diviseurs de $60$, puis celle des diviseurs de $45$. En déduire leurs diviseurs communs et leur PGCD.
\item Écrire la division euclidienne de $237$ par $15$. Le nombre $15$ divise-t-il $237$ ?
\item Un fleuriste de Bonabéri veut composer des bouquets identiques avec $60$ roses et $45$ œillets, sans reste. Quel est le plus grand nombre de bouquets possibles, et quelle sera la composition de chacun ?
\item Écrire les huit premiers multiples non nuls de $15$ et de $20$. Quel est leur plus petit multiple commun non nul ?
\end{enumerate}
\end{exercice}

\begin{corrige}[Exercice « Pour vérifier tes acquis »]
\begin{enumerate}
\item Diviseurs de $60$ : $60 = 1 \times 60 = 2 \times 30 = 3 \times 20 = 4 \times 15 = 5 \times 12 = 6 \times 10$, donc $\{1 ; 2 ; 3 ; 4 ; 5 ; 6 ; 10 ; 12 ; 15 ; 20 ; 30 ; 60\}$. Diviseurs de $45$ : $45 = 1 \times 45 = 3 \times 15 = 5 \times 9$, donc $\{1 ; 3 ; 5 ; 9 ; 15 ; 45\}$. Diviseurs communs : $\{1 ; 3 ; 5 ; 15\}$ ; PGCD : $15$.
\item $237 = 15 \times 15 + 12$, avec $0 \leq 12 < 15$. Le reste n'est pas nul, donc $15$ ne divise pas $237$.
\item Le nombre de bouquets doit diviser à la fois $60$ et $45$ : le plus grand possible est le PGCD, soit $15$ bouquets. Chaque bouquet contient $60 \div 15 = 4$ roses et $45 \div 15 = 3$ œillets.
\item Multiples de $15$ : $15 ; 30 ; 45 ; 60 ; 75 ; 90 ; 105 ; 120$. Multiples de $20$ : $20 ; 40 ; 60 ; 80 ; 100 ; 120 ; 140 ; 160$. Le plus petit multiple commun non nul est $60$.
\end{enumerate}
\end{corrige}

Dans la prochaine section, nous formaliserons la notion de \textbf{PGCD} et nous verrons comment la méthode des listes, déjà utilisée ici, permet de le calculer — avant de découvrir des méthodes beaucoup plus rapides.

\coursec{PGCD : définition et méthode des listes}

\courssub{Transition depuis les rappels}

Dans la section précédente, nous avons révisé la notion de \cle{diviseur} et la \cle{division euclidienne}. Rappelons l'essentiel : un entier naturel $d$ est diviseur d'un entier naturel $a$ s'il existe un entier naturel $q$ tel que $a = d \times q$. Nous savons aussi lister les diviseurs d'un nombre (par exemple $D(12) = \{1; 2; 3; 4; 6; 12\}$) et effectuer une division euclidienne ($a = bq + r$ avec $0 \le r < b$).

Aujourd'hui, nous allons croiser les listes de diviseurs de \textbf{deux} nombres. C'est le point de départ pour résoudre des problèmes de partage équitable, comme celui des sachets de beignets découvert en accroche : \emph{« Comment faire le maximum de sachets identiques avec 84 beignets et 126 cacahuètes ? »}. La réponse se cache dans les diviseurs \textbf{communs} à 84 et 126.

\courssub{Diviseurs communs et intersection d'ensembles}

Soient $a$ et $b$ deux entiers naturels non nuls. L'ensemble des diviseurs de $a$ est noté $D(a)$, celui de $b$ est noté $D(b)$. Les nombres qui divisent \textbf{à la fois} $a$ et $b$ forment l'intersection de ces deux ensembles.

\begin{definition}[Diviseurs communs]
Les \cle{diviseurs communs} à deux entiers naturels non nuls $a$ et $b$ sont les éléments de l'intersection $D(a) \cap D(b)$.
\end{definition}

Prenons l'exemple de notre accroche : $a = 84$ et $b = 126$.
\begin{align*}
D(84) &= \{1; 2; 3; 4; 6; 7; 12; 14; 21; 28; 42; 84\} \\
D(126) &= \{1; 2; 3; 6; 7; 9; 14; 18; 21; 42; 63; 126\}
\end{align*}
Les diviseurs communs sont ceux qui apparaissent dans \textbf{les deux} listes :
\[ D(84) \cap D(126) = \{1; 2; 3; 6; 7; 14; 21; 42\} \]

\begin{popfigure}
\begin{tikzpicture}[scale=0.9, every node/.style={font=\small}]
  \tikzset{
    venn circle/.style={draw=popblue, thick},
    venn label/.style={font=\bfseries\large, text=popdark},
    elem label/.style={font=\small, text=popdark},
    highlight/.style={font=\bfseries\large, text=poporange, fill=popgoldL, rounded corners=2pt, inner sep=2pt}
  }
  \fill[popblue!15] (-2,0) circle (2.5cm);
  \fill[poporange!15] (2,0) circle (2.5cm);
  \begin{scope}
    \clip (-2,0) circle (2.5cm);
    \fill[popgold!35] (2,0) circle (2.5cm);
  \end{scope}
  \draw[venn circle] (-2,0) circle (2.5cm);
  \draw[venn circle] (2,0) circle (2.5cm);
  \node[venn label] at (-3.5, 2.2) {$D(84)$};
  \node[venn label] at (3.5, 2.2) {$D(126)$};
  \node[elem label] at (-3.8, 1.2) {4};
  \node[elem label] at (-3.5, 0.4) {12};
  \node[elem label] at (-3.8, -0.4) {28};
  \node[elem label] at (-3.5, -1.2) {84};
  \node[elem label] at (3.8, 1.2) {9};
  \node[elem label] at (3.5, 0.4) {18};
  \node[elem label] at (3.8, -0.4) {63};
  \node[elem label] at (3.5, -1.2) {126};
  \node[elem label] at (0, 1.8) {1};
  \node[elem label] at (-0.6, 1.2) {2};
  \node[elem label] at (0.6, 1.2) {3};
  \node[elem label] at (-0.6, 0.4) {6};
  \node[elem label] at (0.6, 0.4) {7};
  \node[elem label] at (-0.6, -0.4) {14};
  \node[elem label] at (0.6, -0.4) {21};
  \node[highlight] at (0, -1.3) {42};
  \node[elem label, align=center] at (0, -2.7) {Le plus grand diviseur commun : \textbf{42}};
\end{tikzpicture}
\legende{Diagramme de Venn des diviseurs de 84 et 126. L'intersection (zone centrale) contient les diviseurs communs ; le plus grand (42) est mis en évidence.}
\end{popfigure}

\courssub{Définition du PGCD}

Parmi les diviseurs communs, il y en a un qui nous intéresse particulièrement : le plus grand. C'est lui qui permet de faire le \textbf{maximum} de groupes identiques (ou de sachets, de paquets, de lots...).

\begin{definition}[PGCD — Plus Grand Commun Diviseur]
Soient $a$ et $b$ deux entiers naturels non nuls. Le \cle{PGCD} de $a$ et $b$, noté $\operatorname{PGCD}(a;b)$, est le plus grand élément de l'ensemble $D(a) \cap D(b)$, c'est-à-dire le plus grand des diviseurs communs à $a$ et $b$.
\end{definition}

\begin{aretenir}[Notation et vocabulaire]
\begin{itemize}
  \item On lit $\operatorname{PGCD}(a;b)$ : « P-G-C-D de $a$ et $b$ ».
  \item Au BEPC, la notation $\operatorname{PGCD}(a;b)$ est la notation standard attendue dans la rédaction.
  \item Le PGCD de deux entiers non nuls est toujours un entier supérieur ou égal à 1, car 1 divise tout entier.
\end{itemize}
\end{aretenir}

Appliquons la définition à notre exemple :
\[ \operatorname{PGCD}(84;126) = \max\{1; 2; 3; 6; 7; 14; 21; 42\} = 42 \]

\courssub{Propriétés immédiates du PGCD}

Quelques propriétés découlent directement de la définition. Elles servent souvent de raccourcis ou de vérifications rapides.

\begin{propriete}[Propriétés élémentaires]
Soient $a$ et $b$ deux entiers naturels non nuls.
\begin{enumerate}
  \item \textbf{Symétrie :} $\operatorname{PGCD}(a;b) = \operatorname{PGCD}(b;a)$.
  \item \textbf{PGCD avec 1 :} $\operatorname{PGCD}(a;1) = 1$.
  \item \textbf{Divisibilité :} si $b$ divise $a$ (c'est-à-dire si $a$ est un multiple de $b$), alors $\operatorname{PGCD}(a;b) = b$.
  \item \textbf{Encadrement :} $\operatorname{PGCD}(a;b) \le a$ et $\operatorname{PGCD}(a;b) \le b$.
\end{enumerate}
\end{propriete}

\begin{experience}[Démonstration guidée : si $b$ divise $a$, alors $\operatorname{PGCD}(a;b)=b$]
\textbf{Objectif :} prouver la propriété 3 en quelques lignes de raisonnement.
\begin{enumerate}
  \item Comme $b$ divise $a$, tout diviseur de $b$ est aussi un diviseur de $a$. Donc $D(b) \subset D(a)$.
  \item L'intersection $D(a) \cap D(b)$ est donc exactement égale à $D(b)$.
  \item Or le plus grand élément de $D(b)$ est $b$ lui-même (tout entier non nul est son propre diviseur).
  \item \textbf{Conclusion :} $\operatorname{PGCD}(a;b) = b$. La propriété est démontrée.
\end{enumerate}
\end{experience}

\begin{exemplebox}[Applications directes des propriétés]
\begin{itemize}
  \item $\operatorname{PGCD}(45;15) = 15$ car $15$ divise $45$ (en effet $45 = 15 \times 3$).
  \item $\operatorname{PGCD}(100;1) = 1$.
  \item $\operatorname{PGCD}(27;81) = \operatorname{PGCD}(81;27) = 27$ (symétrie, puis divisibilité car $81 = 27 \times 3$).
\end{itemize}
\end{exemplebox}

\courssub{Méthode des listes de diviseurs}

C'est la méthode la plus intuitive, celle qui applique exactement la définition. Elle se déroule en trois étapes.

\begin{methode}[Trouver le PGCD par les listes de diviseurs]
Pour calculer $\operatorname{PGCD}(a;b)$ :
\begin{enumerate}
  \item \textbf{Lister} tous les diviseurs de $a$ (ensemble $D(a)$) et tous les diviseurs de $b$ (ensemble $D(b)$). \emph{Astuce : procéder par paires de facteurs $1 \times a$, $2 \times \frac{a}{2}$, etc., pour n'en oublier aucun.}
  \item \textbf{Repérer} les nombres présents dans les deux listes : c'est l'intersection $D(a) \cap D(b)$.
  \item \textbf{Sélectionner} le plus grand de ces diviseurs communs : c'est le PGCD.
\end{enumerate}
\end{methode}

\begin{exoresolu}[Calcul de $\operatorname{PGCD}(84;126)$ par les listes]
\textbf{Étape 1 : listes des diviseurs.} On cherche les paires de facteurs.
\begin{center}
\begin{tabularx}{\linewidth}{|c|X|X|}\hline
\rowcolor{popblueL}
\textbf{Nombre} & \textbf{Paires de facteurs} & \textbf{Liste ordonnée $D(n)$} \\ \hline
84 & $1\times84$, $2\times42$, $3\times28$, $4\times21$, $6\times14$, $7\times12$ & $\{1; 2; 3; 4; 6; 7; 12; 14; 21; 28; 42; 84\}$ \\ \hline
126 & $1\times126$, $2\times63$, $3\times42$, $6\times21$, $7\times18$, $9\times14$ & $\{1; 2; 3; 6; 7; 9; 14; 18; 21; 42; 63; 126\}$ \\ \hline
\end{tabularx}
\end{center}

\textbf{Étape 2 : diviseurs communs (intersection).} On compare les deux listes :
\[ D(84) \cap D(126) = \{1; 2; 3; 6; 7; 14; 21; 42\} \]

\textbf{Étape 3 : le plus grand.} Le maximum de cet ensemble est \textbf{42}.
\[ \operatorname{PGCD}(84;126) = 42 \]
\tcblower
\textbf{Interprétation concrète (problème des sachets) :}
On peut faire \textbf{42 sachets identiques} au maximum.
Chaque sachet contiendra : $84 \div 42 = 2$ beignets et $126 \div 42 = 3$ cacahuètes.
Il n'y a aucun reste, et il est impossible de faire plus de 42 sachets identiques.
\end{exoresolu}

\courssub{Limites de la méthode des listes}

La méthode des listes est parfaite pour de petits nombres (inférieurs à 200 environ). Mais essayons avec les nombres d'un exercice type BEPC : \textbf{1 785} et \textbf{1 275}.

\begin{attention}[Piège : la méthode des listes devient ingérable]
Pour $1\,785$ et $1\,275$ :
\begin{itemize}
  \item Il faut trouver \textbf{tous} les diviseurs de chaque nombre. Par exemple $1\,785 = 3 \times 5 \times 7 \times 17$, ce qui donne déjà 16 diviseurs pour 1 785 !
  \item Le risque d'en oublier (surtout les grands comme 255, 357 ou 595) est énorme.
  \item La comparaison des deux listes est longue et source d'erreurs.
\end{itemize}
\textbf{Conclusion :} pour de grands nombres, nous avons besoin de méthodes plus rapides : l'\textbf{algorithme des soustractions successives} et l'\textbf{algorithme d'Euclide}, que nous étudierons dans les sections 3 et 4.
\end{attention}

\courssub{Erreur fréquente : oublier le diviseur 1}

\begin{attention}[Erreur classique — le diviseur « invisible »]
L'erreur n°1 en 3ème : \textbf{oublier 1 dans la liste des diviseurs}.
\begin{itemize}
  \item \textbf{Rappel :} 1 divise \textbf{tout} entier naturel. Donc $1 \in D(n)$ pour tout entier naturel non nul $n$.
  \item Conséquence : 1 est \textbf{toujours} un diviseur commun. L'intersection $D(a) \cap D(b)$ n'est \textbf{jamais vide}.
  \item Si le seul diviseur commun est 1, alors $\operatorname{PGCD}(a;b) = 1$. On dit alors que $a$ et $b$ sont \textbf{premiers entre eux} (nous y reviendrons en section 5).
\end{itemize}
\textbf{Exemple :} $\operatorname{PGCD}(14;15)$. On a $D(14)=\{1;2;7;14\}$ et $D(15)=\{1;3;5;15\}$. Seul diviseur commun : $1$. Donc $\operatorname{PGCD}(14;15)=1$.
\end{attention}

\courssub{Exemple complet : un problème de carrelage}

\begin{exoresolu}[Dallage d'une cour de récréation]
La cour d'un collège à Yaoundé est un rectangle de \textbf{18 m} sur \textbf{24 m}. On veut la carreler avec des carreaux \textbf{carrés} identiques, dont le côté est un nombre entier de mètres, sans couper aucun carreau. Quel est le \textbf{plus grand} côté possible pour les carreaux ? Combien de carreaux faut-il alors ?

\textbf{Analyse :} le côté du carreau doit diviser la longueur (24 m) \textbf{et} la largeur (18 m) : c'est un diviseur commun à 18 et 24. Pour utiliser le moins de carreaux possible (donc les plus grands possible), il faut le \textbf{PGCD} de 18 et 24.

\textbf{Résolution :}
\begin{align*}
D(18) &= \{1; 2; 3; 6; 9; 18\} \\
D(24) &= \{1; 2; 3; 4; 6; 8; 12; 24\} \\
D(18) \cap D(24) &= \{1; 2; 3; 6\} \\
\operatorname{PGCD}(18;24) &= 6
\end{align*}
\tcblower
\textbf{Réponse :}
\begin{itemize}
  \item Côté maximal du carreau : \textbf{6 m}.
  \item Sur la largeur : $18 \div 6 = 3$ carreaux ; sur la longueur : $24 \div 6 = 4$ carreaux.
  \item Nombre total de carreaux : $3 \times 4 = \textbf{12}$ carreaux.
\end{itemize}
\textbf{Astuce de rédaction :} toujours conclure par une phrase répondant exactement à la question posée : « Le plus grand côté possible est 6 m ; il faut alors 12 carreaux. »
\end{exoresolu}

\courssub{Exercices d'application immédiate}

\begin{exercice}[Entraînement — méthode des listes]
\begin{enumerate}
  \item Déterminer $\operatorname{PGCD}(36;48)$ par la méthode des listes.
  \item Déterminer $\operatorname{PGCD}(13;26)$ en utilisant la propriété de divisibilité.
  \item Un marchand a 56 mangues et 72 oranges. Il veut confectionner des paniers identiques (même nombre de mangues et même nombre d'oranges dans chaque panier), en utilisant tous les fruits. Quel est le nombre maximal de paniers ? Combien de fruits de chaque sorte contient alors un panier ?
  \item Vrai ou Faux ? Justifier la réponse.
  \begin{itemize}
    \item[a)] $\operatorname{PGCD}(a;b) \le a$ et $\operatorname{PGCD}(a;b) \le b$.
    \item[b)] Si $\operatorname{PGCD}(a;b) = a$, alors $a$ divise $b$.
  \end{itemize}
\end{enumerate}
\end{exercice}

\begin{corrige}[Exercice n°1]
\begin{enumerate}
  \item $D(36)=\{1;2;3;4;6;9;12;18;36\}$ et $D(48)=\{1;2;3;4;6;8;12;16;24;48\}$. Diviseurs communs : $\{1;2;3;4;6;12\}$. Donc $\operatorname{PGCD}(36;48)=12$.
  \item $13$ divise $26$ (car $26=13\times2$), donc d'après la propriété de divisibilité : $\operatorname{PGCD}(13;26)=13$.
  \item On cherche $\operatorname{PGCD}(56;72)$. $D(56)=\{1;2;4;7;8;14;28;56\}$ et $D(72)=\{1;2;3;4;6;8;9;12;18;24;36;72\}$. Diviseurs communs : $\{1;2;4;8\}$, donc $\operatorname{PGCD}(56;72)=8$.
  \begin{itemize}
    \item Nombre maximal de paniers : \textbf{8}.
    \item Mangues par panier : $56 \div 8 = 7$.
    \item Oranges par panier : $72 \div 8 = 9$.
  \end{itemize}
  \item \textbf{a) Vrai.} Le PGCD est un diviseur de $a$ et un diviseur de $b$ ; or tout diviseur d'un entier naturel non nul est inférieur ou égal à cet entier. Donc $\operatorname{PGCD}(a;b) \le a$ et $\operatorname{PGCD}(a;b) \le b$.
  \textbf{b) Vrai.} Par définition, $\operatorname{PGCD}(a;b)$ est un diviseur commun à $a$ et $b$. Si $\operatorname{PGCD}(a;b)=a$, alors $a$ est un diviseur de $b$, c'est-à-dire $a$ divise $b$.
\end{enumerate}
\end{corrige}

\courssub{Synthèse de la section}

\begin{aretenir}[Ce qu'il faut retenir]
\begin{itemize}
  \item Les diviseurs communs à $a$ et $b$ forment l'intersection $D(a) \cap D(b)$.
  \item $\operatorname{PGCD}(a;b)$ est le \textbf{plus grand} élément de cette intersection.
  \item \textbf{Méthode des listes :} lister $D(a)$ et $D(b)$ $\rightarrow$ intersection $\rightarrow$ maximum.
  \item \textbf{Propriétés utiles :} symétrie ; $\operatorname{PGCD}(a;1)=1$ ; si $b$ divise $a$ alors $\operatorname{PGCD}(a;b)=b$.
  \item \textbf{Attention :} ne jamais oublier 1 dans les listes. Si le seul diviseur commun est 1, alors le PGCD vaut 1 (nombres premiers entre eux).
  \item \textbf{Limite :} méthode lente et risquée pour les grands nombres $\rightarrow$ on utilisera les algorithmes des sections suivantes.
\end{itemize}
\end{aretenir}

\begin{savaistu}[Le PGCD dans le quotidien camerounais]
Au marché Mokolo ou à la gare routière de Bafoussam, le PGCD s'invite sans qu'on le nomme :
\begin{itemize}
  \item \textbf{Transport :} un car de 54 places et un minibus de 18 places desservent la même ligne. Pour former des groupes identiques sans laisser de place vide, la taille maximale d'un groupe est $\operatorname{PGCD}(54;18)=18$.
  \item \textbf{Électrification rurale :} pour planter des poteaux à intervalles réguliers le long de deux portions de route de 12 m et 18 m, l'espacement maximal commun est $\operatorname{PGCD}(12;18)=6$ m.
  \item \textbf{Téléphonie mobile :} les cartes de recharge (500 F, 1 000 F, 2 000 F, 5 000 F, 10 000 F) ont pour PGCD 500 F : c'est l'unité de base commune à toutes ces valeurs.
\end{itemize}
Les mathématiques ne sont pas abstraites : elles organisent notre vie quotidienne !
\end{savaistu}

\textbf{Transition vers la section 3 :} maintenant que nous savons \emph{ce que} nous cherchons (le PGCD) et comment le trouver pour de petits nombres, nous allons découvrir deux méthodes historiques — l'\textbf{algorithme des soustractions successives} et l'\textbf{algorithme d'Euclide} — qui permettent de calculer $\operatorname{PGCD}(1\,785;1\,275)$ en quelques lignes seulement, sans lister le moindre diviseur.

\coursec{Algorithme des soustractions successives}

Dans la section précédente, nous avons défini le PGCD de deux entiers naturels et appris à le trouver en écrivant la \emph{liste complète} des diviseurs de chaque nombre. Cette méthode est fiable, mais elle devient vite fastidieuse : essaie de dresser la liste des diviseurs de $252$ et celle de $108$... tu comprendras qu'il faut parfois beaucoup de temps ! Heureusement, les mathématiciens ont mis au point des \cle{algorithmes}, c'est-à-dire des procédures mécaniques, étape par étape, qui conduisent au résultat sans jamais avoir besoin de lister tous les diviseurs. Le premier que nous étudions est l'\cle{algorithme des soustractions successives}. Il repose sur une idée géométrique très simple : le découpage de segments.

\courssub{L'idée intuitive : découper des segments}

Imagine deux planches de bois : l'une mesure $126$ cm, l'autre $84$ cm. Un menuisier de Nkongsamba veut les découper en morceaux \emph{tous de la même longueur}, la plus grande possible, sans perdre de bois. Cette longueur commune doit diviser à la fois $126$ et $84$ : c'est exactement le $\mathrm{PGCD}(126\,;84)$.

Voici l'astuce : si une longueur $d$ permet de découper exactement la planche de $126$ cm \emph{et} celle de $84$ cm, alors elle permet aussi de découper le morceau \emph{restant} quand on retire $84$ cm de la grande planche, c'est-à-dire $126-84=42$ cm. Chercher le PGCD de $126$ et $84$ revient donc à chercher le PGCD de $84$ et $42$ : un problème \emph{plus simple}, car les nombres sont plus petits ! En répétant l'opération, les nombres diminuent jusqu'à ce que la réponse devienne évidente.

\begin{popfigure}
\begin{tikzpicture}[scale=0.055, >=Stealth]
% Segment 126
\draw[very thick, popblue] (0,10) -- (126,10);
\draw[popblue] (0,9.2) -- (0,10.8);
\draw[popblue] (126,9.2) -- (126,10.8);
\node[popblue, above] at (63,10.5) {$126$};
% Segment 84
\draw[very thick, poporange] (0,6) -- (84,6);
\draw[poporange] (0,5.2) -- (0,6.8);
\draw[poporange] (84,5.2) -- (84,6.8);
\node[poporange, above] at (42,6.5) {$84$};
% Différence 42
\draw[very thick, popgreen] (84,6) -- (126,6);
\draw[popgreen] (126,5.2) -- (126,6.8);
\draw[<->, popgreen, thick] (84,3.6) -- (126,3.6);
\node[popgreen, below] at (105,3.4) {$126-84=42$};
% Étape suivante : 84 et 42
\draw[very thick, poporange] (0,-2) -- (84,-2);
\draw[poporange] (0,-2.8) -- (0,-1.2);
\draw[poporange] (84,-2.8) -- (84,-1.2);
\node[poporange, above] at (42,-1.5) {$84$};
\draw[very thick, popgreen] (0,-6) -- (42,-6);
\draw[popgreen] (0,-6.8) -- (0,-5.2);
\draw[popgreen] (42,-6.8) -- (42,-5.2);
\node[popgreen, below] at (21,-6.8) {$42$};
\draw[very thick, poppurple] (42,-6) -- (84,-6);
\draw[poppurple] (84,-6.8) -- (84,-5.2);
\draw[<->, poppurple, thick] (42,-8.4) -- (84,-8.4);
\node[poppurple, below] at (63,-8.6) {$84-42=42$};
% Résultat
\node[popdark, font=\bfseries] at (63,-12) {Il reste deux morceaux égaux à $42$ : $\mathrm{PGCD}(126\,;84)=42$};
\end{tikzpicture}
\legende{Interprétation géométrique de l'algorithme des soustractions : on retire la petite longueur de la grande, jusqu'à obtenir deux morceaux égaux.}
\end{popfigure}

\courssub{La propriété fondamentale}

Toute la méthode repose sur le résultat suivant, que tu dois connaître et savoir justifier.

\begin{propriete}[Propriété fondamentale des soustractions]
Soient $a$ et $b$ deux entiers naturels non nuls tels que $a > b$. Alors :
\[ \mathrm{PGCD}(a\,;b) = \mathrm{PGCD}(a-b\,;b). \]
Autrement dit : \textbf{remplacer le plus grand des deux nombres par la différence ne change pas le PGCD}.
\end{propriete}

Cette propriété n'est pas magique : elle se démontre en quelques lignes, et la démonstration est formatrice car elle utilise uniquement la définition d'un diviseur. Suivons-la pas à pas.

\begin{experience}[Démonstration guidée de la propriété fondamentale]
Soient $a$ et $b$ deux entiers naturels non nuls avec $a > b$. Notons $\mathcal{D}$ l'ensemble des diviseurs communs à $a$ et $b$, et $\mathcal{D}'$ l'ensemble des diviseurs communs à $a-b$ et $b$. Nous allons montrer que $\mathcal{D} = \mathcal{D}'$.

\textbf{Étape 1 : tout diviseur commun à $a$ et $b$ divise aussi $a-b$.}

Soit $d$ un diviseur commun à $a$ et $b$. Par définition, il existe deux entiers naturels $k$ et $k'$ tels que :
\[ a = k \times d \qquad \text{et} \qquad b = k' \times d. \]
Alors :
\[ a - b = k \times d - k' \times d = (k - k') \times d. \]
Comme $a > b$, on a $k > k'$, donc $k - k'$ est un entier naturel. Ainsi $a-b$ est un multiple de $d$, c'est-à-dire que $d$ divise $a-b$. Puisque $d$ divise aussi $b$, le nombre $d$ est un diviseur commun à $a-b$ et $b$. Donc $\mathcal{D} \subset \mathcal{D}'$.

\textbf{Étape 2 : tout diviseur commun à $a-b$ et $b$ divise aussi $a$.}

Soit $d'$ un diviseur commun à $a-b$ et $b$. Il existe deux entiers naturels $m$ et $m'$ tels que :
\[ a - b = m \times d' \qquad \text{et} \qquad b = m' \times d'. \]
Or $a = (a-b) + b$, donc :
\[ a = m \times d' + m' \times d' = (m + m') \times d'. \]
Ainsi $d'$ divise $a$. Puisque $d'$ divise aussi $b$, le nombre $d'$ est un diviseur commun à $a$ et $b$. Donc $\mathcal{D}' \subset \mathcal{D}$.

\textbf{Étape 3 : conclusion.}

Les deux inclusions montrent que $\mathcal{D} = \mathcal{D}'$ : les couples $(a\,;b)$ et $(a-b\,;b)$ ont \emph{exactement les mêmes diviseurs communs}. En particulier, le plus grand de ces diviseurs communs est le même :
\[ \mathrm{PGCD}(a\,;b) = \mathrm{PGCD}(a-b\,;b). \qquad \blacksquare \]
\end{experience}

\begin{aretenir}[L'idée-clé]
Soustraire le plus petit nombre du plus grand \textbf{conserve l'ensemble des diviseurs communs}, donc conserve le PGCD. Chaque soustraction rend les nombres plus petits sans changer la réponse : on « descend un escalier » dont la dernière marche est le PGCD.
\end{aretenir}

\courssub{Description de l'algorithme}

Répéter la propriété fondamentale donne un véritable algorithme, c'est-à-dire une recette que l'on peut appliquer mécaniquement — ou faire exécuter par une machine.

\begin{methode}[Algorithme des soustractions successives]
Pour déterminer $\mathrm{PGCD}(a\,;b)$ avec $a$ et $b$ entiers naturels non nuls :
\begin{enumerate}
\item \textbf{Comparer} les deux nombres : s'ils sont égaux, c'est terminé, le PGCD est cette valeur commune.
\item \textbf{Sinon}, calculer la différence entre le plus grand et le plus petit.
\item \textbf{Remplacer} le plus grand nombre par cette différence (le plus petit ne change pas).
\item \textbf{Recommencer} à l'étape 1 avec le nouveau couple, jusqu'à obtenir deux nombres égaux.
\item \textbf{Conclure} : le PGCD cherché est ce nombre commun final.
\end{enumerate}
\end{methode}

\begin{popfigure}
\begin{tikzpicture}[node distance=0.9cm, >=Stealth,
  startstop/.style={rectangle, rounded corners=8pt, draw=popdark, fill=popgoldL, thick, minimum width=4.6cm, minimum height=0.9cm, align=center, font=\small},
  process/.style={rectangle, draw=popblue, fill=popblueL, thick, minimum width=5.2cm, minimum height=0.9cm, align=center, font=\small},
  decision/.style={diamond, draw=poppink, fill=poppinkL, thick, aspect=2.4, align=center, font=\small, inner sep=1pt}]
\node[startstop] (start) {Début : entiers $a$ et $b$ non nuls};
\node[decision, below=of start] (test) {$a = b$ ?};
\node[process, below=of test, align=center] (soustraction){Remplacer le plus grand\\ par la différence (grand $-$ petit)};
\node[startstop, right=3.2cm of test] (fin) {Le PGCD est la valeur commune $a=b$};
\draw[->, thick, popdark] (start) -- (test);
\draw[->, thick, popdark] (test) -- node[right, font=\small]{non} (soustraction);
\draw[->, thick, popdark] (test) -- node[above, font=\small]{oui} (fin);
\draw[->, thick, popdark] (soustraction.west) -- ++(-1.3,0) |- (test.west);
\end{tikzpicture}
\legende{Organigramme de l'algorithme des soustractions successives : tant que les deux nombres sont différents, on soustrait ; dès qu'ils sont égaux, on a trouvé le PGCD.}
\end{popfigure}

\begin{exemplebox}[Exemple déroulé pas à pas : $\mathrm{PGCD}(126\,;84)$]
Appliquons l'algorithme en rédigeant chaque étape sous forme d'égalité de PGCD :
\begin{align*}
\mathrm{PGCD}(126\,;84) &= \mathrm{PGCD}(126-84\,;84) = \mathrm{PGCD}(42\,;84)\\
&= \mathrm{PGCD}(84-42\,;42) = \mathrm{PGCD}(42\,;42).
\end{align*}
Les deux nombres sont égaux à $42$ : on s'arrête.
\[ \boxed{\mathrm{PGCD}(126\,;84) = 42} \]
Vérification : $126 = 3 \times 42$ et $84 = 2 \times 42$, et $3$ et $2$ n'ont aucun diviseur commun autre que $1$. Le résultat est cohérent.
\end{exemplebox}

\begin{exoresolu}[Calculer $\mathrm{PGCD}(252\,;108)$ par soustractions successives]
Un fleuriste de Douala dispose de $252$ roses rouges et $108$ roses blanches. Il veut composer des bouquets \emph{tous identiques}, en utilisant toutes les fleurs, avec le plus grand nombre possible de bouquets. Combien de bouquets réalisera-t-il ?

\tcblower

\textbf{Solution détaillée.}

Le nombre de bouquets doit diviser à la fois $252$ et $108$ : le nombre maximal de bouquets est $\mathrm{PGCD}(252\,;108)$. Appliquons l'algorithme des soustractions successives :
\begin{align*}
\mathrm{PGCD}(252\,;108) &= \mathrm{PGCD}(252-108\,;108) = \mathrm{PGCD}(144\,;108)\\
&= \mathrm{PGCD}(144-108\,;108) = \mathrm{PGCD}(36\,;108)\\
&= \mathrm{PGCD}(108-36\,;36) = \mathrm{PGCD}(72\,;36)\\
&= \mathrm{PGCD}(72-36\,;36) = \mathrm{PGCD}(36\,;36).
\end{align*}
Les deux nombres sont égaux : $\mathrm{PGCD}(252\,;108) = 36$.

\textbf{Conclusion :} le fleuriste réalisera $36$ bouquets identiques, chacun contenant $252 \div 36 = 7$ roses rouges et $108 \div 36 = 3$ roses blanches.
\end{exoresolu}

\courssub{Pièges et erreurs fréquentes}

\begin{attention}[Ne pas s'arrêter trop tôt !]
L'erreur la plus fréquente consiste à \textbf{s'arrêter dès qu'un nombre devient $1$ (ou dès qu'un nombre divise l'autre)} et à conclure hâtivement. Même si le PGCD est bien $1$ quand $1$ apparaît, \textbf{la règle de l'algorithme impose de continuer jusqu'à l'égalité des deux nombres} (c'est-à-dire jusqu'au couple $(1\,;1)$). S'arrêter avant l'égalité, c'est ne pas suivre l'algorithme jusqu'au bout.

Par exemple, pour $\mathrm{PGCD}(48\,;18)$ :
\[ \mathrm{PGCD}(48\,;18) = \mathrm{PGCD}(30\,;18) = \mathrm{PGCD}(12\,;18) = \mathrm{PGCD}(6\,;12) = \mathrm{PGCD}(6\,;6) = 6. \]
Si l'on s'était arrêté au couple $(6\,;12)$ en pensant « $6$ divise $12$, c'est fini » sans écrire l'étape finale, la réponse reste juste ici — mais cette habitude est dangereuse. La discipline est simple : \textbf{on poursuit les soustractions jusqu'à l'égalité des deux nombres}, et ce nombre commun est le PGCD.

Autre piège : ne pas oublier de \emph{remplacer le plus grand} nombre. Écrire $\mathrm{PGCD}(84\,;126) = \mathrm{PGCD}(84-126\,;\ldots)$ n'a pas de sens avec des entiers naturels : on calcule toujours « grand $-$ petit ».
\end{attention}

\courssub{Efficacité et limite de la méthode}

L'algorithme des soustractions successives est simple à comprendre et ne demande que des soustractions : aucune table de diviseurs, aucune division. C'est sa grande force. Mais il a une faiblesse : quand les deux nombres sont \emph{très inégaux}, il faut soustraire le petit nombre de très nombreuses fois.

\begin{exemplebox}[Un cas défavorable : $\mathrm{PGCD}(1\,000\,;7)$]
En appliquant l'algorithme, on obtient successivement :
\[ (1\,000\,;7) \to (993\,;7) \to (986\,;7) \to \cdots \to (9\,;7) \to (2\,;7) \to (2\,;5) \to (2\,;3) \to (2\,;1) \to (1\,;1). \]
Il faut plus de $140$ soustractions pour conclure que $\mathrm{PGCD}(1\,000\,;7) = 1$ ! Chaque soustraction de $7$ est en réalité « contenue » dans une seule division : $1\,000 = 142 \times 7 + 6$. Cette remarque est le point de départ d'une méthode beaucoup plus rapide, l'\emph{algorithme d'Euclide}, que nous étudierons dans la prochaine section et qui remplace les soustractions répétées par des divisions euclidiennes.
\end{exemplebox}

\begin{savaistu}[L'un des plus vieux algorithmes de l'histoire]
Cette méthode des soustractions successives est décrite dans les \emph{Éléments} d'\textbf{Euclide}, vers $300$ avant J.-C., il y a donc plus de $2\,300$ ans ! Euclide la présentait géométriquement, exactement comme notre découpage de segments : il s'agissait de trouver la plus grande « mesure commune » de deux longueurs. Les historiens des mathématiques la considèrent comme l'un des tout premiers \emph{algorithmes} écrits de l'humanité. Aujourd'hui encore, des variantes de cet algorithme tournent dans les ordinateurs et les téléphones : quand ton smartphone simplifie une fraction ou sécurise une communication, il utilise des descendants directs de cette idée vieille de vingt-trois siècles.
\end{savaistu}

\begin{exercice}[À toi de jouer]
\begin{enumerate}
\item En utilisant l'algorithme des soustractions successives, détermine $\mathrm{PGCD}(135\,;90)$ puis $\mathrm{PGCD}(96\,;60)$.
\item Un commerçant de Yaoundé a $135$ savons et $90$ tubes de dentifrice. Il veut préparer des colis identiques en utilisant tout le stock. Quel est le nombre maximal de colis ? Que contient chaque colis ?
\item Sans faire tous les calculs, explique pourquoi l'algorithme des soustractions serait très long pour calculer $\mathrm{PGCD}(2\,024\,;3)$.
\end{enumerate}
\end{exercice}

\begin{corrige}[Exercice 1]
\textbf{1.} Pour $\mathrm{PGCD}(135\,;90)$ :
\[ \mathrm{PGCD}(135\,;90) = \mathrm{PGCD}(45\,;90) = \mathrm{PGCD}(45\,;45) = 45. \]
Pour $\mathrm{PGCD}(96\,;60)$ :
\[ \mathrm{PGCD}(96\,;60) = \mathrm{PGCD}(36\,;60) = \mathrm{PGCD}(36\,;24) = \mathrm{PGCD}(12\,;24) = \mathrm{PGCD}(12\,;12) = 12. \]

\textbf{2.} Le nombre maximal de colis est $\mathrm{PGCD}(135\,;90) = 45$. Chaque colis contient $135 \div 45 = 3$ savons et $90 \div 45 = 2$ tubes de dentifrice.

\textbf{3.} Le nombre $3$ est très petit devant $2\,024$ : il faudrait soustraire $3$ plus de $670$ fois avant de ramener le grand nombre en dessous de $3$. L'algorithme des soustractions est donc très long ici ; l'algorithme d'Euclide, basé sur la division euclidienne, donnera la réponse en quelques étapes seulement.
\end{corrige}

\begin{aretenir}[Bilan de la section]
\begin{itemize}
\item \textbf{Propriété fondamentale} : si $a > b$, alors $\mathrm{PGCD}(a\,;b) = \mathrm{PGCD}(a-b\,;b)$, car les deux couples ont exactement les mêmes diviseurs communs.
\item \textbf{Algorithme} : remplacer le plus grand nombre par la différence, répéter, et s'arrêter \emph{uniquement} quand les deux nombres sont égaux : cette valeur commune est le PGCD.
\item \textbf{Limite} : la méthode peut exiger un très grand nombre d'étapes quand les nombres sont très inégaux — d'où l'algorithme d'Euclide, objet de la section suivante.
\end{itemize}
\end{aretenir}

\coursec{Algorithme d'Euclide}

\courssub{Transition depuis les soustractions successives}

Dans la section précédente, nous avons vu que l'algorithme des soustractions successives permet de trouver le PGCD de deux nombres en remplaçant répétitivement le plus grand par la différence entre le plus grand et le plus petit. Cette méthode est juste, mais elle peut devenir très longue lorsque les nombres sont grands et proches l'un de l'autre (par exemple, pour \cle{PGCD(1000;999)}, il faudrait faire 999 soustractions !).

L'algorithme d'Euclide, que nous étudions ici, est une version \emph{accélérée} de cette idée. Au lieu de soustraire $b$ une seule fois à $a$, on soustrait $b$ le maximum de fois possible en une seule étape : c'est exactement ce que fait la \cle{division euclidienne}. Rappelons que pour deux entiers naturels $a$ et $b$ ($b \neq 0$), il existe un unique couple $(q;r)$ tel que :
\[ a = b \times q + r \quad \text{avec} \quad 0 \le r < b \]
Ici, $q$ est le quotient et $r$ le reste. L'algorithme d'Euclide repose sur une propriété fondamentale qui lie le PGCD du couple de départ $(a;b)$ au couple suivant $(b;r)$.

\begin{propriete}[Propriété fondamentale : invariance du PGCD par division euclidienne]
Soient $a$ et $b$ deux entiers naturels non nuls, avec $a \ge b$. Effectuons la division euclidienne de $a$ par $b$ :
\[ a = bq + r \qquad (0 \le r < b) \]
Alors :
\[ \cle{PGCD(a;b) = PGCD(b;r)} \]
\end{propriete}

\begin{experience}[Démonstration guidée de la propriété]
Comprenons \emph{pourquoi} cette égalité est vraie. C'est une conséquence directe de la propriété des soustractions successives que nous venons d'étudier.

\begin{enumerate}
    \item \textbf{Écriture de la division :} $a = bq + r$. On peut réécrire $r = a - bq$.
    \item \textbf{Un diviseur commun à $a$ et $b$ divise $r$ :} Soit $d$ un diviseur commun à $a$ et $b$. Alors $d$ divise $a$ et $d$ divise $b$. Donc $d$ divise $bq$ (multiple de $b$). Par conséquent, $d$ divise la différence $a - bq$, c'est-à-dire $d$ divise $r$. Ainsi, $d$ est un diviseur commun à $b$ et $r$.
    \item \textbf{Un diviseur commun à $b$ et $r$ divise $a$ :} Réciproquement, soit $d'$ un diviseur commun à $b$ et $r$. Alors $d'$ divise $b$ et $d'$ divise $r$. Donc $d'$ divise $bq$ et $d'$ divise la somme $bq + r = a$. Ainsi, $d'$ est un diviseur commun à $a$ et $b$.
    \item \textbf{Conclusion :} Les couples $(a;b)$ et $(b;r)$ ont \emph{exactement les mêmes diviseurs communs}. Leur plus grand diviseur commun est donc identique : $PGCD(a;b) = PGCD(b;r)$.
\end{enumerate}
\end{experience}

\courssub{Description de l'algorithme}

Grâce à cette propriété, on peut remplacer le calcul de $PGCD(a;b)$ par celui de $PGCD(b;r)$ où $r < b$. Les nombres deviennent strictement plus petits. On réitère le processus : on divise $b$ par $r$, on obtient un nouveau reste $r_1$, et ainsi de suite. Comme les restes forment une suite strictement décroissante d'entiers naturels, on finit \emph{obligatoirement} par obtenir un reste nul.

\begin{propriete}[Algorithme d'Euclide — Théorème]
Soient $a$ et $b$ deux entiers naturels non nuls ($a \ge b$). On construit la suite des divisions euclidiennes suivantes :
\begin{align*}
a &= b q_1 + r_1 \quad (0 \le r_1 < b) \\
b &= r_1 q_2 + r_2 \quad (0 \le r_2 < r_1) \\
r_1 &= r_2 q_3 + r_3 \quad (0 \le r_3 < r_2) \\
&\vdots \\
r_{k-2} &= r_{k-1} q_k + r_k \quad (0 \le r_k < r_{k-1}) \\
r_{k-1} &= r_k q_{k+1} + 0
\end{align*}
Le dernier reste \textbf{non nul} $r_k$ est le PGCD de $a$ et $b$ :
\[ \cle{PGCD(a;b) = r_k} \]
\end{propriete}

\begin{aretenir}[Pourquoi le dernier reste non nul ?]
À l'étape finale, on a $r_{k-1} = r_k \times q_{k+1} + 0$. Cela signifie que $r_k$ divise $r_{k-1}$. En remontant les égalités grâce à la propriété $PGCD(x;y)=PGCD(y;r)$, on montre que $r_k$ divise tous les restes précédents, donc $b$ et $a$. C'est un diviseur commun. Comme tout diviseur commun divise les restes (démonstration guidée ci-dessus), il divise $r_k$. Donc $r_k$ est le \emph{plus grand}.
\end{aretenir}

\courssub{Exemple détaillé : PGCD(\num{1785} ; \num{1275})}

Appliquons l'algorithme pour calculer $PGCD(1785;1275)$. Nous allons présenter les calculs dans un tableau clair, format standard au BEPC.

\begin{exemplebox}[Calcul de PGCD(\num{1785} ; \num{1275}) par l'algorithme d'Euclide]
\textbf{Étape 1 :} On divise le plus grand (\num{1785}) par le plus petit (\num{1275}).
\[ 1785 = 1275 \times \mathbf{1} + \mathbf{510} \quad (r_1 = 510) \]
\textbf{Étape 2 :} On divise l'ancien diviseur (\num{1275}) par le reste (\num{510}).
\[ 1275 = 510 \times \mathbf{2} + \mathbf{255} \quad (r_2 = 255) \]
\textbf{Étape 3 :} On divise l'ancien diviseur (\num{510}) par le nouveau reste (\num{255}).
\[ 510 = 255 \times \mathbf{2} + \mathbf{0} \quad (r_3 = 0) \]
\textbf{Arrêt :} Le reste est nul. Le dernier reste \textbf{non nul} est \textbf{255}.
\[ \boxed{PGCD(1785;1275) = 255} \]
\end{exemplebox}

\begin{popfigure}
\begin{tikzpicture}[
    table/.style={matrix of nodes, nodes={draw, minimum height=1.2cm, minimum width=2.8cm, anchor=center, font=\small}, column sep=-\pgflinewidth, row sep=-\pgflinewidth, nodes in empty cells},
    header/.style={fill=popblue, text=white, font=\bfseries\small},
    highlight/.style={fill=popgoldL, font=\bfseries\small},
    lastnonzero/.style={fill=popgreenL, font=\bfseries\small, draw=popgreen, very thick}
]
\matrix (T) [table] {
|[header]| Dividende & |[header]| Diviseur & |[header]| Quotient & |[header]| Reste \\
|[highlight]| 1785 & 1275 & 1 & 510 \\
1275 & |[highlight]| 510 & 2 & |[lastnonzero]| 255 \\
510 & 255 & 2 & 0 \\
};
% Flèches pour montrer le décalage
\draw[poporange, thick, -{Stealth[length=3mm]}] (T-1-2.south) -- ++(0,-0.2) -| (T-2-1.north);
\draw[poporange, thick, -{Stealth[length=3mm]}] (T-1-4.south) -- ++(0,-0.2) -| (T-2-2.north);
\draw[poporange, thick, -{Stealth[length=3mm]}] (T-2-2.south) -- ++(0,-0.2) -| (T-3-1.north);
\draw[poporange, thick, -{Stealth[length=3mm]}] (T-2-4.south) -- ++(0,-0.2) -| (T-3-2.north);
% Encadrement du PGCD final
\node[below=0.3cm of T-2-4, popgreen, font=\bfseries\small] {$\leftarrow$ Dernier reste non nul = PGCD};
\end{tikzpicture}
\legende{Tableau des divisions successives pour PGCD(1785;1275). Les flèches oranges montrent la règle : le diviseur devient dividende, le reste devient diviseur.}
\end{popfigure}

\courssub{Organigramme de l'algorithme (Vision algorithmique)}

Pour bien ancrer la logique itérative, voici l'organigramme (diagramme de flux) de l'algorithme d'Euclide. C'est la représentation formelle que l'on retrouve en algorithmique (programmation).

\begin{popfigure}
\begin{tikzpicture}[
    node distance=1.5cm and 2cm,
    startstop/.style={rectangle, rounded corners, minimum width=3cm, minimum height=1cm, text centered, draw=popblue, fill=popblueL, font=\small},
    process/.style={rectangle, minimum width=3.5cm, minimum height=1cm, text centered, draw=popdark, fill=popcream, font=\small},
    decision/.style={diamond, aspect=2, minimum width=3cm, minimum height=1.5cm, text centered, draw=poporange, fill=poporangeL, font=\small},
    arrow/.style={thick, -{Stealth[length=3mm]}, popdark},
    io/.style={trapezium, trapezium left angle=70, trapezium right angle=110, minimum width=3cm, minimum height=1cm, text centered, draw=poppurple, fill=poppurpleL, font=\small}
]
% Nœuds
\node (start) [startstop] {DÉBUT};
\node (input) [io, below of=start] {Entrer $a$, $b$ ($a \ge b > 0$)};
\node (div) [process, below of=input] {Calculer $r$ = reste de $a \div b$};
\node (test) [decision, below of=div] {$r = 0$ ?};
\node (assign) [process, below of=test, yshift=-0.5cm] {$a \leftarrow b$ ; $b \leftarrow r$};
\node (output) [io, below of=assign, yshift=-0.5cm] {Afficher $b$ (PGCD)};
\node (stop) [startstop, below of=output] {FIN};
% Flèches
\draw [arrow] (start) -- (input);
\draw [arrow] (input) -- (div);
\draw [arrow] (div) -- (test);
\draw [arrow] (test) -- node[anchor=east, font=\small, xshift=-2mm] {Non} (assign);
\draw [arrow] (assign) |- (div);
\draw [arrow] (test) -- node[anchor=west, font=\small, xshift=2mm] {Oui} (output);
\draw [arrow] (output) -- (stop);
\end{tikzpicture}
\legende{Organigramme de l'algorithme d'Euclide. La boucle \og Tant que $r \neq 0$ \fg\ correspond au chemin \og Non \fg.}
\end{popfigure}

\courssub{Méthode : Organiser les divisions successives en tableau}

Au BEPC, la présentation en tableau est souvent exigée pour la clarté et la rapidité. Voici la méthode pas à pas.

\begin{methode}[Calculer un PGCD par l'algorithme d'Euclide (présentation tabulaire)]
\begin{enumerate}
    \item \textbf{Initialisation :} Écrire les deux nombres $a$ et $b$ ($a \ge b$) dans la première ligne : \textbf{Dividende} = $a$, \textbf{Diviseur} = $b$.
    \item \textbf{Division :} Effectuer la division euclidienne du dividende par le diviseur. Noter le \textbf{Quotient} et le \textbf{Reste} sur la même ligne.
    \item \textbf{Test :} Si le \textbf{Reste} = 0, \textbf{STOP}. Le PGCD est le \textbf{Diviseur} de cette ligne (qui est le dernier reste non nul).
    \item \textbf{Itération :} Si le Reste $\neq 0$, passer à la ligne suivante :
          \begin{itemize}
              \item Nouveau \textbf{Dividende} = ancien \textbf{Diviseur}.
              \item Nouveau \textbf{Diviseur} = ancien \textbf{Reste}.
          \end{itemize}
    \item Revenir à l'étape 2.
\end{enumerate}
\end{methode}

\begin{attention}[Piège classique : le reste nul]
\textbf{Erreur très fréquente :} Lire le PGCD sur la \emph{dernière ligne} du tableau (celle où le reste est 0) et répondre 0.
\begin{itemize}
    \item \textbf{Règle d'or :} Le PGCD est le \textbf{dernier reste NON NUL}.
    \item Dans le tableau, c'est le nombre écrit dans la colonne \textbf{Diviseur} (ou \textbf{Reste} de la ligne d'avant) de la ligne \emph{précédant} celle où le reste devient 0.
    \item \emph{Astuce :} Entourez le dernier reste non nul dès qu'il apparaît pour ne pas l'oublier.
\end{itemize}
\end{attention}

\courssub{Comparaison des trois méthodes}

Nous disposons maintenant de trois méthodes pour calculer un PGCD. Il est crucial de savoir choisir la plus adaptée.

\begin{center}
\begin{tabularx}{\linewidth}{|l|X|X|X|}\hline
\rowcolor{popblueL}
\textbf{Méthode} & \textbf{Principe} & \textbf{Avantages} & \textbf{Inconvénients} \\ \hline
\textbf{Listes de diviseurs} & Lister tous les diviseurs de chaque nombre et prendre le plus grand commun. & Simple pour de \textbf{très petits nombres} ($< 100$). Visuel. & \textbf{Impossible} pour grands nombres (ex: 1785). Risque d'oubli de diviseurs. \\ \hline
\textbf{Soustractions successives} & Remplacer le plus grand par la différence jusqu'à égalité. & Pas de division à faire. Montre bien le lien avec la définition. & \textbf{Très lent} si nombres grands et proches (ex: 1000 et 999 $\to$ 999 étapes). \\ \hline
\textbf{\textbf{Algorithme d'Euclide}} & Divisions euclidiennes successives jusqu'à reste nul. & \textbf{La plus rapide} pour tous les nombres (même très grands). Standard au BEPC et au-delà. & Nécessite de savoir poser une division euclidienne. \\ \hline
\end{tabularx}
\end{center}

\begin{aretenir}[Choix stratégique]
\begin{itemize}
    \item Pour des nombres $\le 100$ : Listes ou Euclide (au choix).
    \item Pour des nombres $> 100$ : \textbf{Uniquement l'algorithme d'Euclide}.
    \item En épreuve (BEPC) : L'algorithme d'Euclide est \textbf{la méthode de référence}. Maîtrisez sa présentation en tableau.
\end{itemize}
\end{aretenir}

\courssub{Exercice résolu : Application concrète}

\begin{exoresolu}[Emballage de mangues — PGCD(\num{1785};\num{1275}) revisité]
\textbf{Énoncé :} Un producteur de mangues à Njombé a récolté \num{1785} mangues de variété \og Kent \fg{} et \num{1275} mangues de variété \og Brooks \fg{}. Il veut les emballer dans des caisses identiques, chaque caisse ne contenant qu'une seule variété, de manière à utiliser le \textbf{moins grand nombre possible de caisses} (donc le \textbf{plus grand nombre de mangues par caisse}). Combien de mangues mettra-t-il par caisse ? Combien de caisses au total ?
\tcblower
\textbf{Solution :}
\begin{enumerate}
    \item \textbf{Modélisation :} Soit $x$ le nombre de mangues par caisse. $x$ doit diviser 1785 (pour les Kent) et diviser 1275 (pour les Brooks). Pour minimiser le nombre de caisses, il faut \textbf{maximiser $x$}. Donc $x = PGCD(1785;1275)$.
    \item \textbf{Calcul du PGCD (Algorithme d'Euclide) :}
    \begin{center}
    \begin{tabular}{|c|c|c|c|}\hline
    \rowcolor{popblueL} \textbf{Dividende} & \textbf{Diviseur} & \textbf{Quotient} & \textbf{Reste} \\ \hline
    1785 & 1275 & 1 & 510 \\ \hline
    1275 & 510 & 2 & \cellcolor{popgreenL}\textbf{255} \\ \hline
    510 & 255 & 2 & 0 \\ \hline
    \end{tabular}
    \end{center}
    Dernier reste non nul = \textbf{255}. Donc $PGCD = 255$.
    \item \textbf{Réponse à la question 1 :} Il mettra \textbf{255 mangues par caisse}.
    \item \textbf{Nombre de caisses :}
    \begin{itemize}
        \item Variété Kent : $1785 \div 255 = 7$ caisses.
        \item Variété Brooks : $1275 \div 255 = 5$ caisses.
    \end{itemize}
    Total caisses = $7 + 5 = \textbf{12 caisses}$.
\end{enumerate}
\textbf{Conclusion :} Le producteur fera des caisses de 255 mangues, soit 12 caisses au total.
\end{exoresolu}

\begin{savaistu}[Euclide et l'informatique]
L'algorithme d'Euclide est l'un des \textbf{plus anciens algorithmes} encore en usage aujourd'hui (vers 300 av. J.-C., \emph{Éléments}, Livre VII). C'est aussi l'algorithme de base de la \textbf{cryptographie moderne} (RSA, chiffrement des cartes bancaires, HTTPS, MTN Mobile Money, Orange Money). Calculer un PGCD de nombres à 300 chiffres (clés RSA) prend quelques millisecondes à un ordinateur grâce à cet algorithme, alors que la méthode des listes prendrait plus de temps que l'âge de l'univers !
\end{savaistu}

\begin{exercice}[Entraînement : PGCD par Euclide]
Calculer les PGCD suivants en présentant les divisions dans un tableau (format BEPC).
\begin{enumerate}
    \item $PGCD(468; 198)$
    \item $PGCD(1071; 459)$
    \item $PGCD(2310; 1540)$
    \item Un carreleur doit poser des carreaux carrés identiques (le plus grands possible) pour recouvrir exactement un rectangle de dimensions $135$ cm par $90$ cm. Quelle est la longueur du côté du carreau ? Combien de carreaux faut-il ?
\end{enumerate}
\end{exercice}

\begin{corrige}[Exercice n°1]
\begin{enumerate}
    \item \textbf{PGCD(468; 198)} :
    \begin{tabular}{|c|c|c|c|}\hline \rowcolor{popblueL} Dividende & Diviseur & Quotient & Reste \\ \hline 468 & 198 & 2 & 72 \\ \hline 198 & 72 & 2 & \cellcolor{popgreenL}\textbf{54} \\ \hline 72 & 54 & 1 & 18 \\ \hline 54 & 18 & 3 & 0 \\ \hline \end{tabular} $\Rightarrow PGCD = \mathbf{18}$.
    \item \textbf{PGCD(1071; 459)} :
    \begin{tabular}{|c|c|c|c|}\hline \rowcolor{popblueL} Dividende & Diviseur & Quotient & Reste \\ \hline 1071 & 459 & 2 & 153 \\ \hline 459 & 153 & 3 & \cellcolor{popgreenL}\textbf{0} \\ \hline \end{tabular} $\Rightarrow PGCD = \mathbf{153}$ (Reste nul dès la 2ème division, le diviseur 153 est le PGCD).
    \item \textbf{PGCD(2310; 1540)} :
    \begin{tabular}{|c|c|c|c|}\hline \rowcolor{popblueL} Dividende & Diviseur & Quotient & Reste \\ \hline 2310 & 1540 & 1 & 770 \\ \hline 1540 & 770 & 2 & \cellcolor{popgreenL}\textbf{0} \\ \hline \end{tabular} $\Rightarrow PGCD = \mathbf{770}$.
    \item \textbf{Carreleur} : Côté du carreau = $PGCD(135; 90)$.
    \begin{tabular}{|c|c|c|c|}\hline \rowcolor{popblueL} Dividende & Diviseur & Quotient & Reste \\ \hline 135 & 90 & 1 & 45 \\ \hline 90 & 45 & 2 & \cellcolor{popgreenL}\textbf{0} \\ \hline \end{tabular} $\Rightarrow$ Côté = \textbf{45 cm}.
    Nombre de carreaux : $\frac{135 \times 90}{45 \times 45} = \frac{135}{45} \times \frac{90}{45} = 3 \times 2 = \textbf{6 carreaux}$.
\end{enumerate}
\end{corrige}

\coursec{Nombres premiers entre eux}

Dans la section précédente, tu as appris à calculer le PGCD de deux entiers naturels grâce à l'algorithme d'Euclide, une méthode rapide et sûre. Mais à quoi sert ce PGCD, concrètement ? L'une de ses applications les plus importantes concerne les \emph{fractions} : quand un commerçant de Douala te dit qu'il a vendu « les deux tiers » de sa marchandise, il utilise une fraction \emph{irréductible}, c'est-à-dire qu'on ne peut plus simplifier. Derrière cette idée se cache une notion fondamentale de l'arithmétique : celle des \cle{nombres premiers entre eux}. C'est elle que nous allons découvrir maintenant.

\courssub{Définition et premiers exemples}

\textbf{Une intuition pour commencer.} Prends deux entiers, par exemple 15 et 28. Cherchons leurs diviseurs communs. Les diviseurs de 15 sont : 1 ; 3 ; 5 ; 15. Les diviseurs de 28 sont : 1 ; 2 ; 4 ; 7 ; 14 ; 28. Le seul diviseur commun est 1. Autrement dit, 15 et 28 ne partagent \emph{rien} d'autre que 1 : on dit qu'ils sont « premiers entre eux ».

\begin{definition}[Nombres premiers entre eux]
Deux entiers naturels $a$ et $b$ (non nuls) sont dits \cle{premiers entre eux} lorsque leur PGCD est égal à 1 :
\[ \mathrm{PGCD}(a\,;b) = 1. \]
Autrement dit, le seul diviseur commun de $a$ et $b$ est 1.
\end{definition}

\begin{exemplebox}[Exemples et contre-exemples]
\begin{itemize}
\item $15$ et $28$ : $\mathrm{PGCD}(15\,;28) = 1$, donc 15 et 28 sont \textbf{premiers entre eux}.
\item $7$ et $12$ : $\mathrm{PGCD}(7\,;12) = 1$, donc 7 et 12 sont premiers entre eux.
\item $12$ et $18$ : $\mathrm{PGCD}(12\,;18) = 6 \neq 1$, donc 12 et 18 ne sont \textbf{pas} premiers entre eux.
\item $14$ et $21$ : $\mathrm{PGCD}(14\,;21) = 7 \neq 1$, donc 14 et 21 ne sont pas premiers entre eux.
\end{itemize}
\end{exemplebox}

\begin{attention}[Piège classique : « premiers entre eux » ne veut pas dire « premiers »]
L'erreur la plus fréquente au BEPC est de croire que deux nombres premiers entre eux doivent chacun être des \emph{nombres premiers}. C'est faux !

Reprenons 15 et 28 : $15 = 3 \times 5$ n'est pas premier, et $28 = 4 \times 7$ n'est pas premier non plus. Pourtant 15 et 28 \textbf{sont premiers entre eux}, car ils n'ont aucun diviseur commun autre que 1.

Être « premiers entre eux » est une propriété du \emph{couple} $(a\,;b)$, pas de chaque nombre séparément. Bien sûr, deux nombres premiers distincts (comme 7 et 13) sont toujours premiers entre eux, mais la réciproque est fausse.
\end{attention}

\courssub{Propriété fondamentale : réduire par le PGCD}

Voici maintenant le résultat central de cette section. L'idée intuitive est simple : quand on divise deux nombres par leur PGCD, on « retire » tout ce qu'ils avaient en commun ; il ne reste donc plus rien à partager, et les quotients obtenus sont premiers entre eux.

\begin{propriete}[Propriété fondamentale]
Soient $a$ et $b$ deux entiers naturels non nuls, et $d = \mathrm{PGCD}(a\,;b)$. Alors il existe deux entiers naturels $a'$ et $b'$ tels que :
\[ a = d \times a' \quad \text{et} \quad b = d \times b', \]
avec $a'$ et $b'$ \textbf{premiers entre eux} : $\mathrm{PGCD}(a'\,;b') = 1$.
\end{propriete}

Cette propriété se démontre par un raisonnement \emph{par l'absurde} : on suppose le contraire de ce qu'on veut prouver, et on montre que cela conduit à une contradiction. Suivons la démonstration pas à pas.

\begin{experience}[Démonstration guidée (raisonnement par l'absurde)]
On pose $d = \mathrm{PGCD}(a\,;b)$. Comme $d$ divise $a$ et $b$, on peut écrire $a = d \times a'$ et $b = d \times b'$.

\textbf{Étape 1 — L'hypothèse absurde.} Supposons que $a'$ et $b'$ ne soient \emph{pas} premiers entre eux. Alors ils admettent un diviseur commun $k$ avec $k > 1$.

\textbf{Étape 2 — Conséquence.} Puisque $k$ divise $a'$, on peut écrire $a' = k \times a''$. De même $b' = k \times b''$. Alors :
\[ a = d \times a' = d \times k \times a'' \quad \text{et} \quad b = d \times b' = d \times k \times b''. \]
Donc $d \times k$ est un \textbf{diviseur commun} de $a$ et de $b$.

\textbf{Étape 3 — La contradiction.} Or $k > 1$, donc $d \times k > d$. On aurait ainsi trouvé un diviseur commun de $a$ et $b$ \emph{strictement plus grand} que $d$... alors que $d$ est par définition le \emph{plus grand} diviseur commun de $a$ et $b$ ! C'est impossible.

\textbf{Conclusion.} L'hypothèse de l'étape 1 est fausse : $a'$ et $b'$ n'ont aucun diviseur commun autre que 1, donc $\mathrm{PGCD}(a'\,;b') = 1$. \hfill$\blacksquare$
\end{experience}

\begin{exemplebox}[Vérification sur un exemple]
Prenons $a = 84$ et $b = 126$. On a calculé (par l'algorithme d'Euclide) : $d = \mathrm{PGCD}(84\,;126) = 42$.

Alors $84 = 42 \times 2$ et $126 = 42 \times 3$, donc $a' = 2$ et $b' = 3$.

Et effectivement, $\mathrm{PGCD}(2\,;3) = 1$ : les quotients 2 et 3 sont bien premiers entre eux. La propriété est vérifiée.
\end{exemplebox}

\courssub{Application aux fractions irréductibles}

\begin{definition}[Fraction irréductible]
Une fraction $\dfrac{a}{b}$ est dite \cle{irréductible} lorsque son numérateur $a$ et son dénominateur $b$ sont premiers entre eux, c'est-à-dire lorsque $\mathrm{PGCD}(a\,;b) = 1$. On ne peut alors plus la simplifier.
\end{definition}

Grâce à la propriété précédente, simplifier une fraction « au maximum » devient un travail mécanique : il suffit de diviser le numérateur et le dénominateur par leur PGCD.

\begin{methode}[Rendre une fraction irréductible grâce au PGCD]
Pour écrire la fraction $\dfrac{a}{b}$ sous forme irréductible :
\begin{enumerate}
\item \textbf{Calculer} $d = \mathrm{PGCD}(a\,;b)$, de préférence par l'algorithme d'Euclide.
\item \textbf{Diviser} le numérateur et le dénominateur par $d$ : $\dfrac{a}{b} = \dfrac{a \div d}{b \div d} = \dfrac{a'}{b'}$.
\item \textbf{Conclure} : comme $a'$ et $b'$ sont premiers entre eux (propriété fondamentale), la fraction $\dfrac{a'}{b'}$ est irréductible.
\end{enumerate}
\end{methode}

\begin{attention}[Simplifier « petit à petit » ne suffit pas toujours]
On peut simplifier une fraction en plusieurs étapes (par 2, puis par 3, etc.), mais il faut alors \emph{vérifier à la fin} que la fraction obtenue est bien irréductible. Le risque : s'arrêter trop tôt ! Par exemple, si tu simplifies $\dfrac{84}{126}$ seulement par 2, tu obtiens $\dfrac{42}{63}$, qui est encore simplifiable. L'avantage de la méthode du PGCD est qu'elle donne la forme irréductible \emph{en une seule fois}, avec la garantie mathématique qu'on ne peut plus rien simplifier.
\end{attention}

\begin{exoresolu}[Écrire $\dfrac{84}{126}$ sous forme irréductible]
\textbf{Énoncé.} Un champ de maïs de la région de l'Ouest est partagé en 126 parcelles, dont 84 sont déjà récoltées. Écrire la fraction $\dfrac{84}{126}$ sous forme irréductible.
\tcblower
\textbf{Solution détaillée.}

\textbf{Étape 1 : calcul du PGCD par l'algorithme d'Euclide.}
\begin{align*}
126 &= 84 \times 1 + 42 \\
84 &= 42 \times 2 + 0
\end{align*}
Le dernier reste non nul est 42, donc $\mathrm{PGCD}(84\,;126) = 42$.

\textbf{Étape 2 : division du numérateur et du dénominateur par 42.}
\[ \frac{84}{126} = \frac{84 \div 42}{126 \div 42} = \frac{2}{3}. \]

\textbf{Étape 3 : conclusion.} Comme $\mathrm{PGCD}(2\,;3) = 1$, la fraction $\dfrac{2}{3}$ est irréductible.

On peut vérifier la simplification étape par étape : $42 = 2 \times 3 \times 7$, donc on peut diviser successivement par 2, par 3 puis par 7 :
\[ \frac{84}{126} \xrightarrow{\div 2} \frac{42}{63} \xrightarrow{\div 3} \frac{14}{21} \xrightarrow{\div 7} \frac{2}{3}. \]
Les $\dfrac{2}{3}$ du champ sont donc déjà récoltés.
\end{exoresolu}

\begin{popfigure}
\begin{center}
\begin{tikzpicture}[
  every node/.style={font=\large},
  frac/.style={draw=popblue, fill=popblueL, rounded corners=3pt, inner sep=5pt},
  fleche/.style={-{Stealth[length=3mm]}, thick, poporange},
  etiq/.style={midway, right, font=\small\itshape, text=popdark}
]
\node[frac] (f1) at (0,0) {$\dfrac{84}{126}$};
\node[frac] (f2) at (0,-1.8) {$\dfrac{42}{63}$};
\node[frac] (f3) at (0,-3.6) {$\dfrac{14}{21}$};
\node[frac, draw=popgreen, fill=popgreenL] (f4) at (0,-5.4) {$\dfrac{2}{3}$};
\draw[fleche] (f1) -- (f2) node[etiq] {$\div\,2$};
\draw[fleche] (f2) -- (f3) node[etiq] {$\div\,3$};
\draw[fleche] (f3) -- (f4) node[etiq] {$\div\,7$};
\node[font=\small, text=poppurple, align=center] at (5.2,-2.7)
  {$\mathrm{PGCD}(84\,;126) = 42$\\[2pt] $42 = 2 \times 3 \times 7$};
\draw[fleche, poppurple, dashed] (f1.east) to[bend left=15] node[midway, above, font=\small\itshape, text=poppurple]{$\div\,42$} (f4.east);
\end{tikzpicture}
\end{center}
\legende{Simplification de $\frac{84}{126}$ : en trois étapes ($\div 2$, $\div 3$, $\div 7$) ou en une seule grâce au PGCD ($\div 42$).}
\end{popfigure}

\courssub{Pourquoi cette notion est-elle si utile ?}

\begin{aretenir}[L'essentiel de la section]
\begin{itemize}
\item Deux entiers $a$ et $b$ sont \cle{premiers entre eux} si $\mathrm{PGCD}(a\,;b) = 1$.
\item Deux nombres premiers entre eux ne sont pas forcément premiers chacun (exemple : 15 et 28).
\item Si $d = \mathrm{PGCD}(a\,;b)$, alors $a = d \times a'$ et $b = d \times b'$ avec $a'$ et $b'$ premiers entre eux.
\item Pour rendre une fraction irréductible, on divise le numérateur et le dénominateur par leur PGCD : $\dfrac{84}{126} = \dfrac{2}{3}$ car $\mathrm{PGCD}(84\,;126) = 42$.
\end{itemize}
\end{aretenir}

\begin{savaistu}[Les fractions irréductibles dans la vie quotidienne]
Quand un élève obtient 36 points sur 48 à une évaluation, son professeur peut dire qu'il a réussi « les trois quarts » de l'épreuve : en effet $\dfrac{36}{48} = \dfrac{3}{4}$ car $\mathrm{PGCD}(36\,;48) = 12$. De même, en musique, les rythmes s'écrivent avec des fractions, et les artisans (menuisiers, tailleurs) simplifient leurs mesures pour travailler plus facilement. Les nombres premiers entre eux jouent aussi un rôle central en informatique moderne : la sécurité des transactions Mobile Money repose sur des méthodes de chiffrement qui utilisent de grands nombres premiers entre eux !
\end{savaistu}

\begin{exercice}[À toi de jouer]
\begin{enumerate}
\item Les couples suivants sont-ils formés de nombres premiers entre eux ? Justifier.
\begin{enumerate}
\item 24 et 35 \qquad b) 27 et 45 \qquad c) 8 et 15
\end{enumerate}
\item Écrire sous forme irréductible les fractions : $\dfrac{48}{72}$ ; $\dfrac{105}{147}$ ; $\dfrac{90}{126}$.
\end{enumerate}
\end{exercice}

\begin{corrige}[Exercice 1]
\begin{enumerate}
\item a) $\mathrm{PGCD}(24\,;35) = 1$ : oui, premiers entre eux. \quad b) $\mathrm{PGCD}(27\,;45) = 9 \neq 1$ : non. \quad c) $\mathrm{PGCD}(8\,;15) = 1$ : oui.
\item $\dfrac{48}{72} = \dfrac{2}{3}$ (PGCD $= 24$) ; $\dfrac{105}{147} = \dfrac{5}{7}$ (PGCD $= 21$) ; $\dfrac{90}{126} = \dfrac{5}{7}$ (PGCD $= 18$).
\end{enumerate}
\end{corrige}

\coursec{PPCM et relation avec le PGCD}

Dans la section précédente, nous avons découvert les nombres \cle{premiers entre eux} : deux entiers dont le PGCD vaut $1$, comme $8$ et $15$. Nous savons donc désormais trouver le plus grand diviseur commun de deux entiers. Mais la vie quotidienne pose souvent la question « miroir » : non pas « qu'est-ce qui divise à la fois $a$ et $b$ ? », mais « quel est le plus petit nombre que $a$ et $b$ divisent tous les deux ? ». C'est exactement la situation des deux cars de notre accroche : un car part pour Douala toutes les $12$ minutes, un autre part pour Bafoussam toutes les $18$ minutes ; à quelle heure repartiront-ils ensemble ? La réponse est un \cle{multiple commun} à $12$ et à $18$, et plus précisément le plus petit d'entre eux : le \cle{PPCM}.

\courssub{Multiples communs et définition du PPCM}

\begin{experience}[Les départs des cars à la gare routière]
À la gare routière de Mboppi, un car part pour Douala toutes les $12$ minutes à partir de 6 h 00, et un car part pour Bafoussam toutes les $18$ minutes à partir de 6 h 00. Écrivons les horaires de départ (en minutes après 6 h) :
\begin{itemize}
\item départs pour Douala : $0$ ; $12$ ; $24$ ; $36$ ; $48$ ; $60$ ; $72$ ; \ldots
\item départs pour Bafoussam : $0$ ; $18$ ; $36$ ; $54$ ; $72$ ; $90$ ; \ldots
\end{itemize}
Les instants de départ simultané sont les nombres présents dans les deux listes : $0$ ; $36$ ; $72$ ; \ldots\ Le premier départ commun après 6 h 00 a lieu $36$ minutes plus tard. Ce nombre $36$ est le plus petit multiple commun non nul de $12$ et $18$.
\end{experience}

\begin{definition}[PPCM de deux entiers naturels]
Soient $a$ et $b$ deux entiers naturels non nuls. Un \cle{multiple commun} de $a$ et $b$ est un entier qui est à la fois multiple de $a$ et multiple de $b$. Parmi tous les multiples communs \emph{non nuls} de $a$ et $b$, le plus petit est appelé le \cle{PPCM} de $a$ et $b$ (Plus Petit Commun Multiple). On le note $\mathrm{PPCM}(a\,;\,b)$.
\end{definition}

\begin{exemplebox}[Multiples communs de 12 et 18]
Multiples de $12$ : $12$ ; $24$ ; $\mathbf{36}$ ; $48$ ; $60$ ; $\mathbf{72}$ ; \ldots

Multiples de $18$ : $18$ ; $\mathbf{36}$ ; $54$ ; $\mathbf{72}$ ; $90$ ; \ldots

Les multiples communs non nuls sont $36$ ; $72$ ; $108$ ; \ldots\ Le plus petit est $36$ :
\[ \mathrm{PPCM}(12\,;\,18) = 36. \]
Remarque : les multiples communs de $12$ et $18$ sont exactement les multiples de $36$.
\end{exemplebox}

\begin{popfigure}
\begin{tikzpicture}[scale=1]
% Axe des minutes
\draw[->,thick,popdark] (0,0) -- (13.2,0) node[below left=2pt and -10pt] {\footnotesize minutes après 6 h};
\foreach \x in {0,1,...,12}{
  \draw[popmuted] (\x,0.06) -- (\x,-0.06);
}
\foreach \x/\lab in {0/0, 3/12, 6/24, 9/36, 12/48}{
  \node[below,popdark,font=\footnotesize] at (\x,-0.08) {\lab};
}
% Multiples de 12 (tous les 3 unités)
\foreach \x in {0,3,6,9,12}{
  \draw[popblue,thick] (\x,0) -- (\x,0.45);
  \fill[popblue] (\x,0.5) circle (0.09);
}
\node[popblue,font=\footnotesize,anchor=west] at (12.4,0.5) {départs Douala (12 min)};
% Multiples de 18 (tous les 4.5 unités)
\foreach \x in {0,4.5,9}{
  \draw[poporange,thick] (\x,0) -- (\x,-0.45);
  \fill[poporange] (\x,-0.55) circle (0.09);
}
\node[poporange,font=\footnotesize,anchor=west] at (10.2,-0.75) {départs Bafoussam (18 min)};
% Coïncidence à 36
\draw[popgreen,very thick,dashed] (9,0.85) -- (9,-0.95);
\fill[popgreen] (9,0) circle (0.13);
\node[popgreen,font=\footnotesize\bfseries,align=center] at (9,1.15) {départ commun : 36 min};
\end{tikzpicture}
\legende{Les graduations bleues (multiples de $12$) et oranges (multiples de $18$) coïncident pour la première fois à $36$ minutes : $\mathrm{PPCM}(12\,;\,18)=36$.}
\end{popfigure}

\courssub{Deux méthodes directes pour trouver le PPCM}

\begin{methode}[Méthode 1 : les listes de multiples]
Pour trouver $\mathrm{PPCM}(a\,;\,b)$ :
\begin{enumerate}
\item écrire la liste des premiers multiples de $a$ ;
\item écrire la liste des premiers multiples de $b$ ;
\item repérer le premier nombre (non nul) présent dans les deux listes : c'est le PPCM.
\end{enumerate}
Cette méthode est rapide pour de petits nombres, mais devient pénible dès que les nombres grandissent.
\end{methode}

\begin{methode}[Méthode 2 : la décomposition en facteurs premiers]
Pour trouver $\mathrm{PPCM}(a\,;\,b)$ :
\begin{enumerate}
\item décomposer $a$ et $b$ en produits de facteurs premiers ;
\item pour \textbf{chaque} facteur premier apparaissant dans l'une \emph{ou} l'autre décomposition, retenir \textbf{le plus grand exposant} ;
\item multiplier tous ces facteurs : le produit obtenu est le PPCM.
\end{enumerate}
À comparer avec le PGCD, où l'on ne retient que les facteurs \emph{communs} avec le \emph{plus petit} exposant.
\end{methode}

\begin{exemplebox}[PPCM de 12 et 18 par les facteurs premiers]
$12 = 2^2 \times 3$ et $18 = 2 \times 3^2$.
\begin{itemize}
\item facteur $2$ : exposants $2$ et $1$, on retient $2^2$ ;
\item facteur $3$ : exposants $1$ et $2$, on retient $3^2$.
\end{itemize}
\[ \mathrm{PPCM}(12\,;\,18) = 2^2 \times 3^2 = 4 \times 9 = 36. \]
\end{exemplebox}

\courssub{Le théorème fondamental : lien entre PGCD et PPCM}

Voici maintenant le résultat le plus important de cette section, celui qui relie les deux grandeurs étudiées dans cette leçon.

\begin{propriete}[Théorème : relation entre PGCD et PPCM]
Pour tous entiers naturels non nuls $a$ et $b$ :
\[ \mathrm{PGCD}(a\,;\,b) \times \mathrm{PPCM}(a\,;\,b) = a \times b. \]
Cette égalité est \emph{exigible} au BEPC : tu dois savoir l'énoncer et l'utiliser. Sa justification dans le cas général est admise, mais nous allons la comprendre sur un exemple grâce aux décompositions en facteurs premiers.
\end{propriete}

\begin{experience}[Pourquoi ça marche ? L'idée des exposants min et max]
Prenons $a = 84$ et $b = 126$. Leurs décompositions en facteurs premiers sont :
\[ 84 = 2^2 \times 3^1 \times 7^1 \qquad\text{et}\qquad 126 = 2^1 \times 3^2 \times 7^1. \]
Imaginons chaque décomposition comme un « sac » contenant des facteurs premiers.

\textbf{Étape 1.} Le PGCD prend, pour chaque facteur premier commun, le \emph{plus petit} exposant (le « minimum ») :
\[ \mathrm{PGCD}(84\,;\,126) = 2^1 \times 3^1 \times 7^1 = 42. \]

\textbf{Étape 2.} Le PPCM prend, pour chaque facteur premier, le \emph{plus grand} exposant (le « maximum ») :
\[ \mathrm{PPCM}(84\,;\,126) = 2^2 \times 3^2 \times 7^1 = 252. \]

\textbf{Étape 3.} Multiplions PGCD et PPCM. Pour chaque facteur premier, les exposants s'additionnent :
\[ \underbrace{\min(p,q)}_{\text{dans le PGCD}} + \underbrace{\max(p,q)}_{\text{dans le PPCM}} = p + q, \]
car entre deux nombres $p$ et $q$, l'un est le plus petit et l'autre le plus grand : leur somme vaut toujours $p+q$. Ainsi, pour le facteur $2$ : $\min(2,1)+\max(2,1) = 1+2 = 3 = 2+1$. Pour le facteur $3$ : $1+2 = 3$. Pour le facteur $7$ : $1+1 = 2$.

\textbf{Étape 4.} On obtient donc :
\[ \mathrm{PGCD} \times \mathrm{PPCM} = 2^{1+2} \times 3^{1+2} \times 7^{1+1} = 2^3 \times 3^3 \times 7^2. \]
Or $a \times b = (2^2 \times 3 \times 7) \times (2 \times 3^2 \times 7) = 2^3 \times 3^3 \times 7^2$. C'est le même produit !

\textbf{Conclusion :} $\mathrm{PGCD}(84\,;\,126) \times \mathrm{PPCM}(84\,;\,126) = 84 \times 126$. Vérifions : $42 \times 252 = \num{10584}$ et $84 \times 126 = \num{10584}$. Le même raisonnement vaut pour n'importe quels entiers $a$ et $b$ : le théorème est admis dans le cas général.
\end{experience}

\begin{popfigure}
\begin{tikzpicture}[scale=1]
% Sac 84
\draw[thick,popblue,rounded corners=6pt] (0,0) rectangle (3.4,2.6);
\node[popblue,font=\small\bfseries] at (1.7,2.25) {$84 = 2^2 \times 3 \times 7$};
\node[draw,popblue,fill=popblueL,circle,font=\footnotesize] (a1) at (0.8,1.5) {$2$};
\node[draw,popblue,fill=popblueL,circle,font=\footnotesize] (a2) at (1.7,1.5) {$2$};
\node[draw,popgreen,fill=popgreenL,circle,font=\footnotesize] (a3) at (2.6,1.5) {$3$};
\node[draw,poporange,fill=poporangeL,circle,font=\footnotesize] (a4) at (1.25,0.6) {$7$};
% Sac 126
\draw[thick,poppurple,rounded corners=6pt] (4.4,0) rectangle (7.8,2.6);
\node[poppurple,font=\small\bfseries] at (6.1,2.25) {$126 = 2 \times 3^2 \times 7$};
\node[draw,popblue,fill=popblueL,circle,font=\footnotesize] (b1) at (5.2,1.5) {$2$};
\node[draw,popgreen,fill=popgreenL,circle,font=\footnotesize] (b2) at (6.1,1.5) {$3$};
\node[draw,popgreen,fill=popgreenL,circle,font=\footnotesize] (b3) at (7.0,1.5) {$3$};
\node[draw,poporange,fill=poporangeL,circle,font=\footnotesize] (b4) at (5.65,0.6) {$7$};
% Résultats
\node[draw,popdark,fill=popcream,rounded corners=4pt,align=center,font=\footnotesize] at (1.7,-1.3)
 {PGCD : on prend le \textbf{minimum}\\ $2^1 \times 3^1 \times 7^1 = 42$};
\node[draw,popdark,fill=popcream,rounded corners=4pt,align=center,font=\footnotesize] at (6.1,-1.3)
 {PPCM : on prend le \textbf{maximum}\\ $2^2 \times 3^2 \times 7^1 = 252$};
\draw[->,thick,popdark] (1.7,-0.15) -- (1.7,-0.75);
\draw[->,thick,popdark] (6.1,-0.15) -- (6.1,-0.75);
\node[align=center,font=\footnotesize,popdark] at (3.9,-2.35)
 {Pour chaque facteur : $\min(p,q) + \max(p,q) = p+q$ \quad$\Rightarrow$\quad $42 \times 252 = 84 \times 126$};
\end{tikzpicture}
\legende{Les « sacs de facteurs premiers » de $84$ et $126$ : le PGCD garde le plus petit exposant de chaque facteur, le PPCM le plus grand. Réunis, ils reconstituent exactement le produit $a \times b$.}
\end{popfigure}

\courssub{Méthode 3 : calculer le PPCM à partir du PGCD}

Le théorème précédent fournit un raccourci très puissant : si l'on connaît déjà le PGCD (par exemple grâce à l'algorithme d'Euclide, vu dans la section 4), le PPCM s'obtient par une simple division.

\begin{methode}[Calculer le PPCM grâce au PGCD]
Pour trouver $\mathrm{PPCM}(a\,;\,b)$ :
\begin{enumerate}
\item calculer $\mathrm{PGCD}(a\,;\,b)$ (par l'algorithme d'Euclide, le plus souvent) ;
\item appliquer la formule :
\[ \mathrm{PPCM}(a\,;\,b) = \frac{a \times b}{\mathrm{PGCD}(a\,;\,b)}. \]
\end{enumerate}
C'est la méthode la plus efficace pour de grands nombres, car elle évite les décompositions en facteurs premiers.
\end{methode}

\begin{exoresolu}[PPCM de 84 et 126]
Énoncé. Sachant que $\mathrm{PGCD}(84\,;\,126) = 42$, calculer $\mathrm{PPCM}(84\,;\,126)$.
\tcblower
Solution détaillée. On applique la relation fondamentale :
\[ \mathrm{PPCM}(84\,;\,126) = \frac{84 \times 126}{\mathrm{PGCD}(84\,;\,126)} = \frac{84 \times 126}{42}. \]
Astuce de calcul : simplifions avant de multiplier. Comme $84 = 42 \times 2$, on a :
\[ \mathrm{PPCM}(84\,;\,126) = \frac{42 \times 2 \times 126}{42} = 2 \times 126 = 252. \]
Conclusion : $\mathrm{PPCM}(84\,;\,126) = 252$. Vérification avec la méthode 2 : $84 = 2^2 \times 3 \times 7$ et $126 = 2 \times 3^2 \times 7$, donc $\mathrm{PPCM} = 2^2 \times 3^2 \times 7 = 4 \times 9 \times 7 = 252$. Les deux méthodes concordent.
\end{exoresolu}

\begin{attention}[Piège : le PPCM n'est pas toujours le produit $a \times b$]
Beaucoup d'élèves écrivent trop vite $\mathrm{PPCM}(a\,;\,b) = a \times b$. C'est \textbf{faux en général} ! Par exemple :
\[ \mathrm{PPCM}(12\,;\,18) = 36 \quad\text{et non pas}\quad 12 \times 18 = 216. \]
D'après le théorème, $\mathrm{PPCM}(a\,;\,b) = \dfrac{a \times b}{\mathrm{PGCD}(a\,;\,b)}$. Le PPCM est donc égal au produit $a \times b$ \textbf{uniquement} lorsque $\mathrm{PGCD}(a\,;\,b) = 1$, c'est-à-dire lorsque $a$ et $b$ sont \cle{premiers entre eux} (section précédente). Exemple : $\mathrm{PPCM}(8\,;\,15) = 8 \times 15 = 120$, car $8$ et $15$ sont premiers entre eux.
\end{attention}

\courssub{Application : résolution du problème des cars}

Revenons à la situation de l'accroche, que nous pouvons maintenant résoudre complètement et rédiger proprement, comme on te le demandera au BEPC.

\begin{exoresolu}[Les cars de la gare routière]
Énoncé. À la gare routière, un car part pour Douala toutes les $12$ minutes et un car part pour Bafoussam toutes les $18$ minutes. À 6 h 00, deux cars partent en même temps. À quelle heure deux cars repartiront-ils ensemble pour la première fois ?
\tcblower
Solution détaillée.

\textbf{Analyse.} Les départs pour Douala ont lieu aux multiples de $12$ (minutes après 6 h), ceux pour Bafoussam aux multiples de $18$. Un départ simultané correspond à un multiple commun de $12$ et $18$ ; le \emph{premier} départ commun correspond au \emph{plus petit} d'entre eux : il faut calculer $\mathrm{PPCM}(12\,;\,18)$.

\textbf{Calcul.} Décomposons : $12 = 2^2 \times 3$ et $18 = 2 \times 3^2$. En retenant le plus grand exposant de chaque facteur :
\[ \mathrm{PPCM}(12\,;\,18) = 2^2 \times 3^2 = 4 \times 9 = 36. \]

\textbf{Conclusion.} Les deux cars repartent ensemble $36$ minutes après 6 h 00, c'est-à-dire à \textbf{6 h 36}. Les départs communs suivants auront lieu toutes les $36$ minutes : 7 h 12, 7 h 48, \ldots
\end{exoresolu}

\begin{aretenir}[L'essentiel sur le PPCM]
\begin{itemize}
\item Le \cle{PPCM} de deux entiers non nuls $a$ et $b$ est le plus petit multiple commun \emph{non nul} de $a$ et $b$.
\item Trois méthodes : les listes de multiples ; la décomposition en facteurs premiers (plus grand exposant pour chaque facteur) ; la formule $\mathrm{PPCM}(a\,;\,b) = \dfrac{a \times b}{\mathrm{PGCD}(a\,;\,b)}$.
\item Relation fondamentale : $\mathrm{PGCD}(a\,;\,b) \times \mathrm{PPCM}(a\,;\,b) = a \times b$.
\item $\mathrm{PPCM}(a\,;\,b) = a \times b$ seulement si $a$ et $b$ sont premiers entre eux.
\item Le PPCM répond aux questions du type : « au bout de combien de temps deux événements périodiques se reproduisent-ils en même temps ? ».
\end{itemize}
\end{aretenir}

\begin{savaistu}[Des engrenages aux cycles : le PPCM autour de nous]
Le PPCM sert à prévoir les coïncidences de cycles dans bien des domaines de la vie courante : dans un engrenage, deux dents marquées se retrouvent face à face au bout d'un nombre de tours donné par un PPCM ; il permet aussi de planifier des retrouvailles périodiques (festivals, grandes rencontres traditionnelles) ou de synchroniser des feux tricolores. Les astronomes utilisent le même principe pour prévoir le retour des éclipses, qui combinent plusieurs cycles célestes. Plus tard, dans tes études de mathématiques, tu découvriras que le PGCD et le PPCM sont à la base d'outils plus avancés qui servent aujourd'hui à sécuriser les transactions Mobile Money (MTN MoMo, Orange Money) grâce à la cryptographie. \emph{(Ces notions sont hors-programme en 3ème : c'est une ouverture culturelle.)}
\end{savaistu}

\begin{exercice}[À toi de jouer]
\begin{enumerate}
\item Calculer $\mathrm{PPCM}(15\,;\,20)$ par la méthode des listes, puis vérifier avec la relation fondamentale sachant que $\mathrm{PGCD}(15\,;\,20) = 5$.
\item Deux guirlandes électriques clignotent : l'une s'allume toutes les $9$ secondes, l'autre toutes les $14$ secondes. Elles s'allument ensemble à un instant donné. Au bout de combien de secondes s'allumeront-elles de nouveau ensemble ? Que remarques-tu sur $9$ et $14$ ?
\item Sachant que $\mathrm{PGCD}(360\,;\,840) = 120$, calculer $\mathrm{PPCM}(360\,;\,840)$ sans décomposer en facteurs premiers.
\end{enumerate}
\end{exercice}

\begin{corrige}[Exercice n°1]
\begin{enumerate}
\item Multiples de $15$ : $15$ ; $30$ ; $45$ ; $\mathbf{60}$ ; \ldots\ Multiples de $20$ : $20$ ; $40$ ; $\mathbf{60}$ ; \ldots\ Donc $\mathrm{PPCM}(15\,;\,20) = 60$. Vérification : $\dfrac{15 \times 20}{5} = \dfrac{300}{5} = 60$.
\item Il faut calculer $\mathrm{PPCM}(9\,;\,14)$. Or $9 = 3^2$ et $14 = 2 \times 7$ n'ont aucun facteur premier commun : ils sont premiers entre eux, donc $\mathrm{PPCM}(9\,;\,14) = 9 \times 14 = 126$ secondes (soit $2$ min $6$ s).
\item $\mathrm{PPCM}(360\,;\,840) = \dfrac{360 \times 840}{120} = 360 \times 7 = \num{2520}$.
\end{enumerate}
\end{corrige}

\coursec{Résolution de problèmes concrets et synthèse}

Après avoir maîtrisé le calcul du \cle{PGCD} par les deux algorithmes et la détermination du \cle{PPCM} grâce à la relation fondamentale $\text{PGCD}(a,b) \times \text{PPCM}(a,b) = a \times b$, nous abordons maintenant l'étape décisive : \textbf{savoir choisir le bon outil et rédiger une solution complète, justifiée et concluante}. C'est cette compétence que le jury du \cle{BEPC} évalue prioritairement : non seulement le calcul correct, mais la démonstration que vous comprenez \emph{why} vous calculez un PGCD ou un PPCM.

\courssub{Grille de décision : PGCD ou PPCM ?}

Face à un énoncé concret, deux grandes familles de situations se présentent. Les reconnaître instantanément évite les erreurs d'orientation.

\begin{definition}[Les deux familles de situations types]
\emph{Famille A — Partage équitable en parts identiques maximales (\cle{PGCD}).} On doit diviser une quantité (longueur, surface, nombre d'objets, somme d'argent) en parts \textbf{identiques} et \textbf{aussi grandes que possible}, sans reste. Mots-clés : « plus grand », « plus grande », « maximum », « identique », « sans reste », « pavage », « carrés les plus grands », « groupes de même effectif maximal ».

\emph{Famille B — Coïncidence d'événements périodiques (\cle{PPCM}).} Des phénomènes se répètent avec des périodes différentes ; on cherche le premier instant (ou la plus petite quantité) où ils se produisent \textbf{simultanément}. Mots-clés : « en même temps », « ensemble », « coïncident », « se retrouvent », « même jour », « alignement », « synchronisation », « plus petit commun multiple ».
\end{definition}

\begin{popfigure}
\begin{tikzpicture}[scale=0.85, every node/.style={font=\small}]
% Titre
\node[popdark, font=\bfseries\large] at (0,4.2) {Grille de décision : PGCD vs PPCM};

% Colonne gauche : PGCD
\begin{scope}[xshift=-5.5cm]
\node[popblue, font=\bfseries\large] at (0,3) {PGCD};
\node[popblue, font=\bfseries] at (0,2.5) {« Partage maximal »};
\draw[popblue, thick] (-2.5,2.2) rectangle (2.5,-2.5);
\node[align=center, text width=4.5cm] at (0,1.8) {\textbf{Mots-clés :}\\
plus grand / grande\\
maximum / maximal\\
identique / même\\
sans reste / exact\\
pavage / carrés\\
groupes de même effectif};
\node[align=center, text width=4.5cm, popgreen] at (0,0.2) {\textbf{Question type :}\\
« Quelle est la plus grande\\taille possible ? »\\
« Combien de parts ? »};
\node[align=center, text width=4.5cm, poporange] at (0,-1.4) {\textbf{Exemples :}\\
Pavage de dalles carrées\\
Coupage de fils/bandelettes\\
Groupes d'élèves identiques\\
Partage de bonbons\\sans reste};
\end{scope}

% Colonne droite : PPCM
\begin{scope}[xshift=5.5cm]
\node[poporange, font=\bfseries\large] at (0,3) {PPCM};
\node[poporange, font=\bfseries] at (0,2.5) {« Coïncidence de cycles »};
\draw[poporange, thick] (-2.5,2.2) rectangle (2.5,-2.5);
\node[align=center, text width=4.5cm] at (0,1.8) {\textbf{Mots-clés :}\\
en même temps / ensemble\\
se retrouvent / réunis\\
coïncident / alignés\\
synchronisation\\
premier jour commun\\
périodes / cycles};
\node[align=center, text width=4.5cm, popgreen] at (0,0.2) {\textbf{Question type :}\\
« Quand se retrouveront-ils ? »\\
« Au bout de combien de jours ? »\\
« Première coïncidence ? »};
\node[align=center, text width=4.5cm, popblue] at (0,-1.4) {\textbf{Exemples :}\\
Marchés / rendez-vous\\
Feux tricolores synchronisés\\
Planètes alignées\\
Maintenance machines\\
Horaires de bus};
\end{scope}

% Flèche centrale
\draw[popdark, very thick, ->] (-2.8,0) -- node[above, popdark, font=\bfseries] {ANALYSE DE L'ÉNONCÉ} (2.8,0);
\end{tikzpicture}
\legende{Tableau de synthèse : identifier la famille de situation pour choisir l'outil adapté.}
\end{popfigure}

\begin{methode}[Choisir entre PGCD et PPCM en 3 questions]
\begin{enumerate}
\item \textbf{Que cherche-t-on ?} Une \emph{taille maximale de part} (PGCD) ou un \emph{instant de rencontre} (PPCM) ?
\item \textbf{Y a-t-il « partage sans reste » ?} Oui $\rightarrow$ PGCD. \textbf{Y a-t-il « répétition périodique » ?} Oui $\rightarrow$ PPCM.
\item \textbf{Test rapide :} Remplacez les nombres par des valeurs très petites (ex. 2 et 3). Si la réponse logique est « 1 » (plus grand diviseur commun), c'est PGCD ; si c'est « 6 » (premier multiple commun), c'est PPCM.
\end{enumerate}
\end{methode}

\courssub{Problème type 1 — PGCD : Pavage d'une cour rectangulaire}

\begin{exemplebox}[Pavage par des dalles carrées maximales]
Une cour rectangulaire mesure $5,40\,\text{m}$ sur $3,60\,\text{m}$. On veut la paver entièrement avec des dalles \textbf{carrées}, toutes de la \textbf{même dimension}, la \textbf{plus grande possible}, sans couper de dalle.
\begin{enumerate}
\item Quelle est la longueur de côté de ces dalles (en cm) ?
\item Combien de dalles faut-il ?
\end{enumerate}
\end{exemplebox}

\textbf{Résolution complète et justifiée.}\par

\begin{enumerate}
\item \textbf{Traduction mathématique.}
Les dimensions en mètres donnent des décimales ; on passe en centimètres pour travailler avec des entiers :
\[
5,40\,\text{m} = 540\,\text{cm},\qquad 3,60\,\text{m} = 360\,\text{cm}.
\]
Paver avec des carrés identiques de côté $c$ (entier) sans reste signifie que $c$ divise $540$ et $c$ divise $360$. De plus, on veut $c$ \textbf{maximum}. Donc $c = \text{PGCD}(540,\,360)$.

\item \textbf{Choix de l'algorithme et calcul.}
Les nombres sont de taille moyenne ; l'algorithme d'Euclide est le plus efficace.
\begin{align*}
540 &= 360 \times 1 + 180 \\
360 &= 180 \times 2 + 0
\end{align*}
Le dernier reste non nul est $180$. Donc $\text{PGCD}(540,\,360) = 180$.

\item \textbf{Interprétation.}
La plus grande dimension possible pour les dalles carrées est \textbf{$180\,\text{cm}$} (soit $1,80\,\text{m}$).

\item \textbf{Nombre de dalles.}
Sur la longueur $540\,\text{cm}$ : $540 \div 180 = 3$ dalles.
Sur la largeur $360\,\text{cm}$ : $360 \div 180 = 2$ dalles.
Total : $3 \times 2 = \textbf{6 dalles}$.
\end{enumerate}

\begin{popfigure}
\begin{tikzpicture}[scale=0.0065, every node/.style={font=\tiny}]
% Rectangle global 540 x 360
\draw[popdark, very thick] (0,0) rectangle (540,360);

% Quadrillage des dalles 180 x 180 (3 par 2)
\foreach \x in {180,360} {
  \draw[popblue!60, thick] (\x,0) -- (\x,360);
}
\foreach \y in {180} {
  \draw[popblue!60, thick] (0,\y) -- (540,\y);
}

% Hachures légères sur quelques dalles pour montrer le pavage
\foreach \i/\j in {0/0, 1/0, 2/0, 0/1, 2/1} {
  \fill[popblue!10, pattern=north east lines, pattern color=popblue!40]
    (\i*180,\j*180) rectangle (\i*180+180,\j*180+180);
}

% Cotes
\draw[|<->|, popdark] (0,-25) -- (540,-25) node[midway, below=3pt, font=\small\bfseries] {540 cm};
\draw[|<->|, popdark] (-25,0) -- (-25,360) node[midway, left=3pt, font=\small\bfseries, rotate=90] {360 cm};

% Cote d'une dalle
\draw[|<->|, poporange, thick] (0,385) -- (180,385) node[midway, above=3pt, font=\small\bfseries, poporange] {180 cm};
\draw[|<->|, poporange, thick] (565,0) -- (565,180) node[midway, right=3pt, font=\small\bfseries, poporange, rotate=-90] {180 cm};

% Légende
\node[popdark, font=\bfseries\small] at (270, -80) {Pavage de la cour par 6 dalles carrées de 180 cm de côté (PGCD = 180)};
\end{tikzpicture}
\legende{La cour de $540\times360$ cm est exactement pavée par $3\times2=6$ carrés de côté $180$ cm (le PGCD).}
\end{popfigure}

\begin{exemplebox}[Vérification du résultat]
\begin{itemize}
\item $180$ divise $540$ ? $540 = 180 \times 3$ ✓
\item $180$ divise $360$ ? $360 = 180 \times 2$ ✓
\item Existe-t-il un diviseur commun plus grand ? Non, car $180$ est le PGCD.
\item Nombre de dalles entier ? $3\times2=6$ ✓
\end{itemize}
\end{exemplebox}

\begin{attention}[Piège classique à l'examen]
\textbf{Ne vous arrêtez JAMAIS à « PGCD = 180 ».} La question demande souvent : « Combien de dalles faut-il ? » ou « Quelle est la longueur de côté ? » avec l'unité.
\begin{itemize}
\item \textbf{Mauvaise réponse :} « PGCD(540,360) = 180. » (Incomplète, pas de conclusion, pas d'unité.)
\item \textbf{Bonne réponse :} « La plus grande dimension possible pour les dalles est \textbf{180 cm}. Il faudra \textbf{6 dalles} (3 dans la longueur, 2 dans la largeur). »
\end{itemize}
\textbf{À retenir :} Une solution complète = Calcul justifié + Conclusion en phrase + Unité.
\end{attention}

\courssub{Problème type 2 — PPCM : Rendez-vous au marché}

\begin{exemplebox}[Trois amis au marché]
Trois amis, Amadou, Binta et Chantal, se rendent au même marché. Amadou y va \textbf{tous les 4 jours}, Binta \textbf{tous les 6 jours}, Chantal \textbf{tous les 10 jours}. Ils s'y sont rencontrés \textbf{aujourd'hui} (jour 0). Dans combien de jours se retrouveront-ils à nouveau tous les trois au marché ?
\end{exemplebox}

\textbf{Résolution complète et justifiée.}\par

\begin{enumerate}
\item \textbf{Analyse de la situation.}
Chacun a une périodicité : $4$, $6$, $10$ jours. Ils se retrouvent quand le nombre de jours écoulés est un multiple de $4$, de $6$ et de $10$ \textbf{simultanément}. On cherche le \textbf{plus petit} de ces multiples communs (strictement positif). C'est le \textbf{PPCM}(4, 6, 10).

\item \textbf{Calcul du PPCM.}
Méthode par décomposition en facteurs premiers (efficace pour plusieurs nombres) :
\begin{align*}
4 &= 2^2 \\
6 &= 2 \times 3 \\
10 &= 2 \times 5
\end{align*}
Le PPCM prend chaque facteur premier avec son \textbf{plus grand exposant} :
\[
\text{PPCM}(4,6,10) = 2^2 \times 3 \times 5 = 4 \times 3 \times 5 = 60.
\]

\item \textbf{Alternative via la relation PGCD-PPCM (deux nombres à la fois).}
$\text{PPCM}(4,6) = \frac{4\times6}{\text{PGCD}(4,6)} = \frac{24}{2} = 12$.
$\text{PPCM}(12,10) = \frac{12\times10}{\text{PGCD}(12,10)} = \frac{120}{2} = 60$.

\item \textbf{Conclusion.}
Ils se retrouveront tous les trois au marché \textbf{dans 60 jours}.
\end{enumerate}

\begin{methode}[Rédiger une solution complète en 4 étapes (standard BEPC)]
\begin{enumerate}
\item \textbf{TRADUIRE} : Identifier les grandeurs, passer en unités homogènes (entiers), nommer l'inconnue, écrire la condition mathématique (diviseur commun maximal $\leftrightarrow$ PGCD ; multiple commun minimal $\leftrightarrow$ PPCM).
\item \textbf{CHOISIR ET CALCULER} : Citer la propriété ou l'algorithme utilisé (Euclide, soustractions, décomposition en facteurs, relation PGCD$\times$PPCM). Présenter le calcul \textbf{étape par étape}, aligné, lisible.
\item \textbf{VÉRIFIER} : Contrôler que le résultat satisfait la condition (divise bien les deux nombres pour PGCD ; est bien multiple des deux pour PPCM). Vérifier l'ordre de grandeur.
\item \textbf{CONCLURE} : Répondre \textbf{précisément à la question posée}, en \textbf{phrase affirmative}, avec l'\textbf{unité} adaptée (cm, jours, dalles, francs CFA...).
\end{enumerate}
\end{methode}

\courssub{Synthèse : Les trois méthodes de calcul du PGCD}

\begin{popfigure}
\begin{center}
\begin{tabularx}{\linewidth}{|l|X|X|X|X|}
\hline
\rowcolor{popblueL}
\textbf{Méthode} & \textbf{Principe} & \textbf{Domaines d'efficacité} & \textbf{Avantages} & \textbf{Inconvénients / Limites} \\
\hline
Listes de diviseurs &
Lister tous les diviseurs de chaque nombre, garder les communs, prendre le max &
Nombres \textbf{très petits} ($<50$) ; nombres à diviseurs évidents &
Simple, visuel, pas d'algorithme &
Impossible pour grands nombres ; risque d'oubli ; lent \\
\hline
Algorithme des soustractions &
Remplacer le plus grand par la différence jusqu'à égalité : PGCD$(a,b)=$PGCD$(a-b,b)$ &
Nombres \textbf{proches} ; compréhension du concept ; examen sans calculatrice &
Montre la logique du PGCD ; pas de division requise &
Très lent si nombres d'échelles très différentes (ex. 10000 et 3) \\
\hline
Algorithme d'Euclide &
Divisions successives : $a=bq+r$ ; PGCD$(a,b)=$PGCD$(b,r)$ jusqu'à $r=0$ &
\textbf{Tous les cas} ; \emph{méthode de référence} au BEPC ; grands nombres &
Rapide (logarithmique) ; standard mondial ; base du RSA &
Nécessite la division euclidienne ; moins intuitif au début \\
\hline
\end{tabularx}
\end{center}
\legende{Tableau récapitulatif des trois méthodes de PGCD. À l'examen, l'algorithme d'Euclide est attendu pour les nombres de taille courante.}
\end{popfigure}

\courssub{Exercices résolus types BEPC}

\begin{exoresolu}[Exercice 1 — Partage de fruits]
Un marchand a $240$ mangues et $180$ oranges. Il veut faire des paniers identiques (même nombre de mangues, même nombre d'oranges par panier), en utilisant tous les fruits, avec le \textbf{maximum de paniers possibles}.
\begin{enumerate}
\item Combien de paniers peut-il faire ?
\item Combien de mangues et d'oranges par panier ?
\end{enumerate}
\tcblower
\textbf{Solution.}
\begin{enumerate}
\item \textbf{Analyse :} « Paniers identiques », « tous les fruits », « maximum de paniers » $\rightarrow$ partage maximal $\rightarrow$ \textbf{PGCD}.
\item \textbf{Calcul :} PGCD$(240,180)$ par Euclide.
$240 = 180 \times 1 + 60$
$180 = 60 \times 3 + 0$
PGCD $= 60$.
\item \textbf{Interprétation :} $60$ est le \textbf{nombre maximal de paniers}.
\item \textbf{Contenu par panier :} Mangues : $240 \div 60 = 4$. Oranges : $180 \div 60 = 3$.
\item \textbf{Conclusion :} Le marchand peut faire \textbf{60 paniers}, chacun contenant \textbf{4 mangues et 3 oranges}.
\end{enumerate}
\end{exoresolu}

\begin{exoresolu}[Exercice 2 — Feux tricolores synchronisés]
Deux feux tricolores sont sur la même route. Le premier passe au vert toutes les $45$ secondes, le second toutes les $75$ secondes. Ils sont verts \textbf{en même temps} à $08{:}00{:}00$ précises. À quelle heure seront-ils à nouveau verts simultanément ?
\tcblower
\textbf{Solution.}
\begin{enumerate}
\item \textbf{Analyse :} Périodicités $45$ s et $75$ s, recherche coïncidence $\rightarrow$ \textbf{PPCM}.
\item \textbf{Calcul :} PGCD$(45,75)$ par Euclide : $75=45\times1+30$ ; $45=30\times1+15$ ; $30=15\times2+0$ $\Rightarrow$ PGCD $=15$.
PPCM $= \frac{45\times75}{15} = 45\times5 = 225$ secondes.
\item \textbf{Conversion :} $225\,\text{s} = 3\,\text{min}\,45\,\text{s}$.
\item \textbf{Conclusion :} Les deux feux seront verts ensemble à \textbf{08{:}03{:}45}.
\end{enumerate}
\end{exoresolu}

\courssub{Bilan de l'unité Arithmétique}

\begin{aretenir}[Ce qu'il faut retenir pour le BEPC]
\begin{itemize}
\item \textbf{PGCD} : Plus Grand Commun Diviseur. Se calcule par \textbf{listes} (petits nombres), \textbf{soustractions} (nombres proches), \textbf{Euclide} (cas général, méthode de référence).
\item \textbf{PPCM} : Plus Petit Commun Multiple. Relation clé : $\text{PGCD}(a,b) \times \text{PPCM}(a,b) = a \times b$.
\item \textbf{Nombres premiers entre eux} : PGCD $= 1$ $\Leftrightarrow$ PPCM $= a \times b$.
\item \textbf{Rédiger une solution} : 1) Traduire $\rightarrow$ 2) Choisir l'outil + calcul justifié $\rightarrow$ 3) Vérifier $\rightarrow$ 4) Conclure (phrase + unité).
\item \textbf{Choix PGCD/PPCM} : Partage maximal / pavage $\rightarrow$ PGCD. Périodicités / rendez-vous $\rightarrow$ PPCM.
\end{itemize}
\end{aretenir}

\begin{savaistu}[PGCD, RSA et Mobile Money au Cameroun]
Saviez-vous que l'algorithme d'Euclide que vous venez d'apprendre est le \textbf{cœur mathématique du chiffrement RSA} qui protège vos transactions \textbf{Mobile Money (MTN MoMo, Orange Money)} et vos communications WhatsApp ?

Le principe : pour créer une clé publique, on choisit deux nombres premiers \emph{très grands} (plusieurs centaines de chiffres), $p$ et $q$. Leur produit $n=pq$ est public. La sécurité repose sur la difficulté de \textbf{retrouver $p$ et $q$ à partir de $n$} (factorisation). Mais pour \textbf{générer} les clés, il faut calculer des PGCD géants (ex. PGCD$(e, \varphi(n))=1$) et des inverses modulaires — ce qui utilise l'algorithme d'Euclide \emph{étendu}.

Chaque fois que vous envoyez 5\,000 FCFA par MoMo, votre téléphone exécute en quelques millisecondes des dizaines d'algorithmes d'Euclide sur des entiers de 2048 bits pour établir une session sécurisée. \textbf{Les mathématiques de 3ème sécurisent l'économie numérique du Cameroun !}
\end{savaistu}

\begin{exercice}[Entraînement BEPC — À vous de jouer !]
\begin{enumerate}
\item Un terrain rectangulaire de $72\,\text{m}$ sur $48\,\text{m}$ doit être divisé en carrés identiques les plus grands possibles. Donnez le côté du carré (en m) et le nombre de carrés.
\item Trois bus quittent la gare routière de Yaoundé à 6h00. Le bus A revient toutes les 12 min, le B toutes les 18 min, le C toutes les 30 min. À quelle heure repartiront-ils ensemble ?
\item Deux nombres ont pour PGCD 12 et pour PPCM 360. L'un des nombres est 60. Trouvez l'autre.
\item Un boulanger a 360 g de farine et 270 g de beurre. Il veut faire des parts de pâte identiques (même masse farine, même masse beurre), le plus grand nombre possible. Masse d'une part ? Nombre de parts ?
\end{enumerate}
\end{exercice}

\begin{corrige}[Exercice 1]
PGCD$(72,48)$ : $72=48\times1+24$ ; $48=24\times2+0$ $\Rightarrow$ PGCD $=24$.
Côté du carré : \textbf{24 m}. Nombre : $(72\div24)\times(48\div24)=3\times2=\textbf{6 carrés}$.
\end{corrige}

\begin{corrige}[Exercice 2]
PPCM$(12,18,30)$. Décompositions : $12=2^2\times3$, $18=2\times3^2$, $30=2\times3\times5$.
PPCM $=2^2\times3^2\times5=180$ min $=3$ h.
Heure : $6\text{h}00 + 3\text{h} = \textbf{9h00}$.
\end{corrige}

\begin{corrige}[Exercice 3]
Relation : $a\times b = \text{PGCD}\times\text{PPCM}$.
$60 \times b = 12 \times 360 = 4320 \Rightarrow b = 4320 \div 60 = \textbf{72}$.
Vérification : PGCD$(60,72)=12$ ✓ ; PPCM$(60,72)=360$ ✓.
\end{corrige}

\begin{corrige}[Exercice 4]
PGCD$(360,270)$ : $360=270\times1+90$ ; $270=90\times3+0$ $\Rightarrow$ PGCD $=90$ g.
Nombre de parts : $90$.
Masse farine par part : $360\div90=4$ g.
Masse beurre par part : $270\div90=3$ g.
Masse totale d'une part : $7$ g.
\textbf{90 parts de 7 g (4 g farine + 3 g beurre)}.
\end{corrige}

\newpage

\coursec{Activité d'intégration}

\courssub{Objectifs de l'activité}
Cette activité d'intégration vise à \cle{mobiliser} l'ensemble des savoirs et savoir-faire de l'unité d'arithmétique (PGCD par l'algorithme d'Euclide, PPCM, relation $\text{PGCD} \times \text{PPCM} = a \times b$, nombres premiers entre eux) pour résoudre un \cle{problème authentique} en plusieurs étapes. Vous devrez :
\begin{itemize}
    \item \textbf{Traduire} une situation concrète en langage mathématique (identifier les nombres-clés, choisir l'outil adapté : PGCD ou PPCM).
    \item \textbf{Mettre en œuvre} les algorithmes de calcul avec rigueur (divisions euclidiennes successives, décomposition en facteurs premiers si nécessaire, application de la relation fondamentale).
    \item \textbf{Justifier} chaque étape par l'énoncé précis de la propriété ou du théorème utilisé.
    \item \textbf{Interpréter} les résultats numériques dans le contexte du problème (composition des lots, date calendaire).
    \item \textbf{Communiquer} votre démarche par écrit sous forme d'un \cle{mémo professionnel} clair, structuré et argumenté, à destination d'une destinataire non mathématicienne.
\end{itemize}
Cette activité s'inscrit pleinement dans la compétence « \emph{Résoudre des problèmes} » et « \emph{Communiquer} » du programme de 3ème, en préparation aux épreuves écrites du BEPC où la qualité de la rédaction et du raisonnement est notée.

\courssub{Situation : Le grand marché de Mokolo — Partages et rendez-vous}
\emph{Contexte.} Au cœur de Yaoundé, le marché \cle{Mokolo} bat son plein. La \textbf{coopérative « Maman Mokolo »}, regroupant une cinquantaine de revendeuses de céréales et légumineuses, vient de recevoir une livraison groupée pour la saison sèche : \textbf{1\,260 kg de maïs} et \textbf{1\,050 kg de haricots}. La présidente, Madame \textsc{Fouda}, souhaite organiser le stockage et la revente de manière équitable entre les membres actifs.

\medskip
\noindent \textbf{Partie 1 — Constitution des lots identiques (Gestion des stocks).}
Madame \textsc{Fouda} veut constituer le \textbf{plus grand nombre possible de lots identiques} pour les distribuer aux membres. Chaque lot doit contenir \emph{exactement la même masse de maïs} et \emph{exactement la même masse de haricots}. Il ne doit \textbf{rester aucun kilogramme} non réparti (ni maïs, ni haricots). Combien de lots peut-on faire ? Quelle sera la composition exacte (en kg) de chaque lot ?

\medskip
\noindent \textbf{Partie 2 — Synchronisation des approvisionnements (Logistique).}
Trois transporteurs indépendants assurent la liaison entre les zones de production (Nord et Ouest) et le marché Mokolo :
\begin{itemize}
    \item Le camion de \textsc{Papa Ngono} passe \textbf{tous les 10 jours}.
    \item Le camion de \textsc{Maman Biloa} passe \textbf{tous les 15 jours}.
    \item Le camion de \textsc{Tonton Essomba} passe \textbf{tous les 18 jours}.
\end{itemize}
Par un heureux hasard, \textbf{le 1er mars}, les trois camions sont présents simultanément au marché pour décharger. Madame \textsc{Fouda} doit planifier la prochaine réception groupée. \textbf{Quelle sera la date exacte de la prochaine venue simultanée des trois transporteurs ?} (On suppose que les mois de mars, avril, mai ont respectivement 31, 30, 31 jours).

\medskip
\noindent \textbf{Partie 3 — Communication professionnelle.}
Chaque groupe de travail (3 ou 4 élèves) doit rédiger un \textbf{mémo adressé à Madame \textsc{Fouda}}. Ce document, d'une demi-page maximum, présentera : la démarche de calcul du nombre de lots et leur composition, la détermination de la date de prochaine livraison groupée, et une conclusion rassurante sur l'équité du partage (lots identiques, pas de reste). La clarté du langage, la justesse des termes mathématiques (\emph{diviseur, multiple, PGCD, PPCM, premiers entre eux}) et la logique de l'enchaînement seront évaluées.

\courssub{Consignes}
Travaillez par groupe de 3 ou 4. Répondez aux questions dans l'ordre, en montrant \textbf{toutes les étapes de calcul} et en \textbf{écrivant les justifications} (propriétés, théorèmes) pour chaque étape.

\medskip
\noindent \textbf{Partie A : Les lots de maïs et de haricots (PGCD).}
\begin{enumerate}
    \item Quelles grandeurs mathématiques représentent les masses totales de maïs et de haricots ? Quel objet mathématique cherchez-vous pour déterminer le \emph{nombre maximal de lots identiques sans reste} ? Justifiez votre choix en citant la définition du PGCD.
    \item Calculez le PGCD de 1\,260 et 1\,050 \textbf{uniquement par l'algorithme d'Euclide} (divisions euclidiennes successives). Présentez les divisions alignées.
    \item Déduisez le nombre de lots constitués. Calculez la masse de maïs et la masse de haricots dans \emph{un seul lot}.
    \item Vérifiez que les deux masses obtenues pour un lot sont \textbf{premiers entre eux}. Que signifie cette propriété dans le contexte du problème ? (Rappelez la définition de nombres premiers entre eux).
\end{enumerate}

\medskip
\noindent \textbf{Partie B : Le rendez-vous des trois transporteurs (PPCM).}
\begin{enumerate}
    \setcounter{enumi}{4}
    \item Quelle notion mathématique (PGCD ou PPCM) permet de déterminer le délai en jours avant la prochaine rencontre simultanée ? Justifiez en utilisant la notion de \cle{multiple commun}.
    \item Calculez le PPCM de 10, 15 et 18. \textbf{Utilisez la relation $\text{PGCD}(a,b) \times \text{PPCM}(a,b) = a \times b$} en procédant par étapes : d'abord PPCM(10 ; 15), puis PPCM(résultat ; 18). Montrez les calculs de PGCD intermédiaires (par Euclide ou soustractions).
    \item Déterminez la date calendaire exacte (jour et mois) correspondant à ce délai après le 1er mars. Détaillez le décompte jour par jour ou mois par mois.
\end{enumerate}

\medskip
\noindent \textbf{Partie C : Rédaction du mémo (Synthèse).}
\begin{enumerate}
    \setcounter{enumi}{7}
    \item Rédigez le mémo final à l'attention de Madame \textsc{Fouda}. Structurez-le ainsi :
    \begin{itemize}
        \item \textbf{Objet :} Organisation des stocks et planning des livraisons.
        \item \textbf{1. Constitution des lots :} Rappel des quantités, méthode (PGCD/Euclide), résultat (nombre de lots, composition), vérification (premiers entre eux), phrase de conclusion sur l'équité.
        \item \textbf{2. Prochaine livraison groupée :} Périodicités, méthode (PPCM via relation), calcul du délai, conversion en date.
        \item \textbf{3. Conclusion générale.}
    \end{itemize}
\end{enumerate}

\courssub{Ressources à mobiliser}
\begin{aretenir}[Boîte à outils — Arithmétique 3ème]
\textbf{Division euclidienne :} $\forall a \in \mathbb{N}, \forall b \in \mathbb{N}^*, \exists! (q,r) \in \mathbb{N}^2,\ a = bq + r \text{ et } 0 \le r < b$.
\vspace{2mm}

\textbf{PGCD (Définition) :} Le Plus Grand Commun Diviseur de $a$ et $b$ entiers naturels non nuls est le plus grand entier qui divise $a$ et $b$. Noté $\text{PGCD}(a;b)$.
\vspace{2mm}

\textbf{Algorithme d'Euclide :} Pour $a \ge b > 0$, on effectue les divisions successives :
$a = bq_1 + r_1$ ; $b = r_1q_2 + r_2$ ; ... jusqu'au reste nul $r_{k+1}=0$. Alors $\text{PGCD}(a;b) = r_k$ (dernier reste non nul).
\vspace{2mm}

\textbf{Algorithme des soustractions successives :} Pour $a \ge b > 0$, on remplace le couple $(a;b)$ par $(a-b;b)$ (ou $(b;a-b)$) tant que les deux nombres sont distincts. Lorsque les deux nombres sont égaux, leur valeur commune est le PGCD. Propriété : $\text{PGCD}(a;b) = \text{PGCD}(a-b;b)$.
\vspace{2mm}

\textbf{PPCM (Définition) :} Le Plus Petit Commun Multiple de $a$ et $b$ entiers naturels non nuls est le plus petit entier non nul qui est multiple de $a$ et de $b$. Noté $\text{PPCM}(a;b)$.
\vspace{2mm}

\textbf{Relation fondamentale :} $\forall a,b \in \mathbb{N}^*,\ \text{PGCD}(a;b) \times \text{PPCM}(a;b) = a \times b$.
\vspace{2mm}

\textbf{Nombres premiers entre eux :} $a$ et $b$ sont premiers entre eux $\iff \text{PGCD}(a;b) = 1$.
\vspace{2mm}

\textbf{Propriété utile :} Si $d = \text{PGCD}(a;b)$, alors $\frac{a}{d}$ et $\frac{b}{d}$ sont premiers entre eux.
\vspace{2mm}

\textbf{Extension à trois nombres :} $\text{PGCD}(a;b;c) = \text{PGCD}(\text{PGCD}(a;b);c)$ et $\text{PPCM}(a;b;c) = \text{PPCM}(\text{PPCM}(a;b);c)$.
\end{aretenir}

\courssub{Démarche et corrigé commenté}
\begin{exoresolu}[Corrigé détaillé et commenté de l'activité d'intégration]
\textbf{PARTIE A — Constitution des lots (PGCD)}

\noindent \textbf{Question 1. Traduction et choix de l'outil.}
Les masses totales sont $a = 1\,260$ (maïs) et $b = 1\,050$ (haricots), deux entiers naturels non nuls.
On cherche le \textbf{nombre maximal de lots identiques} sans reste. Soit $n$ ce nombre de lots. Chaque lot contient $x$ kg de maïs et $y$ kg de haricots.
On a : $1\,260 = n \times x$ et $1\,050 = n \times y$.
Donc $n$ divise $1\,260$ et $n$ divise $1\,050$ : $n$ est un \cle{diviseur commun} de 1\,260 et 1\,050.
Comme on veut le \emph{maximum} de lots, $n$ doit être le \textbf{plus grand} diviseur commun.
\emph{Justification :} D'après la \textbf{définition du PGCD}, le nombre de lots cherché est $n = \text{PGCD}(1\,260 ; 1\,050)$.

\noindent \textbf{Question 2. Algorithme d'Euclide.}
On applique l'algorithme d'Euclide à $a=1\,260$ et $b=1\,050$ ($a > b$).
\begin{align*}
1\,260 &= 1\,050 \times \mathbf{1} + \mathbf{210} \quad (0 \le 210 < 1\,050) \\
1\,050 &= 210 \times \mathbf{5} + \mathbf{0} \quad (0 \le 0 < 210)
\end{align*}
Le dernier reste non nul est \textbf{210}.
\emph{Conclusion :} $\text{PGCD}(1\,260 ; 1\,050) = \mathbf{210}$.

\noindent \textbf{Question 3. Nombre de lots et composition.}
\begin{itemize}
    \item Nombre maximal de lots : $n = \mathbf{210}$ lots.
    \item Masse de maïs par lot : $x = \frac{1\,260}{210} = \mathbf{6 \text{ kg}}$.
    \item Masse de haricots par lot : $y = \frac{1\,050}{210} = \mathbf{5 \text{ kg}}$.
\end{itemize}
Chaque lot contient donc \textbf{6 kg de maïs et 5 kg de haricots}. Vérification : $210 \times 6 = 1\,260$ ; $210 \times 5 = 1\,050$. Pas de reste.

\noindent \textbf{Question 4. Vérification : nombres premiers entre eux.}
Les masses unitaires sont 6 et 5.
Calculons $\text{PGCD}(6 ; 5)$ :
$6 = 5 \times 1 + 1$ ; $5 = 1 \times 5 + 0 \implies \text{PGCD}(6;5) = 1$.
\emph{Propriété :} Deux entiers sont \textbf{premiers entre eux} si leur PGCD vaut 1.
Ici $\text{PGCD}(6;5)=1$, donc 6 et 5 sont premiers entre eux.
\emph{Sens contextuel :} Cela confirme que le partage est \textbf{irréductible} : on ne peut pas subdiviser davantage les lots en gardant des masses entières identiques pour les deux produits. C'est la garantie mathématique du « nombre maximal de lots ».

\bigskip
\noindent \textbf{PARTIE B — Rendez-vous des transporteurs (PPCM)}

\noindent \textbf{Question 5. Choix de l'outil : PPCM.}
Les camions passent tous les 10, 15 et 18 jours. Les jours de passage de chaque camion sont les \textbf{multiples} de sa période.
Le camion A passe aux jours : $10, 20, 30, \dots$ (multiples de 10).
Le camion B passe aux jours : $15, 30, 45, \dots$ (multiples de 15).
Le camion C passe aux jours : $18, 36, 54, \dots$ (multiples de 18).
Une présence \textbf{simultanée} correspond à un jour qui est multiple de 10 \textbf{ET} de 15 \textbf{ET} de 18, c'est-à-dire un \cle{multiple commun} à 10, 15 et 18.
Le 1er mars correspond au jour 0 (ou jour de départ). La \textbf{prochaine} rencontre correspond au \textbf{plus petit} multiple commun \textbf{strictement positif}.
\emph{Justification :} D'après la \textbf{définition du PPCM}, le délai en jours cherché est $D = \text{PPCM}(10 ; 15 ; 18)$.

\noindent \textbf{Question 6. Calcul du PPCM par la relation fondamentale (étapes).}
On calcule par étapes : $m_1 = \text{PPCM}(10;15)$, puis $D = \text{PPCM}(m_1;18)$.
Relation : $\text{PPCM}(a;b) = \frac{a \times b}{\text{PGCD}(a;b)}$.

\emph{Étape 1 : PPCM(10 ; 15).}
Calcul de $\text{PGCD}(10;15)$ par Euclide :
$15 = 10 \times 1 + 5$ ; $10 = 5 \times 2 + 0 \implies \text{PGCD}(10;15) = 5$.
Application de la relation :
$\text{PPCM}(10;15) = \frac{10 \times 15}{5} = \frac{150}{5} = \mathbf{30}$.
On note $m_1 = 30$.

\emph{Étape 2 : PPCM(30 ; 18).}
Calcul de $\text{PGCD}(30;18)$ par Euclide :
$30 = 18 \times 1 + 12$ ; $18 = 12 \times 1 + 6$ ; $12 = 6 \times 2 + 0 \implies \text{PGCD}(30;18) = 6$.
Application de la relation :
$\text{PPCM}(30;18) = \frac{30 \times 18}{6} = \frac{540}{6} = \mathbf{90}$.
\emph{Conclusion :} Le délai avant la prochaine rencontre simultanée est \textbf{90 jours}.

\noindent \textbf{Question 7. Conversion en date calendaire.}
Départ : 1er mars (jour 0). On ajoute 90 jours.
\begin{itemize}
    \item \textbf{Mars :} 31 jours au total. Jours restants en mars après le 1er mars : 30 jours (du 2 mars au 31 mars). Reste $90 - 30 = 60$ jours.
    \item \textbf{Avril :} 30 jours. Reste $60 - 30 = 30$ jours.
    \item \textbf{Mai :} 31 jours. Les 30 jours restants mènent au \textbf{30 mai}.
\end{itemize}
\emph{Vérification :} 30 (mars) + 30 (avril) + 30 (mai) = 90 jours après le 1er mars.
\emph{Réponse :} La prochaine livraison groupée aura lieu le \textbf{30 mai}.

\bigskip
\noindent \textbf{PARTIE C — Mémo type (Proposition de rédaction)}

\begin{center}
\begin{tabularx}{\linewidth}{|X|}\hline
\rowcolor{popblueL}
\textbf{COOPÉRATIVE « MAMAN MOKOLO » — MEMO INTERNE} \\ \hline
\textbf{À :} Madame la Présidente, \textsc{Fouda} \\
\textbf{De :} Commission Gestion des Stocks \& Logistique \\
\textbf{Date :} 1er mars 202X \\
\textbf{Objet :} Organisation des lots (maïs/haricots) et planning livraisons groupées \\ \hline
\end{tabularx}
\end{center}

\vspace{2mm}
\noindent \textbf{1. Constitution des lots équitables}
La livraison reçue porte sur 1\,260 kg de maïs et 1\,050 kg de haricots. Pour constituer le maximum de lots identiques sans reste, nous avons déterminé le PGCD de ces deux masses par l'algorithme d'Euclide :
$1\,260 = 1\,050 \times 1 + 210$ ; $1\,050 = 210 \times 5 + 0 \Rightarrow \text{PGCD} = 210$.
\textbf{Résultat :} \textbf{210 lots} pourront être constitués.
Chaque lot contiendra : $\frac{1\,260}{210} = \mathbf{6 \text{ kg de maïs}}$ et $\frac{1\,050}{210} = \mathbf{5 \text{ kg de haricots}}$.
Nous avons vérifié que 6 et 5 sont premiers entre eux (PGCD = 1), ce qui garantit que ce partage est le plus fin possible : aucune subdivision supplémentaire en lots identiques n'est possible. L'équité entre les 210 membres bénéficiaires est donc mathématiquement assurée.

\vspace{2mm}
\noindent \textbf{2. Prochaine livraison groupée des trois transporteurs}
Les périodes sont de 10, 15 et 18 jours. La prochaine coïncidence correspond au PPCM de ces trois nombres.
Calcul par la relation $\text{PGCD} \times \text{PPCM} = a \times b$ :
\begin{itemize}
    \item $\text{PGCD}(10;15)=5 \Rightarrow \text{PPCM}(10;15)=30$.
    \item $\text{PGCD}(30;18)=6 \Rightarrow \text{PPCM}(30;18)=90$.
\end{itemize}
Le délai est de \textbf{90 jours} après le 1er mars.
Décompte : 30 jours (reste mars) + 30 jours (avril) + 30 jours (mai) = 90 jours.
\textbf{Date prévue :} \textbf{30 mai}.

\vspace{2mm}
\noindent \textbf{3. Conclusion}
Le stock est prêt pour une distribution immédiate en 210 lots équitables (6 kg maïs + 5 kg haricots). La prochaine réception groupée des trois fournisseurs est programmée pour le \textbf{30 mai}. Nous restons à disposition pour la mise en place physique des lots.
\vspace{1mm}
\noindent \textit{La Commission (Groupe 3\`eme A)}
\end{exoresolu}

\courssub{Bilan de l'activité}
Cette activité a permis de mettre en évidence plusieurs points clés du programme de 3ème :
\begin{itemize}
    \item \textbf{Modélisation :} La traduction du langage courant (« lots identiques, maximum, sans reste » $\to$ PGCD ; « rendez-vous simultané, périodicités » $\to$ PPCM) est l'étape cruciale. L'élève doit justifier son choix par la définition (diviseur commun vs multiple commun).
    \item \textbf{Maîtrise algorithmique :} L'algorithme d'Euclide s'est révélé plus rapide et moins source d'erreurs que la méthode des soustractions ou la décomposition en facteurs premiers pour des nombres de cette taille (1\,260 ; 1\,050). La présentation alignée des divisions euclidiennes est la norme attendue au BEPC.
    \item \textbf{Articulation PGCD/PPCM :} L'utilisation de la relation $\text{PGCD} \times \text{PPCM} = a \times b$ pour calculer le PPCM de trois nombres (par itération sur deux nombres) est une technique puissante qu'il faut savoir expliciter : « Je calcule le PGCD du couple, j'en déduis le PPCM, je recombine avec le troisième nombre ».
    \item \textbf{Sens des résultats :} Vérifier que les quotients $\frac{a}{\text{PGCD}}$ et $\frac{b}{\text{PGCD}}$ sont premiers entre eux n'est pas un simple exercice de style : c'est la \textbf{preuve} que le nombre de lots est bien maximal (irréductibilité du partage).
    \item \textbf{Compétence communication :} Le mémo final exige de passer du langage mathématique (symboles, égalités) au langage courant précis, sans perdre la rigueur (citer « PGCD », « algorithme d'Euclide », « PPCM », « premiers entre eux »). C'est un entraînement direct à l'\emph{épreuve de résolution de problèmes} du BEPC où la rédaction compte pour une part significative de la note.
    \item \textbf{Contexte camerounais :} L'ancrage dans la réalité du marché Mokolo (masses en kg, périodicités en jours, calendrier grégorien) donne du sens aux nombres et motive la recherche de la solution exacte (pas d'arrondi, date précise).
\end{itemize}
\medskip
\noindent \textbf{Pistes de remédiation / Différenciation :}
\begin{itemize}
    \item \emph{En difficulté :} Fournir une fiche aide-mémoire avec la structure de l'algorithme d'Euclide (tableau à trous) et la formule de la relation PGCD$\times$PPCM. Proposer des nombres plus petits pour l'entraînement (ex: 126 et 105).
    \item \emph{Avancés :} Demander de résoudre la Partie B par décomposition en facteurs premiers (arbre de factorisation) et de comparer les deux méthodes. Poser la question inverse : « Si la prochaine rencontre est le 30 mai, quelles pourraient être les périodes des camions ? » (Recherche de triplets dont le PPCM est 90).
\end{itemize}
Cette activité clôture l'unité d'arithmétique en validant la capacité de l'élève à \textbf{choisir, appliquer, justifier et communiquer} — le cœur du socle commun de compétences visé par le MINESEC.

\newpage

\coursec{Méthodes et exercices résolus}

\courssub{Méthodes et exercices résolus}

\begin{methode}[Rappel : division euclidienne et vocabulaire de la divisibilité]
\textbf{Objectif :} Maîtriser le vocabulaire (diviseur, multiple, dividende, quotient, reste) et poser une division euclidienne.
\begin{enumerate}
    \item \textbf{Identifier les termes :} pour une division euclidienne $a = b \times q + r$ avec $0 \le r < b$ :
    \begin{itemize}
        \item $a$ est le \cle{dividende} ;
        \item $b$ est le \cle{diviseur} ;
        \item $q$ est le \cle{quotient} ;
        \item $r$ est le \cle{reste}.
    \end{itemize}
    \item \textbf{Tester la divisibilité :} $b$ divise $a$ (on note $b \mid a$) si et seulement si le reste $r = 0$. On dit alors aussi que $a$ est un \cle{multiple} de $b$.
    \item \textbf{Poser la division :} utiliser la technique de la division posée ou la calculatrice pour trouver $q$ et $r$.
    \item \textbf{Vérifier :} recalculer $b \times q + r$ ; on doit retrouver $a$, et contrôler que $0 \le r < b$.
\end{enumerate}
\textbf{Réflexe :} si $r=0$, dire « $b$ divise $a$ » ou « $a$ est un multiple de $b$ ». Ne pas confondre « diviseur » (celui qui divise) et « dividende » (celui qui est divisé).
\end{methode}

\begin{exoresolu}[Division euclidienne dans un contexte camerounais]
Un transporteur de \textbf{Bamenda} doit charger \num{1358} sacs de café dans des camions. Chaque camion peut contenir exactement \num{42} sacs.
\begin{enumerate}
    \item Effectuer la division euclidienne de \num{1358} par \num{42}.
    \item Combien de camions pleins peut-on charger ?
    \item Combien de sacs reste-t-il à charger dans un dernier camion non plein ?
    \item Écrire l'égalité de la division euclidienne.
\end{enumerate}
\tcblower
\textbf{Solution détaillée :}
\begin{enumerate}
    \item \textbf{Recherche du quotient et du reste :} on cherche combien de fois 42 « rentre » dans 1358.
    $42 \times 30 = 1260$ et $1358 - 1260 = 98$ ; puis $42 \times 2 = 84$ et $98 - 84 = 14$.
    Donc $1358 = 42 \times (30 + 2) + 14$, c'est-à-dire :
    \[ 1358 = 42 \times 32 + 14 \quad \text{avec } 0 \le 14 < 42. \]
    Le quotient est $q = 32$ et le reste est $r = 14$.
    \item Le nombre de camions \textbf{pleins} est le quotient : \textbf{32 camions}.
    \item Le nombre de sacs restants est le reste : \textbf{14 sacs}.
    \item L'égalité de la division euclidienne est :
    \[ 1358 = 42 \times 32 + 14. \]
    \textbf{Vérification :} $42 \times 32 = 1344$ ; $1344 + 14 = 1358$. C'est exact.
\end{enumerate}
\textbf{Conclusion :} il faut 32 camions pleins et un 33\textsuperscript{e} camion pour les 14 sacs restants.
\end{exoresolu}

\begin{methode}[PGCD par la méthode des listes de diviseurs (petits nombres)]
\textbf{Objectif :} trouver le Plus Grand Commun Diviseur (PGCD) de deux entiers naturels en listant leurs diviseurs.
\begin{enumerate}
    \item \textbf{Lister les diviseurs} de chaque nombre (on peut s'arrêter à la racine carrée et compléter par paires pour ne rien oublier).
    \item \textbf{Repérer les diviseurs communs} (présents dans les deux listes).
    \item \textbf{Choisir le plus grand} de ces diviseurs communs : c'est le PGCD.
    \item \textbf{Notation :} $\text{PGCD}(a; b)$.
\end{enumerate}
\textbf{Limite :} cette méthode devient longue et risquée pour des nombres supérieurs à 200 ou 300. On utilisera alors l'algorithme des soustractions ou celui d'Euclide.
\end{methode}

\begin{exoresolu}[PGCD par listes pour un découpage de terrain]
Un agriculteur de \textbf{Dschang} possède deux parcelles rectangulaires de dimensions \num{48} m sur \num{60} m et \num{36} m sur \num{72} m. Il veut les diviser en carrés identiques de côté entier le plus grand possible, sans perte de terrain.
\begin{enumerate}
    \item Déterminer le PGCD de 48 et 60 par la méthode des listes.
    \item Déterminer le PGCD de 36 et 72 par la méthode des listes.
    \item Quelle est la longueur du côté des carrés pour chaque parcelle ?
\end{enumerate}
\tcblower
\textbf{Solution détaillée :}
\begin{enumerate}
    \item \textbf{Diviseurs de 48 :} $1, 2, 3, 4, 6, 8, 12, 16, 24, 48$.\\
    \textbf{Diviseurs de 60 :} $1, 2, 3, 4, 5, 6, 10, 12, 15, 20, 30, 60$.\\
    \textbf{Diviseurs communs :} $1, 2, 3, 4, 6, 12$.\\
    Le plus grand est \textbf{12}. Donc $\text{PGCD}(48; 60) = 12$.
    \item \textbf{Diviseurs de 36 :} $1, 2, 3, 4, 6, 9, 12, 18, 36$.\\
    \textbf{Diviseurs de 72 :} $1, 2, 3, 4, 6, 8, 9, 12, 18, 24, 36, 72$.\\
    \textbf{Diviseurs communs :} $1, 2, 3, 4, 6, 9, 12, 18, 36$.\\
    Le plus grand est \textbf{36}. Donc $\text{PGCD}(36; 72) = 36$ (ce qui était prévisible : 36 divise 72, car $72 = 36 \times 2$).
    \item Pour la 1\textsuperscript{re} parcelle : côté du carré = \textbf{12 m} (on obtient $48 \div 12 = 4$ carrés en largeur et $60 \div 12 = 5$ en longueur, soit 20 carrés).\\
    Pour la 2\textsuperscript{e} parcelle : côté du carré = \textbf{36 m} (soit $1 \times 2 = 2$ carrés).
\end{enumerate}
\textbf{Conclusion :} les carrés feront 12 m de côté pour la première parcelle et 36 m pour la seconde.
\end{exoresolu}

\begin{methode}[PGCD par l'algorithme des soustractions successives]
\textbf{Objectif :} trouver le PGCD de deux entiers $a$ et $b$ ($a \ge b$) par soustractions répétées.
\begin{enumerate}
    \item \textbf{Ordonner :} soit $a$ le plus grand, $b$ le plus petit.
    \item \textbf{Soustraire :} remplacer $a$ par $a - b$. On considère le nouveau couple $(a-b \,; b)$.
    \item \textbf{Réordonner :} remettre le plus grand des deux nombres en première position.
    \item \textbf{Répéter} les étapes 2 et 3 jusqu'à obtenir deux nombres \textbf{égaux}.
    \item \textbf{Résultat :} ce nombre commun est le PGCD.
\end{enumerate}
\textbf{Justification :} tout diviseur commun à $a$ et $b$ divise aussi $a - b$, et réciproquement ; donc $\text{PGCD}(a; b) = \text{PGCD}(a-b; b)$ : les diviseurs communs ne changent pas à chaque étape.
\textbf{Présentation :} on peut disposer les calculs dans un tableau à deux colonnes.
\end{methode}

\begin{exoresolu}[Algorithme des soustractions : câblage réseau]
Un technicien \textbf{CAMTEL} doit couper deux câbles de longueurs \num{216} cm et \num{156} cm en morceaux de même longueur entière maximale, sans perte. On utilise l'algorithme des soustractions.
\begin{enumerate}
    \item Appliquer l'algorithme des soustractions pour trouver le PGCD de 216 et 156.
    \item Quelle sera la longueur de chaque morceau ?
    \item Combien de morceaux obtient-on au total ?
\end{enumerate}
\tcblower
\textbf{Solution détaillée :}
\begin{enumerate}
    \item \textbf{Tableau des soustractions successives :}
    \begin{center}
    \begin{tabular}{|c|c|}\hline
    \textbf{Nombre 1} & \textbf{Nombre 2} \\ \hline
    216 & 156 \\ \hline
    60 ($=216-156$) & 156 \\ \hline
    156 & 60 \\ \hline
    96 ($=156-60$) & 60 \\ \hline
    60 & 96 \\ \hline
    60 & 36 ($=96-60$) \\ \hline
    24 ($=60-36$) & 36 \\ \hline
    36 & 24 \\ \hline
    12 ($=36-24$) & 24 \\ \hline
    24 & 12 \\ \hline
    \textbf{12} & \textbf{12} \\ \hline
    \end{tabular}
    \end{center}
    Les deux nombres deviennent égaux à \textbf{12}.
    Donc $\text{PGCD}(216; 156) = 12$.
    \item La longueur maximale de chaque morceau est \textbf{12 cm}.
    \item Nombre de morceaux :
    \begin{itemize}
        \item câble 1 : $216 \div 12 = 18$ morceaux ;
        \item câble 2 : $156 \div 12 = 13$ morceaux.
    \end{itemize}
    Total : $18 + 13 = \textbf{31 morceaux}$.
\end{enumerate}
\textbf{Conclusion :} on coupe en morceaux de 12 cm et l'on obtient 31 morceaux au total.
\end{exoresolu}

\begin{methode}[PGCD par l'algorithme d'Euclide — méthode de référence]
\textbf{Objectif :} trouver le PGCD rapidement, même pour de grands nombres, grâce aux divisions euclidiennes successives.
\begin{enumerate}
    \item \textbf{Initialiser :} soit $a$ et $b$ les deux entiers ($a \ge b > 0$).
    \item \textbf{Diviser :} effectuer la division euclidienne de $a$ par $b$ : $a = b \times q_1 + r_1$ avec $0 \le r_1 < b$.
    \item \textbf{Tester le reste :}
    \begin{itemize}
        \item si $r_1 = 0$ : \textbf{STOP}, le PGCD est $b$ ;
        \item si $r_1 \neq 0$ : recommencer l'étape 2 avec le couple $(b \,; r_1)$.
    \end{itemize}
    \item \textbf{Boucler :} tant que le reste n'est pas nul, on divise le diviseur précédent par le reste obtenu.
    \item \textbf{Écrire la chaîne de divisions} (rédaction exigée au BEPC) :
    \[
    \begin{cases}
    a = b \times q_1 + r_1 \\
    b = r_1 \times q_2 + r_2 \\
    r_1 = r_2 \times q_3 + r_3 \\
    \vdots \\
    r_{k-1} = r_k \times q_{k+1} + 0
    \end{cases}
    \implies \text{PGCD}(a; b) = r_k
    \]
\end{enumerate}
\textbf{Réflexe BEPC :} toujours écrire la suite des égalités de divisions euclidiennes. Le PGCD est le \textbf{dernier reste non nul}.
\textbf{Pourquoi ça marche :} comme pour les soustractions, les diviseurs communs à $a$ et $b$ sont exactement les diviseurs communs à $b$ et $r_1$ (car $r_1 = a - b q_1$). On réduit ainsi le problème jusqu'au reste nul.
\end{methode}

\begin{exoresolu}[Algorithme d'Euclide : numéros d'abonnés]
Deux numéros d'identification d'abonnés sont $a = 1428$ et $b = 1071$. On veut configurer un groupe fermé d'utilisateurs dont la taille maximale est le PGCD de ces deux numéros.
\begin{enumerate}
    \item Déterminer $\text{PGCD}(1428; 1071)$ par l'algorithme d'Euclide en détaillant toutes les divisions.
    \item En déduire la taille maximale du groupe.
\end{enumerate}
\tcblower
\textbf{Solution détaillée :}
\begin{enumerate}
    \item \textbf{Chaîne des divisions euclidiennes :}
    \begin{align*}
    1428 &= 1071 \times \mathbf{1} + \mathbf{357} \quad (0 \le 357 < 1071) \\
    1071 &= 357 \times \mathbf{3} + \mathbf{0}
    \end{align*}
    Le dernier reste non nul est \textbf{357}.
    Donc $\text{PGCD}(1428; 1071) = 357$.
    \item La taille maximale du groupe est de \textbf{357 abonnés}.
\end{enumerate}
\textbf{Vérification rapide :} $1428 \div 357 = 4$ (entier) et $1071 \div 357 = 3$ (entier). C'est correct.
\textbf{Conclusion :} le PGCD est 357 ; le groupe pourra contenir au maximum 357 membres.
\end{exoresolu}

\begin{methode}[Nombres premiers entre eux et fractions irréductibles]
\textbf{Objectif :} reconnaître des nombres premiers entre eux et simplifier une fraction au maximum.
\begin{enumerate}
    \item \textbf{Définition :} deux entiers $a$ et $b$ sont \cle{premiers entre eux} si $\text{PGCD}(a; b) = 1$, c'est-à-dire si leur seul diviseur commun est 1.
    \item \textbf{Application courante :} pour simplifier une fraction $\dfrac{a}{b}$, on divise le numérateur et le dénominateur par leur PGCD. La fraction obtenue est \cle{irréductible} : son numérateur et son dénominateur sont premiers entre eux.
    \item \textbf{Vérification :} après simplification, on peut refaire l'algorithme d'Euclide sur le nouveau couple pour contrôler que le PGCD vaut bien 1.
\end{enumerate}
\textbf{Attention :} « premiers entre eux » ne veut pas dire « nombres premiers ». Par exemple, 8 et 15 sont premiers entre eux ($\text{PGCD}(8;15)=1$) mais ni 8 ni 15 n'est un nombre premier.
\end{methode}

\begin{exoresolu}[Simplification de fraction et nombres premiers entre eux]
On considère la fraction $\dfrac{561}{391}$ issue d'un calcul de débit moyen sur une liaison \textbf{MTN Fibre}.
\begin{enumerate}
    \item Montrer que 561 et 391 ne sont pas premiers entre eux en calculant leur PGCD par l'algorithme d'Euclide.
    \item Simplifier la fraction pour la rendre irréductible.
    \item Vérifier que le numérateur et le dénominateur de la fraction simplifiée sont premiers entre eux.
\end{enumerate}
\tcblower
\textbf{Solution détaillée :}
\begin{enumerate}
    \item \textbf{Algorithme d'Euclide :}
    \begin{align*}
    561 &= 391 \times 1 + 170 \\
    391 &= 170 \times 2 + 51 \\
    170 &= 51 \times 3 + 17 \\
    51 &= 17 \times 3 + 0
    \end{align*}
    Le dernier reste non nul est \textbf{17}.
    $\text{PGCD}(561; 391) = 17 \neq 1$ : les nombres 561 et 391 ne sont \textbf{pas} premiers entre eux.
    \item \textbf{Simplification :} on divise le numérateur et le dénominateur par 17.
    \[ \frac{561}{391} = \frac{561 \div 17}{391 \div 17} = \frac{33}{23}. \]
    \item \textbf{Vérification :} calculons $\text{PGCD}(33; 23)$.
    \[ 33 = 23 \times 1 + 10 \ ;\quad 23 = 10 \times 2 + 3 \ ;\quad 10 = 3 \times 3 + 1 \ ;\quad 3 = 1 \times 3 + 0. \]
    Le dernier reste non nul est 1, donc $\text{PGCD}(33; 23) = 1$ : les nombres 33 et 23 sont \textbf{premiers entre eux} et la fraction $\dfrac{33}{23}$ est \textbf{irréductible}.
\end{enumerate}
\textbf{Conclusion :} la fraction irréductible équivalente est $\dfrac{33}{23}$.
\end{exoresolu}

\begin{methode}[PPCM et relation fondamentale $\text{PGCD} \times \text{PPCM} = a \times b$]
\textbf{Objectif :} calculer le Plus Petit Commun Multiple (PPCM) et utiliser sa relation avec le PGCD.
\begin{enumerate}
    \item \textbf{Définition :} le PPCM de deux entiers naturels non nuls $a$ et $b$ (noté $\text{PPCM}(a; b)$) est le plus petit entier naturel \textbf{non nul} qui est multiple à la fois de $a$ et de $b$.
    \item \textbf{Méthode 1 (listes) :} lister les multiples de chaque nombre jusqu'à trouver le premier multiple commun (réservée aux petits nombres).
    \item \textbf{Méthode 2 (relation fondamentale) :} pour tous entiers naturels non nuls $a$ et $b$ :
    \[ \text{PGCD}(a; b) \times \text{PPCM}(a; b) = a \times b \qquad \text{donc} \qquad \text{PPCM}(a; b) = \dfrac{a \times b}{\text{PGCD}(a; b)}. \]
    \item \textbf{Étapes avec la formule :}
    \begin{enumerate}
        \item calculer $\text{PGCD}(a; b)$ (algorithme d'Euclide) ;
        \item calculer le produit $a \times b$ ;
        \item diviser ce produit par le PGCD.
    \end{enumerate}
    \item \textbf{Cas particulier :} si $a$ et $b$ sont premiers entre eux, alors $\text{PGCD}(a;b)=1$, donc $\text{PPCM}(a; b) = a \times b$.
\end{enumerate}
\textbf{Réflexe :} toujours calculer le PGCD d'abord. Vérifier que la division du produit par le PGCD tombe juste (reste 0).
\end{methode}

\begin{exoresolu}[PPCM : synchronisation de feux tricolores à Yaoundé]
Deux feux tricolores au carrefour \textbf{Échangeur} passent au vert simultanément à 08h00.
Le feu A a un cycle de \num{120} secondes. Le feu B a un cycle de \num{168} secondes.
\begin{enumerate}
    \item Calculer le PGCD de 120 et 168.
    \item En déduire le PPCM de 120 et 168.
    \item À quelle heure les deux feux repasseront-ils au vert en même temps pour la première fois après 08h00 ?
\end{enumerate}
\tcblower
\textbf{Solution détaillée :}
\begin{enumerate}
    \item \textbf{PGCD(168 ; 120) par l'algorithme d'Euclide :}
    \begin{align*}
    168 &= 120 \times 1 + 48 \\
    120 &= 48 \times 2 + 24 \\
    48 &= 24 \times 2 + 0
    \end{align*}
    Le dernier reste non nul est 24, donc $\text{PGCD}(120; 168) = \mathbf{24}$.
    \item \textbf{PPCM par la relation fondamentale :}
    \[ \text{PPCM}(120; 168) = \frac{120 \times 168}{\text{PGCD}(120; 168)} = \frac{20160}{24} = 840. \]
    (Vérification : $24 \times 840 = 20160$.) Le PPCM est \textbf{840 secondes}.
    \item \textbf{Conversion en minutes :} $840 \text{ s} = \dfrac{840}{60} \text{ min} = \mathbf{14 \text{ minutes}}$.\\
    Heure de synchronisation : 08h00 + 14 min = \textbf{08h14}.
\end{enumerate}
\textbf{Conclusion :} les deux feux repasseront au vert ensemble pour la première fois à \textbf{08h14}.
\end{exoresolu}

\begin{methode}[Résolution de problèmes concrets : « PGCD ou PPCM ? »]
\textbf{Objectif :} choisir la bonne notion (PGCD ou PPCM) selon l'énoncé.
\begin{enumerate}
    \item \textbf{Analyser le verbe ou l'action :}
    \begin{itemize}
        \item \textbf{Partager, diviser, couper, regrouper en paquets identiques, maximiser la taille d'un groupe} $\rightarrow$ \cle{PGCD} (on cherche un diviseur commun le \textbf{plus grand} possible).
        \item \textbf{Synchroniser, répéter un cycle, se retrouver ensemble, s'aligner, trouver une quantité multiple des deux} $\rightarrow$ \cle{PPCM} (on cherche un multiple commun le \textbf{plus petit} possible).
    \end{itemize}
    \item \textbf{Poser les données :} identifier les deux (ou plusieurs) nombres entiers $a$ et $b$.
    \item \textbf{Calculer :} appliquer l'algorithme d'Euclide pour le PGCD, puis la relation fondamentale pour le PPCM si nécessaire.
    \item \textbf{Répondre par une phrase} en reprenant les unités du problème (cm, secondes, sacs, personnes...).
\end{enumerate}
\textbf{Astuce :} si le problème demande « la plus grande longueur possible » ou « le plus grand nombre de parts identiques » $\rightarrow$ PGCD ; s'il demande « le plus petit temps » ou « la plus petite quantité commune » $\rightarrow$ PPCM.
\end{methode}

\begin{exoresolu}[Problème de synthèse : carrelage et rendez-vous]
\textbf{Partie A : carrelage.} On veut carreler une salle rectangulaire de \num{660} cm sur \num{450} cm avec des carreaux carrés identiques de côté entier (en cm), sans les couper.
\begin{enumerate}
    \item Quelle est la dimension du plus grand carreau possible ?
    \item Combien de carreaux faut-il ?
\end{enumerate}
\textbf{Partie B : rendez-vous.} Deux bus partent du même terminus à 6h00. Le bus de la ligne 1 revient au terminus toutes les \num{45} minutes. Le bus de la ligne 2 revient toutes les \num{72} minutes.
\begin{enumerate}
    \item À quelle heure se retrouveront-ils ensemble au terminus pour la première fois après 6h00 ?
    \item Combien de tours aura fait chaque bus à ce moment-là ?
\end{enumerate}
\tcblower
\textbf{Solution détaillée :}

\textbf{Partie A (PGCD : découpage, partage) :}
\begin{enumerate}
    \item Le côté du carreau doit diviser à la fois 660 et 450 ; on cherche le plus grand possible : c'est $\text{PGCD}(660; 450)$.
    \begin{align*}
    660 &= 450 \times 1 + 210 \\
    450 &= 210 \times 2 + 30 \\
    210 &= 30 \times 7 + 0
    \end{align*}
    Le dernier reste non nul est 30, donc $\text{PGCD}(660; 450) = \mathbf{30}$. Le côté du plus grand carreau est \textbf{30 cm}.
    \item \textbf{Nombre de carreaux :} en longueur : $660 \div 30 = 22$ carreaux ; en largeur : $450 \div 30 = 15$ carreaux.
    \[ \text{Nombre total} = 22 \times 15 = \mathbf{330 \text{ carreaux}}. \]
    (Vérification par les aires : $\dfrac{660 \times 450}{30 \times 30} = \dfrac{297000}{900} = 330$.)
\end{enumerate}
\textbf{Partie B (PPCM : synchronisation de cycles) :}
\begin{enumerate}
    \item On cherche le plus petit temps multiple à la fois de 45 et de 72 : c'est $\text{PPCM}(45; 72)$. On calcule d'abord le PGCD :
    \begin{align*}
    72 &= 45 \times 1 + 27 \\
    45 &= 27 \times 1 + 18 \\
    27 &= 18 \times 1 + 9 \\
    18 &= 9 \times 2 + 0
    \end{align*}
    Donc $\text{PGCD}(45; 72) = 9$, puis :
    \[ \text{PPCM}(45; 72) = \frac{45 \times 72}{9} = 45 \times 8 = \mathbf{360 \text{ minutes}} = 6 \text{ heures}. \]
    Heure des retrouvailles : 6h00 + 6h = \textbf{12h00 (midi)}.
    \item \textbf{Nombre de tours :}
    bus 1 : $360 \div 45 = \mathbf{8}$ tours ; bus 2 : $360 \div 72 = \mathbf{5}$ tours.
\end{enumerate}
\textbf{Conclusion :} il faut des carreaux de 30 cm de côté (330 carreaux). Les deux bus se retrouvent au terminus à 12h00, après 8 tours pour la ligne 1 et 5 tours pour la ligne 2.
\end{exoresolu}

\begin{attention}[Les 5 erreurs qui coûtent des points au BEPC]
\begin{enumerate}
    \item \textbf{Confondre PGCD et PPCM dans l'énoncé :} écrire « PGCD » quand il faut « PPCM » (ou l'inverse) fausse tout l'exercice. \textbf{Relisez le verbe :} « couper, diviser, partager » = PGCD ; « se retrouver, synchroniser, multiple commun » = PPCM.
    \item \textbf{Oublier d'écrire la chaîne des divisions euclidiennes :} au BEPC, le résultat seul (même juste) ne suffit pas. Il faut \textbf{impérativement} écrire les égalités $a = bq + r$, $b = rq' + r'$, ... jusqu'au reste nul. Sans cette rédaction, la note est plafonnée.
    \item \textbf{Ne pas vérifier la condition $0 \le r < b$ :} si vous trouvez un reste négatif ou supérieur ou égal au diviseur, votre division est fausse. Contrôlez cette inégalité à chaque ligne.
    \item \textbf{Mal appliquer la relation $\text{PGCD} \times \text{PPCM} = a \times b$ :} l'utiliser sans avoir calculé le PGCD d'abord, ou calculer $a \times b$ et oublier de diviser par le PGCD. \textbf{Ordre obligatoire :} 1. PGCD ; 2. produit $a \times b$ ; 3. division.
    \item \textbf{Réponse sans phrase ni unité :} « 12 » seul ne vaut presque rien. Il faut écrire : « La longueur maximale du côté du carré est \textbf{12 cm} » ou « Ils se retrouvent à \textbf{08h14} ». L'unité et la phrase de conclusion font partie de la note de rédaction.
\end{enumerate}
\end{attention}

\newpage

\coursec{Exercices}

\courssub{Je vérifie mes connaissances}

\begin{exercice}[QCM : Diviseurs et PGCD]
Parmi les affirmations suivantes, une seule est exacte. Laquelle ? Justifier brièvement.
\begin{enumerate}
    \item Les diviseurs de $18$ sont $1 ; 2 ; 3 ; 6 ; 9$.
    \item $\text{PGCD}(24 ; 36) = 6$.
    \item $15$ et $28$ sont premiers entre eux.
    \item Si $a$ est un multiple de $b$, alors $\text{PGCD}(a ; b) = a$.
\end{enumerate}
\end{exercice}

\begin{exercice}[Vrai ou Faux ? Justifier]
Dire si les affirmations suivantes sont vraies ou fausses. Justifier chaque réponse par un contre-exemple ou une propriété du cours.
\begin{enumerate}
    \item « Si deux nombres sont premiers entre eux, alors leur PPCM est égal à leur produit. »
    \item « Deux nombres pairs sont toujours premiers entre eux. »
    \item « $\text{PGCD}(a ; b) = \text{PGCD}(b ; r)$ où $r$ est le reste de la division euclidienne de $a$ par $b$. »
    \item « Le PPCM de deux entiers est toujours plus grand que leur PGCD. »
\end{enumerate}
\end{exercice}

\begin{exercice}[Définitions et notations]
Compléter les phrases suivantes en utilisant le vocabulaire mathématique précis.
\begin{enumerate}
    \item On appelle \ldots d'un entier naturel $a$ tout entier naturel $d$ tel que le reste de la division de $a$ par $d$ soit nul.
    \item Le \ldots de deux entiers naturels non nuls $a$ et $b$ est le plus grand de leurs diviseurs communs. On le note \ldots
    \item Deux entiers naturels non nuls $a$ et $b$ sont dits \ldots si leur PGCD est égal à $1$.
    \item Le \ldots de deux entiers naturels non nuls $a$ et $b$ est le plus petit de leurs multiples communs non nuls. On le note \ldots
\end{enumerate}
\end{exercice}

\begin{exercice}[Calculs directs : PGCD et PPCM]
Calculer les PGCD et PPCM des couples d'entiers suivants. On présentera les divisions euclidiennes pour le PGCD (algorithme d'Euclide).
\begin{enumerate}
    \item $a = 144$ et $b = 60$
    \item $a = 231$ et $b = 154$
    \item $a = 377$ et $b = 221$
    \item $a = 48$ et $b = 180$ (calculer PGCD et PPCM, vérifier la relation $\text{PGCD} \times \text{PPCM} = a \times b$)
\end{enumerate}
\end{exercice}

\courssub{Je m'entraîne}

\begin{exercice}[Algorithme des soustractions]
On veut calculer $\text{PGCD}(437 ; 323)$ par la méthode des soustractions successives.
\begin{enumerate}
    \item Rappeler la propriété qui justifie cette méthode : $\text{PGCD}(a ; b) = \text{PGCD}(a-b ; b)$ pour $a > b$.
    \item Appliquer pas à pas la méthode des soustractions jusqu'à obtenir deux nombres égaux.
    \item Conclure sur la valeur du PGCD.
\end{enumerate}
\end{exercice}

\begin{exercice}[Fraction irréductible]
On considère la fraction $\dfrac{1496}{2244}$.
\begin{enumerate}
    \item Calculer $\text{PGCD}(2244 ; 1496)$ en détaillant l'algorithme d'Euclide sous forme de tableau (colonnes : Dividende, Diviseur, Quotient, Reste).
    \item Écrire la fraction sous sa forme irréductible.
    \item Donner une fraction équivalente à $\dfrac{1496}{2244}$ dont le numérateur est $34$.
\end{enumerate}
\end{exercice}

\begin{exercice}[Problème de bouquets - Contexte camerounais]
Un fleuriste du marché de Nkolbisson à Yaoundé dispose de $180$ roses et de $252$ œillets. Il veut composer des bouquets identiques en utilisant toutes les fleurs, chaque bouquet contenant le même nombre de roses et le même nombre d'œillets. Il souhaite faire le nombre maximal de bouquets possibles.
\begin{enumerate}
    \item Quel est le nombre maximal de bouquets qu'il peut composer ? Justifier.
    \item Combien y aura-t-il de roses et d'œillets dans chaque bouquet ?
    \item Si chaque bouquet est vendu $2\,500$ FCFA, quel sera le chiffre d'affaires total ?
\end{enumerate}
\end{exercice}

\begin{exercice}[Problème de bus - PPCM]
Deux lignes de bus (Ligne A et Ligne B) quittent la gare routière d'Ekounou à Douala à $6$ h $00$ précises.
\begin{itemize}
    \item La Ligne A part toutes les $24$ minutes.
    \item La Ligne B part toutes les $40$ minutes.
\end{itemize}
\begin{enumerate}
    \item Calculer le PPCM de $24$ et $40$.
    \item À quelle heure les deux bus repartiront-ils ensemble pour la première fois après $6$ h $00$ ?
    \item Combien de fois chaque bus est-il parti (y compris le départ de $6$ h $00$) lorsque ils repartent ensemble à cette heure ?
\end{enumerate}
\end{exercice}

\courssub{Je me prépare au Probatoire}

\begin{exercice}[Problème de synthèse : Clôture d'un terrain]
Un propriétaire terrien à Dschang possède un terrain rectangulaire de $294$ mètres de long sur $210$ mètres de large. Il souhaite le clôturer en plantant des poteaux régulièrement espacés, un poteau étant obligatoirement planté à chaque sommet du rectangle. Il veut minimiser le nombre total de poteaux utilisés.
\begin{enumerate}
    \item Quelle doit être la distance entre deux poteaux consécutifs pour que le nombre de poteaux soit minimal ? Justifier en utilisant la notion de PGCD.
    \item Calculer le nombre total de poteaux nécessaires pour clôturer tout le terrain.
    \item Si chaque poteau coûte $1\,200$ FCFA et que la main-d'œuvre pour la pose d'un poteau est de $300$ FCFA, calculer le coût total de la clôture.
\end{enumerate}
\end{exercice}

\begin{exercice}[Relation PGCD-PPCM et recherche de nombres]
Soient $a$ et $b$ deux entiers naturels non nuls tels que :
\begin{itemize}
    \item $\text{PGCD}(a ; b) = 18$
    \item $a \times b = 9\,720$
\end{itemize}
\begin{enumerate}
    \item Calculer $\text{PPCM}(a ; b)$ en utilisant la relation fondamentale entre PGCD et PPCM.
    \item On suppose maintenant que $a < b$ et que $a = 108$. Déterminer la valeur de $b$.
    \item Vérifier que $\text{PGCD}(108 ; b) = 18$ et $\text{PPCM}(108 ; b)$ correspond à la valeur trouvée à la question 1.
\end{enumerate}
\end{exercice}

\begin{exercice}[Démonstration et application algorithmique]
\begin{enumerate}
    \item \textbf{Démonstration.} Soient $a$ et $b$ deux entiers naturels tels que $a > b > 0$. On note $r$ le reste de la division euclidienne de $a$ par $b$ ($a = bq + r$ avec $0 \le r < b$).
    Démontrer que l'ensemble des diviseurs communs à $a$ et $b$ est identique à l'ensemble des diviseurs communs à $b$ et $r$. En déduire que $\text{PGCD}(a ; b) = \text{PGCD}(b ; r)$.
    \item \textbf{Application.} En utilisant la propriété démontrée ci-dessus (algorithme d'Euclide), calculer $\text{PGCD}(1\,274 ; 546)$. Présenter les divisions successives clairement.
    \item On considère la fraction $\dfrac{546}{1\,274}$. Écrire cette fraction sous sa forme irréductible.
\end{enumerate}
\end{exercice}

\newpage

\coursec{Corrigés des exercices}

\courssub{Je vérifie mes connaissances}

\begin{corrige}[Exercice 1 — QCM : Diviseurs et PGCD]
Analysons chaque affirmation.

\textbf{Affirmation 1 : FAUSSE.} La liste des diviseurs de $18$ est incomplète : il manque $18$ lui-même. En effet $18 = 18 \times 1$. Les diviseurs de $18$ sont : $1 ; 2 ; 3 ; 6 ; 9 ; 18$.

\textbf{Affirmation 2 : FAUSSE.} Les diviseurs de $24$ sont $1 ; 2 ; 3 ; 4 ; 6 ; 8 ; 12 ; 24$ et ceux de $36$ sont $1 ; 2 ; 3 ; 4 ; 6 ; 9 ; 12 ; 18 ; 36$. Le plus grand diviseur commun est $12$, donc $\text{PGCD}(24 ; 36) = 12 \neq 6$.

\textbf{Affirmation 3 : EXACTE.} Les diviseurs de $15$ sont $1 ; 3 ; 5 ; 15$ et ceux de $28$ sont $1 ; 2 ; 4 ; 7 ; 14 ; 28$. Le seul diviseur commun est $1$, donc $\text{PGCD}(15 ; 28) = 1$ : les nombres $15$ et $28$ sont \cle{premiers entre eux}.

\textbf{Affirmation 4 : FAUSSE.} Si $a$ est un multiple de $b$, alors $b$ divise $a$, donc $\text{PGCD}(a ; b) = b$ (et non $a$). Contre-exemple : $\text{PGCD}(12 ; 4) = 4$.

\textbf{Conclusion :} la seule affirmation exacte est l'affirmation \textbf{3}.
\end{corrige}

\begin{corrige}[Exercice 2 — Vrai ou Faux ? Justifier]
\begin{enumerate}
    \item \textbf{VRAI.} On sait que $\text{PGCD}(a ; b) \times \text{PPCM}(a ; b) = a \times b$. Si $a$ et $b$ sont premiers entre eux, $\text{PGCD}(a ; b) = 1$, donc $\text{PPCM}(a ; b) = a \times b$.
    \item \textbf{FAUX.} Deux nombres pairs admettent tous deux $2$ comme diviseur, donc $\text{PGCD}(a ; b) \geq 2 \neq 1$. Contre-exemple : $4$ et $6$ sont pairs et $\text{PGCD}(4 ; 6) = 2$.
    \item \textbf{VRAI.} C'est la propriété fondamentale qui justifie l'\cle{algorithme d'Euclide} : tout diviseur commun à $a$ et $b$ est diviseur commun à $b$ et $r$, et réciproquement.
    \item \textbf{FAUX.} Lorsque les deux entiers sont égaux, PGCD et PPCM sont égaux : $\text{PGCD}(6 ; 6) = \text{PPCM}(6 ; 6) = 6$. L'affirmation « toujours plus grand » est donc fausse (en général on a seulement $\text{PPCM}(a;b) \geq \text{PGCD}(a;b)$).
\end{enumerate}
\end{corrige}

\begin{corrige}[Exercice 3 — Définitions et notations]
\begin{enumerate}
    \item On appelle \textbf{diviseur} d'un entier naturel $a$ tout entier naturel $d$ tel que le reste de la division de $a$ par $d$ soit nul.
    \item Le \textbf{PGCD} de deux entiers naturels non nuls $a$ et $b$ est le plus grand de leurs diviseurs communs. On le note \textbf{PGCD$(a ; b)$}.
    \item Deux entiers naturels non nuls $a$ et $b$ sont dits \textbf{premiers entre eux} si leur PGCD est égal à $1$.
    \item Le \textbf{PPCM} de deux entiers naturels non nuls $a$ et $b$ est le plus petit de leurs multiples communs non nuls. On le note \textbf{PPCM$(a ; b)$}.
\end{enumerate}
\end{corrige}

\begin{corrige}[Exercice 4 — Calculs directs : PGCD et PPCM]
\begin{enumerate}
    \item \textbf{$a = 144$, $b = 60$.} Algorithme d'Euclide :
    \[ 144 = 60 \times 2 + 24 \quad ; \quad 60 = 24 \times 2 + 12 \quad ; \quad 24 = 12 \times 2 + 0. \]
    Le dernier reste non nul est $12$, donc $\text{PGCD}(144 ; 60) = 12$.
    \[ \text{PPCM}(144 ; 60) = \dfrac{144 \times 60}{12} = \dfrac{8\,640}{12} = 720. \]

    \item \textbf{$a = 231$, $b = 154$.}
    \[ 231 = 154 \times 1 + 77 \quad ; \quad 154 = 77 \times 2 + 0. \]
    Donc $\text{PGCD}(231 ; 154) = 77$ et $\text{PPCM}(231 ; 154) = \dfrac{231 \times 154}{77} = 3 \times 154 = 462$.

    \item \textbf{$a = 377$, $b = 221$.}
    \begin{align*}
    377 &= 221 \times 1 + 156 \\
    221 &= 156 \times 1 + 65 \\
    156 &= 65 \times 2 + 26 \\
    65 &= 26 \times 2 + 13 \\
    26 &= 13 \times 2 + 0
    \end{align*}
    Donc $\text{PGCD}(377 ; 221) = 13$ et $\text{PPCM}(377 ; 221) = \dfrac{377 \times 221}{13} = 29 \times 221 = 6\,409$.

    \item \textbf{$a = 48$, $b = 180$.}
    \[ 180 = 48 \times 3 + 36 \quad ; \quad 48 = 36 \times 1 + 12 \quad ; \quad 36 = 12 \times 3 + 0. \]
    Donc $\text{PGCD}(48 ; 180) = 12$ et $\text{PPCM}(48 ; 180) = \dfrac{48 \times 180}{12} = 4 \times 180 = 720$.
    Vérification : $\text{PGCD} \times \text{PPCM} = 12 \times 720 = 8\,640$ et $a \times b = 48 \times 180 = 8\,640$. La relation est bien vérifiée.
\end{enumerate}
\end{corrige}

\courssub{Je m'entraîne}

\begin{corrige}[Exercice 5 — Algorithme des soustractions]
\begin{enumerate}
    \item \textbf{Propriété :} si $a > b$, tout diviseur commun à $a$ et $b$ divise aussi $a - b$, et tout diviseur commun à $a - b$ et $b$ divise aussi $a = (a-b) + b$. Les deux couples $(a ; b)$ et $(a-b ; b)$ ont donc les mêmes diviseurs communs, d'où $\text{PGCD}(a ; b) = \text{PGCD}(a-b ; b)$.
    \item Appliquons la méthode (on remplace à chaque étape le plus grand nombre par la différence) :
    \begin{align*}
    \text{PGCD}(437 ; 323) &= \text{PGCD}(114 ; 323) & (437 - 323 &= 114)\\
    &= \text{PGCD}(114 ; 209) & (323 - 114 &= 209)\\
    &= \text{PGCD}(114 ; 95) & (209 - 114 &= 95)\\
    &= \text{PGCD}(19 ; 95) & (114 - 95 &= 19)\\
    &= \text{PGCD}(19 ; 76) & (95 - 19 &= 76)\\
    &= \text{PGCD}(19 ; 57) & (76 - 19 &= 57)\\
    &= \text{PGCD}(19 ; 38) & (57 - 19 &= 38)\\
    &= \text{PGCD}(19 ; 19) & (38 - 19 &= 19)
    \end{align*}
    \item Les deux nombres obtenus sont égaux à $19$, donc $\text{PGCD}(437 ; 323) = 19$.
\end{enumerate}
\textbf{Vérification (rapide) :} $437 = 19 \times 23$ et $323 = 19 \times 17$ ; comme $17$ et $23$ sont premiers entre eux, le PGCD est bien $19$.
\end{corrige}

\begin{corrige}[Exercice 6 — Fraction irréductible]
\begin{enumerate}
    \item Algorithme d'Euclide pour $\text{PGCD}(2\,244 ; 1\,496)$ :
    \begin{center}
    \begin{tabular}{|c|c|c|c|}
    \hline
    \rowcolor{popblueL} \textbf{Dividende} & \textbf{Diviseur} & \textbf{Quotient} & \textbf{Reste} \\
    \hline
    $2\,244$ & $1\,496$ & $1$ & $748$ \\
    \hline
    $1\,496$ & $748$ & $2$ & $0$ \\
    \hline
    \end{tabular}
    \end{center}
    Le dernier reste non nul est $748$, donc $\text{PGCD}(2\,244 ; 1\,496) = 748$.
    \item On simplifie numérateur et dénominateur par $748$ :
    \[ \dfrac{1\,496}{2\,244} = \dfrac{1\,496 \div 748}{2\,244 \div 748} = \dfrac{2}{3}. \]
    Comme $\text{PGCD}(2 ; 3) = 1$, la fraction $\dfrac{2}{3}$ est \cle{irréductible}.
    \item On cherche une fraction équivalente de numérateur $34$ : comme $34 = 2 \times 17$, on multiplie numérateur et dénominateur de $\dfrac{2}{3}$ par $17$ :
    \[ \dfrac{2}{3} = \dfrac{2 \times 17}{3 \times 17} = \dfrac{34}{51}. \]
    La fraction cherchée est $\dfrac{34}{51}$.
\end{enumerate}
\end{corrige}

\begin{corrige}[Exercice 7 — Problème de bouquets (marché de Nkolbisson)]
\begin{enumerate}
    \item Le nombre de bouquets doit diviser à la fois le nombre de roses ($180$) et le nombre d'œillets ($252$), car chaque bouquet contient le même nombre de chaque fleur et toutes les fleurs sont utilisées. Pour un nombre \emph{maximal} de bouquets, on cherche le plus grand diviseur commun :
    \begin{align*}
    252 &= 180 \times 1 + 72 \\
    180 &= 72 \times 2 + 36 \\
    72 &= 36 \times 2 + 0
    \end{align*}
    Donc $\text{PGCD}(252 ; 180) = 36$. Le fleuriste peut composer au maximum \textbf{36 bouquets}.
    \item Chaque bouquet contient :
    \[ \dfrac{180}{36} = 5 \text{ roses} \qquad \text{et} \qquad \dfrac{252}{36} = 7 \text{ œillets.} \]
    \item Chiffre d'affaires total : $36 \times 2\,500 = 90\,000$ FCFA.
\end{enumerate}
\textbf{Point de vigilance :} il faut bien justifier que le nombre de bouquets est un \emph{diviseur commun} aux deux quantités, puis que « maximal » signifie qu'on prend le \emph{plus grand} d'entre eux, d'où le PGCD.
\end{corrige}

\begin{corrige}[Exercice 8 — Problème de bus (PPCM)]
\begin{enumerate}
    \item Les départs communs correspondent aux multiples communs de $24$ et de $40$. Calculons le PGCD puis le PPCM :
    \[ 40 = 24 \times 1 + 16 \quad ; \quad 24 = 16 \times 1 + 8 \quad ; \quad 16 = 8 \times 2 + 0 \]
    donc $\text{PGCD}(24 ; 40) = 8$ et
    \[ \text{PPCM}(24 ; 40) = \dfrac{24 \times 40}{8} = \dfrac{960}{8} = 120. \]
    \item Les deux bus repartent ensemble toutes les $120$ minutes, c'est-à-dire toutes les $2$ heures. Première coïncidence après $6$ h $00$ : $6\text{ h} + 2\text{ h} = \textbf{8 h 00}$.
    \item Nombre de départs (en comptant celui de $6$ h $00$) :
    \begin{itemize}
        \item Ligne A : $\dfrac{120}{24} + 1 = 5 + 1 = 6$ départs ;
        \item Ligne B : $\dfrac{120}{40} + 1 = 3 + 1 = 4$ départs.
    \end{itemize}
\end{enumerate}
\textbf{Point de vigilance :} ne pas oublier le $+1$ : à $8$ h $00$, le départ de $6$ h $00$ compte aussi. La Ligne A est partie à $6$ h $00$, $6$ h $24$, $6$ h $48$, $7$ h $12$, $7$ h $36$ et $8$ h $00$.
\end{corrige}

\courssub{Je me prépare au Probatoire}

\begin{corrige}[Exercice 9 — Problème de synthèse : Clôture d'un terrain]
\begin{enumerate}
    \item Un poteau doit être planté à chaque sommet, et les poteaux sont régulièrement espacés sur la longueur ($294$ m) et sur la largeur ($210$ m). La distance $d$ entre deux poteaux consécutifs doit donc diviser \emph{à la fois} $294$ et $210$. Minimiser le nombre de poteaux revient à maximiser l'écartement, donc $d = \text{PGCD}(294 ; 210)$.
    \begin{align*}
    294 &= 210 \times 1 + 84 \\
    210 &= 84 \times 2 + 42 \\
    84 &= 42 \times 2 + 0
    \end{align*}
    Donc $\text{PGCD}(294 ; 210) = 42$. La distance entre deux poteaux consécutifs doit être de \textbf{42 mètres}.
    \item Le périmètre du terrain est :
    \[ P = 2 \times (294 + 210) = 2 \times 504 = 1\,008 \text{ m}. \]
    Sur un contour fermé, le nombre de poteaux est égal au nombre d'intervalles :
    \[ N = \dfrac{1\,008}{42} = 24 \text{ poteaux.} \]
    (Vérification par côté : $\frac{294}{42} = 7$ intervalles par longueur et $\frac{210}{42} = 5$ par largeur, soit $2 \times 7 + 2 \times 5 = 24$ poteaux, les sommets n'étant comptés qu'une fois.)
    \item Coût par poteau (poteau + pose) : $1\,200 + 300 = 1\,500$ FCFA.
    \[ \text{Coût total} = 24 \times 1\,500 = 36\,000 \text{ FCFA.} \]
\end{enumerate}
\textbf{Barème indicatif (sur 5 pts) :} justification « diviseur commun » et choix du PGCD : $1$ pt ; calcul du PGCD détaillé : $1$ pt ; périmètre et nombre de poteaux : $1{,}5$ pt ; coût total : $1$ pt ; présentation et unités : $0{,}5$ pt.

\textbf{Points de vigilance :} citer explicitement que l'écartement doit diviser les deux dimensions \emph{à cause des sommets} ; sur un contour fermé, ne pas ajouter $1$ au nombre d'intervalles ; ne pas oublier la main-d'œuvre dans le coût unitaire.
\end{corrige}

\begin{corrige}[Exercice 10 — Relation PGCD-PPCM et recherche de nombres]
\begin{enumerate}
    \item D'après la relation fondamentale $\text{PGCD}(a ; b) \times \text{PPCM}(a ; b) = a \times b$ :
    \[ \text{PPCM}(a ; b) = \dfrac{a \times b}{\text{PGCD}(a ; b)} = \dfrac{9\,720}{18} = 540. \]
    \item Si $a = 108$, alors $b = \dfrac{9\,720}{108} = 90$. Comme $108 > 90$, la condition $a < b$ n'est pas respectée : on échange les rôles et on pose $a = 90$ et $b = 108$. Le couple de nombres cherché est donc $\{90 ; 108\}$, avec $b = 108$ si l'on impose $a = 90 < b$.
    \item Vérification avec $108$ et $90$ :
    \[ 108 = 90 \times 1 + 18 \quad ; \quad 90 = 18 \times 5 + 0 \]
    donc $\text{PGCD}(108 ; 90) = 18$ ✓ et
    \[ \text{PPCM}(108 ; 90) = \dfrac{108 \times 90}{18} = \dfrac{9\,720}{18} = 540, \]
    ce qui correspond bien à la valeur trouvée à la question 1. ✓
\end{enumerate}
\textbf{Points de vigilance :} la relation fondamentale n'est valable que pour des entiers naturels \emph{non nuls} — le dire ; vérifier toujours la cohérence des conditions ($a < b$) et ne pas hésiter à échanger $a$ et $b$, le PGCD et le PPCM étant symétriques.
\end{corrige}

\begin{corrige}[Exercice 11 — Démonstration et application algorithmique]
\begin{enumerate}
    \item \textbf{Démonstration.} On a $a = bq + r$ avec $0 \leq r < b$. Montrons la double inclusion.
    \begin{itemize}
        \item Soit $d$ un diviseur commun à $a$ et $b$ : $d$ divise $a$ et $d$ divise $b$, donc $d$ divise toute combinaison $a - bq$, c'est-à-dire $d$ divise $r$. Ainsi $d$ est un diviseur commun à $b$ et $r$.
        \item Réciproquement, soit $d'$ un diviseur commun à $b$ et $r$ : $d'$ divise $bq$ et $r$, donc $d'$ divise $bq + r = a$. Ainsi $d'$ est un diviseur commun à $a$ et $b$.
    \end{itemize}
    L'ensemble des diviseurs communs à $a$ et $b$ est donc \emph{identique} à l'ensemble des diviseurs communs à $b$ et $r$. Ces deux ensembles ayant les mêmes éléments, ils ont le même plus grand élément, d'où :
    \[ \text{PGCD}(a ; b) = \text{PGCD}(b ; r). \]
    \item \textbf{Application} au calcul de $\text{PGCD}(1\,274 ; 546)$ :
    \begin{align*}
    1\,274 &= 546 \times 2 + 182 \\
    546 &= 182 \times 3 + 0
    \end{align*}
    Le dernier reste non nul est $182$, donc $\text{PGCD}(1\,274 ; 546) = 182$.
    \item On simplifie par $182$ : $546 = 182 \times 3$ et $1\,274 = 182 \times 7$, donc
    \[ \dfrac{546}{1\,274} = \dfrac{182 \times 3}{182 \times 7} = \dfrac{3}{7}. \]
    Comme $\text{PGCD}(3 ; 7) = 1$, la forme irréductible est $\dfrac{3}{7}$.
\end{enumerate}
\textbf{Barème indicatif (sur 6 pts) :} double inclusion rédigée : $2$ pts ; conclusion sur l'égalité des PGCD : $1$ pt ; divisions euclidiennes exactes : $1{,}5$ pt ; fraction irréductible justifiée : $1$ pt ; rédaction : $0{,}5$ pt.

\textbf{Points de vigilance :} pour la démonstration, traiter \emph{les deux sens} de l'inclusion (c'est l'oubli le plus fréquent) ; citer la condition $0 \leq r < b$ qui garantit que le procédé s'arrête ; pour la fraction, préciser que diviser numérateur et dénominateur par le PGCD donne deux entiers premiers entre eux, donc une fraction irréductible.
\end{corrige}

\newpage

\coursec{Fiche bilan}

\begin{popfigure}
\centering
\begin{tikzpicture}[
  mindmap,
  grow cyclic,
  every node/.style={concept, execute at begin node=\hskip0pt},
  concept color=popblue,
  level 1/.append style={level distance=4.5cm, sibling angle=72},
  level 2/.append style={level distance=3cm, sibling angle=45},
  texte court/.style={text width=2.8cm, align=center, font=\small\sffamily}
]
\node[texte court, concept color=popblue, align=center]{Arithmétique\\3ème}
  child[concept color=poporange] {node[texte court, align=center]{Division\\Euclidienne\\$a = bq + r$}}
  child[concept color=popgreen] {node[texte court, align=center]{PGCD\\Listes \textbullet\ Soustractions\\Algorithme d'Euclide}}
  child[concept color=poppurple] {node[texte court, align=center]{Premiers\\entre eux\\PGCD = 1}}
  child[concept color=poppink] {node[texte court, align=center]{PPCM\\Multiples communs\\Relation $a\times b$}}
  child[concept color=popgold] {node[texte court, align=center]{Applications\\Simplification\\Pavage \textbullet\ Rendez-vous}};
\end{tikzpicture}
\legende{Carte mentale : Arithmétique -- PGCD, PPCM et applications (3ème)}
\end{popfigure}

\begin{bilan}

\courssub{Définitions clés}

\begin{itemize}
\item \textbf{Division euclidienne} : $\forall a\in\mathbb{N}, \forall b\in\mathbb{N}^*$, $\exists!\ (q,r)\in\mathbb{N}^2$ tel que $a = bq + r$ et $0 \le r < b$. ($q$ quotient, $r$ reste).
\item \textbf{Diviseur} : $d$ divise $a$ ($d\mid a$) ssi $\exists k\in\mathbb{N},\ a = dk$. Equivaut à : le reste de la division de $a$ par $d$ est $0$.
\item \textbf{PGCD} : Plus Grand Commun Diviseur de $a$ et $b$ ($a,b$ non tous deux nuls) = plus grand élément de l'ensemble des diviseurs communs. Noté $\operatorname{PGCD}(a,b)$.
\item \textbf{PPCM} : Plus Petit Commun Multiple de $a$ et $b$ (non nuls) = plus petit élément non nul de l'ensemble des multiples communs. Noté $\operatorname{PPCM}(a,b)$.
\item \textbf{Premiers entre eux} : $a$ et $b$ sont premiers entre eux $\iff \operatorname{PGCD}(a,b)=1$.
\end{itemize}

\courssub{Théorèmes \& Formules à connaître par cœur}

\begin{center}
\begin{tabularx}{\linewidth}{|>{\bfseries}l|X|}\hline
\textbf{Propriété / Théorème} & \textbf{Énoncé / Formule} \\ \hline
Algorithme des soustractions & $\operatorname{PGCD}(a,b) = \operatorname{PGCD}(a-b,b)$ pour $a\ge b>0$. On répète jusqu'à obtenir deux nombres égaux. \\ \hline
Algorithme d'Euclide & $\operatorname{PGCD}(a,b) = \operatorname{PGCD}(b,r)$ où $r$ est le reste de la division euclidienne de $a$ par $b$ ($a\ge b>0$). On répète jusqu'à reste nul. \\ \hline
Relation PGCD -- PPCM & Pour $a,b\in\mathbb{N}^*$ : $\operatorname{PGCD}(a,b) \times \operatorname{PPCM}(a,b) = a \times b$. \\ \hline
Conséquence simplification & La fraction $\frac{a}{b}$ est irréductible $\iff \operatorname{PGCD}(a,b)=1$. Simplification : $\frac{a}{b} = \frac{a/\operatorname{PGCD}(a,b)}{b/\operatorname{PGCD}(a,b)}$. \\ \hline
\end{tabularx}
\end{center}

\courssub{Méthodes express}

\begin{enumerate}
\item \textbf{PGCD par Euclide (recommandé)} : Poser les divisions successives $a = bq_1 + r_1$, $b = r_1q_2 + r_2$, ... jusqu'à $r_k=0$. Le dernier reste non nul est le PGCD.
\item \textbf{PPCM via la relation} : Calculer $d=\operatorname{PGCD}(a,b)$, puis $\operatorname{PPCM}(a,b) = \frac{a\times b}{d}$.
\item \textbf{Problème de « rendez-vous » / pavage} : Chercher un multiple commun $\rightarrow$ PPCM. Problème de « partage / découpe maximale » $\rightarrow$ PGCD.
\item \textbf{Rédaction Euclide} : Aligner les égalités : $a = bq_1 + r_1$, $b = r_1q_2 + r_2$, ... Conclure : « Le dernier reste non nul est $r_k$, donc $\operatorname{PGCD}(a,b)=r_k$. »
\end{enumerate}

\end{bilan}

\begin{aretenir}[Checklist avant le BEPC]
\begin{itemize}
\item $\square$ Je sais poser et vérifier une division euclidienne ($a=bq+r$, $0\le r<b$).
\item $\square$ Je sais lister les diviseurs d'un entier $\le 500$.
\item $\square$ Je maîtrise l'algorithme des soustractions pour trouver le PGCD.
\item $\square$ Je maîtrise l'algorithme d'Euclide (divisions successives) pour trouver le PGCD.
\item $\square$ Je sais déterminer si deux entiers sont premiers entre eux.
\item $\square$ Je connais par cœur la relation $\operatorname{PGCD}(a,b)\times\operatorname{PPCM}(a,b)=a\times b$.
\item $\square$ Je sais calculer le PPCM de deux entiers (via la relation ou par listes de multiples).
\item $\square$ Je sais simplifier une fraction en divisant numérateur et dénominateur par leur PGCD.
\item $\square$ Je sais modéliser un problème concret : « partage maximal » $\rightarrow$ PGCD ; « rendez-vous / synchronisation » $\rightarrow$ PPCM.
\item $\square$ Je rédige mes divisions d'Euclide alignées et je conclus par une phrase.
\item $\square$ Je vérifie la cohérence de mes résultats (PGCD divise les deux nombres ; PPCM est multiple des deux).
\item $\square$ J'utilise le vocabulaire précis : diviseur, multiple, quotient, reste, PGCD, PPCM, premiers entre eux.
\end{itemize}
\end{aretenir}

\findecours

\end{document}
